% =====================================================================
% Chương 3 – Phương pháp nghiên cứu
% =====================================================================
\section{Phương pháp nghiên cứu}
\label{sec:method}

\subsection{Quy trình tổng thể}
\label{subsec:pipeline}

Đề tài thực hiện đủ sáu giai đoạn của quy trình CRISP-DM (Mục~\ref{subsec:crispdm}). Bảng~\ref{tab:pipeline} tóm tắt công việc ở từng giai đoạn, mục trình bày trong báo cáo và mã nguồn tương ứng trong GitHub của nhóm. Mọi bước có yếu tố ngẫu nhiên đều cố định hạt giống \texttt{SEED = 42}.

\begin{table}[H]
  \centering
  \caption{Các bước thực hiện theo quy trình CRISP-DM.}
  \label{tab:pipeline}
  \small
  \begin{tabular}{@{}p{0.17\textwidth}p{0.47\textwidth}p{0.29\textwidth}@{}}
    \toprule
    \textbf{Giai đoạn} & \textbf{Công việc} & \textbf{Mã nguồn} \\
    \midrule
    Hiểu bài toán & Xác định đầu vào, đầu ra, thước đo Macro-F1, mục tiêu (Chương~\ref{sec:intro}) & -- \\
    Hiểu dữ liệu & Phân tích khám phá: cấu trúc, chất lượng, phân phối, khác biệt giữa các tập (Mục~\ref{sec:eda}) & \texttt{eda/} \\
    Chuẩn bị dữ liệu & Làm sạch, lấy mẫu khung, chuẩn hoá, chia tập, tăng cường (Mục~\ref{sec:preprocess}); trích xuất và gộp đặc trưng (Mục~\ref{subsec:features}) & \texttt{processing/}, \texttt{models/extract\_*.ipynb} \\
    Mô hình hoá & Bốn thuật toán cổ điển và CNN-LSTM, tìm siêu tham số bằng cross-validation (Mục~\ref{sec:model}, \ref{subsec:res_tuning}) & \texttt{models/svm/}, \texttt{models/classical/}, \texttt{models/cnn\_lstm/} \\
    Đánh giá & Fold~0 với nhiều chỉ số, tập private trên AIChallenge (Mục~\ref{subsec:eval_design}, \ref{subsec:res_fold0}--\ref{subsec:res_private}); diễn giải mô hình (Mục~\ref{subsec:res_explain}) & như trên \\
    Triển khai & Ứng dụng demo Streamlit (Mục~\ref{subsec:demo_design}); mã nguồn kèm \texttt{README} và \texttt{requirements.txt} (Phụ lục~\ref{app:repo}) & \texttt{demo/} \\
    \bottomrule
  \end{tabular}
\end{table}

\paragraph*{Vòng lặp trong quy trình} Đề tài đi qua các giai đoạn hai lần. Ở lần đầu, các mô hình dùng đặc trưng từ EfficientNet-B0 đóng băng đạt Macro-F1 cao trên fold~0 nhưng giảm mạnh trên tập private (Mục~\ref{subsec:res_private}). Kết quả đánh giá này đưa nhóm quay lại giai đoạn chuẩn bị dữ liệu: trích bộ đặc trưng mới từ CNN đã tinh chỉnh trên dữ liệu ký hiệu (Mục~\ref{subsec:features}), rồi huấn luyện và đánh giá lại cả bốn thuật toán cổ điển trên bộ đặc trưng này.

\subsection{Thu thập dữ liệu}
\label{subsec:data_source}

Nhóm sử dụng bộ dữ liệu công khai của cuộc thi \textit{Olympic Trí tuệ Nhân tạo Sinh viên 2025 -- Vòng Sơ loại}~\cite{ptit2025} do Khoa Trí tuệ Nhân tạo, Học viện Công nghệ Bưu chính Viễn thông (PTIT) tổ chức, tải từ nền tảng AIChallenge của PTIT. Bộ dữ liệu gồm thư mục \texttt{train/} (mỗi lớp một thư mục con), hai thư mục \texttt{public\_test/} và \texttt{private\_test/} không có nhãn, cùng tệp \texttt{label\_mapping.pkl} ánh xạ 100 tên lớp sang chỉ số 0--99. Nhãn của hai tập kiểm tra do ban tổ chức giữ; nhóm nộp tệp dự đoán lên nền tảng để được chấm Macro-F1 (Mục~\ref{subsec:eval_design}). Các thuộc tính mô tả video (số khung, FPS, độ phân giải, dung lượng, mã băm\ldots) được nhóm tự trích xuất bằng OpenCV như trình bày ở mục tiếp theo.

% 3.3 – Phân tích khám phá dữ liệu (EDA)
% =====================================================================
% Mục 3.1 – Phân tích dữ liệu (EDA)  [chèn bằng % =====================================================================
% Mục 3.1 – Phân tích dữ liệu (EDA)  [chèn bằng % =====================================================================
% Mục 3.1 – Phân tích dữ liệu (EDA)  [chèn bằng \input{chapters/eda}]
% Cấu trúc theo khung EDA chuẩn:
%   3.1.1 Mục tiêu và tổng quan dữ liệu
%   3.1.2 Cấu trúc dữ liệu
%   3.1.3 Chất lượng dữ liệu
%   3.1.4 Phân tích đơn biến
%   3.1.5 Phân tích biến mục tiêu
%   3.1.6 Phân tích hai biến và đa biến
%   3.1.7 So sánh phân phối giữa các tập
%   3.1.8 Kết luận và khuyến nghị
% Gói cần: graphicx, booktabs, float, caption, amsmath
% =====================================================================

\subsection{Phân tích dữ liệu (EDA)}
\label{sec:eda}

% \edafile{thư mục}{tệp}: in đường dẫn script, cho phép xuống dòng sau dấu /
\newcommand{\edafile}[2]{\texttt{#1/}\allowbreak\texttt{#2}}
% Cho phép xuống dòng sau dấu gạch dưới trong tên tệp dài (vd. class\_counts\_train\_unique\_vs\_all.csv)
\let\edaOrigUS\_
\renewcommand{\_}{\edaOrigUS\allowbreak}
\sloppy

Mục này trình bày kết quả phân tích khám phá dữ liệu (\textit{Exploratory Data Analysis -- EDA}) trên bộ dữ liệu video ngôn ngữ ký hiệu. EDA chỉ nhằm \emph{quan sát} xu hướng, phân phối và chất lượng dữ liệu. Nhóm sử dụng Python làm công cụ để thực hiện EDA (\texttt{os}, \texttt{pickle}, \texttt{pandas}, \texttt{OpenCV}, \texttt{numpy}, \texttt{hashlib}, \texttt{matplotlib}).

Phần phân tích được tổ chức theo trình tự quy trình EDA: tổng quan và cấu trúc dữ liệu (Mục~\ref{subsec:eda_overview}--\ref{subsec:eda_structure}), chất lượng dữ liệu (Mục~\ref{subsec:eda_quality}), phân tích đơn biến (Mục~\ref{subsec:eda_uni}), phân tích biến mục tiêu (Mục~\ref{subsec:eda_class}), phân tích hai biến và đa biến (Mục~\ref{subsec:eda_multi}), so sánh giữa các tập (Mục~\ref{subsec:eda_shift}) và cuối cùng là kết luận kèm khuyến nghị (Mục~\ref{subsec:eda_conclusion}).

\begin{table}[H]
  \centering
  \caption{Các script EDA, thứ tự chạy và mục sử dụng kết quả.}
  \label{tab:eda_scripts}
  \small
  \begin{tabular}{@{}p{0.17\textwidth}p{0.26\textwidth}p{0.27\textwidth}p{0.18\textwidth}@{}}
    \toprule
    \textbf{Bước (thư mục)} & \textbf{Script -- câu hỏi} & \textbf{Kết quả chính} & \textbf{Dùng ở} \\
    \midrule
    \texttt{01\_overview} & \texttt{01\_scan\_metadata.py} -- dữ liệu gồm những gì? & \texttt{video\_metadata.csv} & 3.1.1 \\
      & \texttt{02\_dataset\_structure.py} -- chia tập, định dạng, tệp nhãn & \texttt{structure\_summary.txt} \newline \texttt{split\_sizes.png} & 3.1.2 \\
    \midrule
    \texttt{02\_data\_quality} & \texttt{01\_integrity\_missing.py} -- giá trị thiếu, tệp lỗi & \texttt{integrity\_summary.txt} & 3.1.3 \\
      & \texttt{02\_duplicates.py} -- trùng lặp, xung đột nhãn, trùng giữa các tập & \texttt{file\_hashes.csv} \newline \texttt{dups\_summary.txt} \newline \texttt{dedup\_summary.txt} \newline \texttt{class\_counts\_train\_unique\_vs\_all.csv} & 3.1.3, 3.1.5 \\
      & \texttt{03\_frame\_scan.py} -- quét nội dung từng khung & \texttt{video\_quality\_flags.csv} & 3.1.3, 3.1.6, 3.1.7 \\
      & \texttt{04\_frame\_quality\_summary.py} -- tổng hợp khung đơn sắc & \texttt{solid\_frame\_summary.txt} \newline \texttt{class\_solid\_frame\_ratio.csv} & 3.1.3, 3.1.6 \\
    \midrule
    \texttt{03\_univariate} & \texttt{01\_numeric\_attributes.py} -- phân phối, ngoại lai & \texttt{desc\_by\_split.csv} \newline \texttt{univariate\_summary.txt} \newline 3 histogram & 3.1.3, 3.1.4 \\
    \midrule
    \texttt{04\_target\_analysis} & \texttt{01\_class\_distribution.py} -- phân phối nhãn & \texttt{class\_summary.txt} \newline \texttt{class\_counts\_train.csv} \newline 3 hình & 3.1.5 \\
    \midrule
    \texttt{05\_bivariate\_multivariate} & \texttt{01\_correlation\_and\_label\_association.py} & \texttt{multivariate\_summary.txt} \newline 3 hình & 3.1.6 \\
    \midrule
    \texttt{06\_split\_comparison} & \texttt{01\_split\_shift.py} -- khác biệt phân phối giữa các tập & \texttt{split\_shift\_summary.txt} \newline \texttt{ecdf\_by\_split.png} & 3.1.7 \\
    \midrule
    \texttt{07\_visual\_inspection} & \texttt{01\_sample\_frames.py} -- khung hình mẫu & 2 lưới khung hình & 3.1.2 \\
      & \texttt{02\_flagged\_videos.py} -- xem lại video bị đánh dấu & \texttt{flagged\_videos\_info.txt} \newline 2 hình & 3.1.3 \\
    \bottomrule
  \end{tabular}
\end{table}

% =====================================================================
\subsubsection{Mục tiêu và tổng quan dữ liệu}
\label{subsec:eda_overview}

\paragraph{Bài toán và thước đo} Bài toán là phân loại video ngôn ngữ ký hiệu (\textit{Sign Language Video Classification}) thành $N=100$ lớp từ vựng tiếng Việt. Đầu vào là một video ghi cử chỉ của một người, được tiền xử lý thành chuỗi $T=16$ khung hình RGB; đầu ra là nhãn dự đoán. Thước đo chính thức là Macro F1-Score. Biến mục tiêu là nhãn lớp (biến phân loại, 100 giá trị), còn các thuộc tính mô tả video (số khung hình, FPS, dung lượng\ldots) chỉ đóng vai trò thông tin phụ trợ cho phân tích.

\paragraph{Nguồn dữ liệu} Bộ dữ liệu là đề chính thức của cuộc thi \textit{Olympic Trí tuệ Nhân tạo Sinh viên 2025 -- Vòng Sơ loại} do Khoa Trí tuệ Nhân tạo, Học viện Công nghệ Bưu chính Viễn thông (PTIT) tổ chức. Các lớp thuộc nhiều nhóm chủ đề: y tế -- dịch bệnh (\textit{Bệnh viện}, \textit{Cách ly}, \textit{Khẩu trang}, \textit{Virus}, \textit{Sốt}, \textit{Ho}\ldots), cảm xúc (\textit{Lo lắng}, \textit{Xúc động}, \textit{Thương}, \textit{Ghét}\ldots), hành động thường ngày (\textit{Ăn}, \textit{Uống}, \textit{Chạy}, \textit{Đi}, \textit{Mua}\ldots), thời gian -- vị trí (\textit{Ban ngày}, \textit{Trưa}, \textit{Hôm nay}, \textit{Phía sau}, \textit{Rẽ trái}\ldots) và phương tiện (\textit{Ô tô}, \textit{Xe máy}, \textit{Xe đạp}).

\paragraph{Từ điển dữ liệu} Vì dữ liệu là video nên không có bảng thuộc tính sẵn. Nhóm duyệt cây thư mục, mở từng video bằng \texttt{cv2.VideoCapture} và gom thông tin vào một bảng metadata duy nhất (\texttt{01\_overview/video\_metadata.csv}, 8.364 dòng). Mã băm MD5 được lưu ở \texttt{02\_data\_quality/file\_hashes.csv}, còn các thuộc tính suy ra từ nội dung khung hình (khung đơn sắc, chuyển động) được lưu ở \texttt{02\_data\_quality/video\_quality\_flags.csv}. Ba bảng có chung khoá (tập, lớp, tên tệp). Bảng~\ref{tab:eda_dict} liệt kê ý nghĩa từng thuộc tính; mọi thống kê phía sau chỉ là phép ghép và \texttt{groupby} trên ba bảng này.

\begin{table}[H]
  \centering
  \caption{Từ điển các thuộc tính dùng trong phân tích.}
  \label{tab:eda_dict}
  \small
  \begin{tabular}{@{}p{0.22\textwidth}p{0.17\textwidth}p{0.5\textwidth}@{}}
    \toprule
    \textbf{Thuộc tính} & \textbf{Kiểu} & \textbf{Ý nghĩa / cách tính} \\
    \midrule
    \texttt{split}        & phân loại & Tập chứa video: \texttt{train}, \texttt{public\_test}, \texttt{private\_test} \\
    \texttt{cls}          & phân loại & Nhãn lớp, lấy từ tên thư mục con của \texttt{train} (không có ở hai tập test) \\
    \texttt{file}         & chuỗi     & Tên tệp \texttt{.mp4}; hậu tố \texttt{\_N} biểu thị bản sao \\
    \texttt{frames}       & số nguyên & Số khung hình (\texttt{CAP\_\allowbreak PROP\_\allowbreak FRAME\_\allowbreak COUNT}) \\
    \texttt{fps}          & số thực   & Tốc độ khung hình ghi trong container \\
    \texttt{w}, \texttt{h}& số nguyên & Chiều rộng, chiều cao khung hình (pixel) \\
    \texttt{dur}          & số thực   & Thời lượng (giây) $=\texttt{frames}/\texttt{fps}$ \\
    \texttt{size\_kb}     & số thực   & Dung lượng tệp (KB) \\
    \texttt{md5}          & chuỗi     & Mã băm nội dung tệp, dùng phát hiện video trùng (lưu trong \texttt{file\_hashes.csv}) \\
    \texttt{solid\_ratio} & số thực   & Tỷ lệ khung đơn sắc trong video (khung có độ lệch chuẩn điểm ảnh $<3$) \\
    \texttt{motion}       & số thực   & Độ chuyển động trung bình $\overline{|f_t-f_{t-1}|}$ giữa các khung liên tiếp (ảnh xám) \\
    \bottomrule
  \end{tabular}
\end{table}

% =====================================================================
\subsubsection{Cấu trúc dữ liệu}
\label{subsec:eda_structure}

\paragraph{Quy mô và cách chia} Sau khi giải nén \texttt{dataset.zip} (khoảng 864~MB), dữ liệu gồm ba tập video và một tệp ánh xạ nhãn (Bảng~\ref{tab:eda_splits}, Hình~\ref{fig:eda_split_sizes}). Tệp \texttt{label\_mapping.pkl} là một \texttt{dict[str, int]} chứa đúng 100 lớp, đánh chỉ số 0--99 theo thứ tự bảng chữ cái tiếng Việt (\textit{An ủi}~$=0$, \textit{Áp dụng}~$=1$, \textit{Ăn}~$=2$, \ldots, \textit{Đồng ý}~$=99$). Khi đối chiếu với tên thư mục lớp, nhóm phát hiện hai bên dùng \emph{hai dạng chuẩn hoá Unicode khác nhau}: khoá trong \texttt{label\_mapping.pkl} ở dạng NFD (dấu tách rời khỏi chữ cái), còn tên thư mục ở dạng NFC (dấu dựng sẵn). Vì vậy 91/100 tên lớp (mọi lớp có dấu) không khớp nếu so nguyên văn, dù hiển thị giống hệt nhau; sau khi chuẩn hoá cả hai về NFC thì 100 lớp khớp hoàn toàn và chỉ số liên tục 0--99 không trùng.

\begin{table}[H]
  \centering
  \caption{Cấu trúc và quy mô các tập dữ liệu.}
  \label{tab:eda_splits}
  \begin{tabular}{@{}lp{0.42\textwidth}rr@{}}
    \toprule
    \textbf{Thành phần} & \textbf{Nội dung} & \textbf{Số video} & \textbf{Dung lượng} \\
    \midrule
    \texttt{train/}         & 100 thư mục con, mỗi thư mục là một lớp (có nhãn) & 3.875 & 323~MB \\
    \texttt{public\_test/}  & Video không nhãn, dùng cho bảng xếp hạng công khai & 1.630 & 181~MB \\
    \texttt{private\_test/} & Video không nhãn, dùng cho xếp hạng cuối cùng      & 2.859 & 325~MB \\
    \texttt{label\_mapping.pkl} & \texttt{dict[str, int]}: tên lớp (Unicode NFD) $\rightarrow$ chỉ số 0--99 & -- & -- \\
    \midrule
    \textbf{Tổng} & & \textbf{8.364} & $\approx$\textbf{829~MB} \\
    \bottomrule
  \end{tabular}
\end{table}

\begin{figure}[H]
  \centering
  \includegraphics[width=0.6\textwidth]{split_sizes.png}
  \caption{Số lượng video trong ba tập \texttt{train}, \texttt{public\_test} và \texttt{private\_test}.}
  \label{fig:eda_split_sizes}
\end{figure}

\paragraph{Định dạng video} Toàn bộ 8.364 tệp đều là \texttt{.mp4} và đọc được; không có video nào có 0 khung hình. Độ phân giải đồng nhất tuyệt đối: 100\% video ở cả ba tập đều là $224\times224$ (Bảng~\ref{tab:eda_res}). Số khung hình gần như cố định ở mức 30 (chi tiết ở Mục~\ref{subsec:eda_uni}).

\begin{table}[H]
  \centering
  \caption{Thống kê độ phân giải video.}
  \label{tab:eda_res}
  \begin{tabular}{@{}lrrr@{}}
    \toprule
    \textbf{Tập} & \textbf{Độ phân giải} & \textbf{Số video} & \textbf{Tỷ lệ} \\
    \midrule
    train         & $224\times224$ & 3.875 & 100\% \\
    public\_test  & $224\times224$ & 1.630 & 100\% \\
    private\_test & $224\times224$ & 2.859 & 100\% \\
    \bottomrule
  \end{tabular}
\end{table}

\paragraph{Quan sát trực quan} Để kiểm chứng bằng mắt, nhóm trích 6--8 khung hình phân bố đều trong video (\texttt{cap.set(cv2.CAP\_\allowbreak PROP\_\allowbreak POS\_\allowbreak FRAMES, i)} với $i$ lấy từ \texttt{np.linspace}) và ghép thành lưới ảnh cho một số lớp (Hình~\ref{fig:eda_samples}). Các video được quay chính diện, người ký ngồi hoặc đứng giữa khung hình, thấy rõ phần thân trên, hai tay và khuôn mặt; nền chủ yếu là tường trơn với ánh sáng trong nhà. Sự đa dạng nằm ở người ký (giới tính, trang phục, kính), màu nền (be, xanh nhạt, trắng) và vị trí tay lúc bắt đầu. Ký hiệu thường chỉ chiếm một phần thời lượng: các khung đầu và cuối thường là tư thế nghỉ. Ở một số video xuất hiện các khung \emph{đơn sắc hoàn toàn} (hàng \textit{Phỏng vấn}, \textit{Uống}); hiện tượng này được định lượng ở Mục~\ref{subsec:eda_quality}.

\begin{figure}[H]
  \centering
  \includegraphics[width=0.85\textwidth]{sample_frames_6classes.png}
  \caption{Khung hình mẫu (lấy đều theo thời gian) của sáu lớp. Hàng \textit{Phỏng vấn} và \textit{Uống} có các khung đơn sắc ở cuối video.}
  \label{fig:eda_samples}
\end{figure}

\noindent\textit{Hệ quả:} khi gán nhãn số cho video phải chuẩn hoá tên lớp bằng \texttt{unicodedata.normalize(\textquotesingle NFC\textquotesingle, ...)} trước khi tra \texttt{label\_mapping.pkl}, nếu không phần lớn video sẽ không tìm được nhãn. Chỉ có tập \texttt{train} mang nhãn nên tập kiểm định (\textit{validation}) phải tách ra từ chính tập này; việc tách cần tính đến trùng lặp (Mục~\ref{subsec:eda_quality}). Thông tin phân biệt nằm ở hình dạng và chuyển động của bàn tay, khẩu hình và tư thế, nên các phép tăng cường không được làm đổi ngữ nghĩa không gian; đặc biệt \emph{không} dùng lật ngang (\textit{Horizontal Flip}) vì có thể đổi nghĩa tay trái/phải.

% =====================================================================
\subsubsection{Chất lượng dữ liệu}
\label{subsec:eda_quality}

\paragraph{Script và cách thực hiện}
\begin{itemize}
  \item \textbf{Script:} bốn script trong \texttt{02\_data\_quality/}, chạy theo thứ tự: \edafile{02\_data\_quality}{01\_integrity\_missing.py}, \edafile{02\_data\_quality}{02\_duplicates.py}, \edafile{02\_data\_quality}{03\_frame\_scan.py}, \edafile{02\_data\_quality}{04\_frame\_quality\_summary.py}. Kiểm tra bằng mắt dùng \edafile{07\_visual\_inspection}{02\_flagged\_videos.py}.
  \item \textbf{Cách khai thác dữ liệu:} 
  
  (i) \emph{Giá trị thiếu và lỗi đọc:} \texttt{01\_integrity\_missing.py} đếm ô trống (\texttt{isna}) theo cột và theo tập, đếm video có 0 khung, FPS hoặc kích thước bằng 0, rồi liệt kê các tệp có giá trị phi lý theo ngưỡng tuyệt đối (dưới 10~KB, dưới 16 khung, FPS dưới 5). 
  
  (ii) \emph{Trùng lặp:} \texttt{02\_duplicates.py} đọc toàn bộ byte của từng tệp, tính \texttt{hashlib.md5} một lần và lưu vào \texttt{file\_hashes.csv}; hai tệp cùng mã băm là cùng nội dung. Gom nhóm theo mã băm để đếm bản sao, tìm nhóm có nhiều hơn một nhãn (xung đột nhãn), lấy giao các tập mã băm giữa train, public test, private test, rồi giữ một tệp cho mỗi mã băm và đếm lại số video mỗi lớp. 
  
  (iii) \emph{Khung đơn sắc:} \texttt{03\_frame\_scan.py} giải mã \emph{từng} khung bằng \texttt{cap.read()}, tính độ lệch chuẩn điểm ảnh của từng kênh màu và coi khung là đơn sắc khi giá trị lớn nhất nhỏ hơn 3; đồng thời đếm số khung đơn sắc liên tiếp ở đầu video và tính độ chuyển động là trung bình $|f_t-f_{t-1}|$ trên ảnh xám. \texttt{04\_frame\_quality\_summary.py} đối chiếu số khung đọc được với metadata, bỏ các bản sao \texttt{\_N}, phân loại vị trí khung đơn sắc (đầu, cuối, giữa, toàn bộ) và mô phỏng việc lấy $T=16$ khung bằng \texttt{np.linspace} để đếm số khung đơn sắc bị lấy trúng. 
  
  (iv) \emph{Kiểm tra bằng mắt:} \texttt{02\_flagged\_videos.py} \emph{tự chọn} video cần xem từ kết quả hai bước trên, không dùng danh sách chọn tay. Tiêu chí chọn gồm: 6 video có tỷ lệ khung đơn sắc cao nhất trong số các video có từ 50\% khung đơn sắc trở lên; mọi video thuộc nhóm MD5 nằm ở hơn một lớp; một ví dụ cho mỗi dạng xuất hiện khung đơn sắc. Ngoại lai theo quy tắc Tukey được tính ở bước đơn biến (Mục~\ref{subsec:eda_uni}).
  \item \textbf{Đầu ra:} \texttt{integrity\_summary.txt}, \texttt{file\_hashes.csv}, \texttt{dups\_summary.txt}, \texttt{dedup\_summary.txt}, \texttt{class\_counts\_train\_unique\_vs\_all.csv}, \texttt{video\_quality\_flags.csv}, \texttt{solid\_frame\_summary.txt}, \texttt{class\_solid\_frame\_ratio.csv}; ở bước kiểm tra bằng mắt là \texttt{flagged\_videos\_info.txt}, \texttt{flagged\_videos\_check.png}, \texttt{solid\_frame\_patterns.png}.
\end{itemize}

Nhóm kiểm tra chất lượng dữ liệu theo các khía cạnh: giá trị thiếu, bản ghi trùng lặp, khung hình đơn sắc/video rỗng, giá trị ngoại lai, nhiễu nhãn và trùng lặp giữa các tập. Bảng~\ref{tab:eda_quality} tóm tắt các vấn đề phát hiện được.

\begin{table}[H]
  \centering
  \caption{Tóm tắt các vấn đề chất lượng dữ liệu.}
  \label{tab:eda_quality}
  \small
  \begin{tabular}{@{}p{0.26\textwidth}p{0.36\textwidth}p{0.30\textwidth}@{}}
    \toprule
    \textbf{Vấn đề} & \textbf{Số lượng} & \textbf{Ví dụ} \\
    \midrule
    Bản sao hậu tố \texttt{\_N} trong train & 1.069 tệp giống hệt bản gốc; còn 2.776 video duy nhất & \texttt{127824.mp4}, \texttt{127824\_1.mp4}, \texttt{127824\_2.mp4} \\
    Video có khung đơn sắc (train) & 758/3.875 tệp; 565/2.776 video duy nhất (20,4\%), ở 78/100 lớp & \texttt{Thất lạc/708519}, \texttt{Bộ y tế/523389} \\
    Video rỗng hoàn toàn & 2 video gốc (4 tệp kể cả bản sao) & \texttt{Biết/214576}, \texttt{Khai báo/885339} \\
    Nhiễu nhãn (cùng nội dung, khác lớp) & 2 nhóm nội dung, 5 tệp & \textit{Chậm lại} $\leftrightarrow$ \textit{Chào} \\
    Trùng train $\rightarrow$ public test & 157 video public giống hệt video train & -- \\
    Trùng public $\leftrightarrow$ private & 896 video & -- \\
    \bottomrule
  \end{tabular}
\end{table}

\paragraph{Giá trị thiếu và lỗi đọc} Bảng metadata không có giá trị thiếu ở bất kỳ thuộc tính nào (nhãn \texttt{cls} chỉ trống ở hai tập test vì các tập này vốn không có nhãn). Mọi video đều mở và giải mã được, và ở cả 8.364 video, số khung giải mã thực tế khớp với số khung ghi trong metadata. Như vậy không có vấn đề về tệp hỏng hay nhãn bị thiếu ở tập train.

\paragraph{Trùng lặp trong tập train} Nhóm tính mã băm MD5 nội dung từng tệp (\texttt{hashlib.md5}) rồi gom nhóm theo mã băm. Có 1.069/3.875 tệp mang hậu tố \texttt{\_1}, \texttt{\_2}, \texttt{\_3}\ldots{} So sánh từng cặp cho thấy cả 1.069 tệp này đều có bản gốc tồn tại, cùng lớp và \emph{giống hệt bản gốc từng byte} (cùng số khung, FPS, dung lượng): đây là bản sao, không phải biến thể tăng cường. Có 176 video gốc có 1 bản sao, 169 video có 2 và 185 video có 3. Tính cả các bản sao khác, tập train có 1.099 tệp dư thừa và chỉ còn \textbf{2.776 video duy nhất}. Các bản sao tập trung ở những lớp vốn đã nhiều mẫu (\textit{Chậm lại} và \textit{Trường học} cùng thêm 25 tệp, \textit{Rau} và \textit{Nhìn} thêm 23, \textit{Khóc} thêm 22\ldots), nên sau khi khử trùng lặp số video mỗi lớp còn từ 5 đến 51 (trung bình 27,8; tỷ lệ chênh $\approx$10,2 lần). Bảng trước/sau khử trùng lặp của từng lớp lưu ở \texttt{class\_counts\_train\_unique\_vs\_all.csv}. Hệ quả là nguy cơ rò rỉ dữ liệu (\textit{Data Leakage}): nếu chia train/validation ngẫu nhiên theo từng tệp, các bản sao của cùng một video rơi vào cả hai phần, làm điểm validation bị thổi phồng.

\paragraph{Khung hình đơn sắc và video rỗng} Nhóm đọc \emph{toàn bộ} khung hình của 8.364 video và tính độ lệch chuẩn giá trị điểm ảnh trên từng kênh màu; khung có độ lệch chuẩn lớn nhất giữa ba kênh nhỏ hơn 3 được coi là \emph{khung đơn sắc}. Kết quả tổng hợp ở Bảng~\ref{tab:eda_solid}.

\begin{table}[H]
  \centering
  \caption{Số video chứa khung hình đơn sắc theo từng tập.}
  \label{tab:eda_solid}
  \begin{tabular}{@{}lrrrr@{}}
    \toprule
    \textbf{Tập} & \textbf{Tổng video} & \textbf{Có $\geq$1 khung đơn sắc} & \textbf{Đơn sắc toàn bộ} & \textbf{Tỷ lệ} \\
    \midrule
    train (tất cả tệp)        & 3.875 & 758 & 4 & 19,6\% \\
    train (video duy nhất)    & 2.776 & 565 & 2 & 20,4\% \\
    public\_test              & 1.630 & 1   & 0 & 0,06\% \\
    private\_test             & 2.859 & 1   & 0 & 0,03\% \\
    \bottomrule
  \end{tabular}
\end{table}

Đây là phát hiện đáng chú ý nhất của EDA. Khoảng một phần năm số video duy nhất trong train (565/2.776, thuộc 78/100 lớp) chứa khung đơn sắc, và số khung này không nhỏ: trung bình 12,2 khung trên khoảng 30 khung mỗi video (trung vị 12, nhiều nhất 23). Các khung có màu thuần (đỏ, xanh lục, xanh lam, xanh lục lam, tím, vàng, trắng, đen, xám) và xuất hiện theo các dạng: chỉ ở đầu video (231 video), chỉ ở cuối (161), cả đầu và cuối (116), xen ở giữa (55) hoặc toàn bộ video (2) (Hình~\ref{fig:eda_solid_examples}). Lớp có tỷ lệ video bị ảnh hưởng cao nhất là \textit{Có thể} (6/15 video duy nhất), \textit{Xin lỗi} (5/14) và \textit{Thất lạc} (7/20) (chi tiết ở \texttt{class\_solid\_frame\_ratio.csv}).

\begin{figure}[H]
  \centering
  \includegraphics[width=\textwidth]{solid_frame_patterns.png}
  \caption{Bốn dạng khung đơn sắc trong tập train (khung chẵn $f_0$--$f_{28}$), mỗi dạng một video được chọn tự động: chỉ ở đầu (\textit{An ủi}/133326), chỉ ở cuối (\textit{Ban ngày}/363626), xen ở giữa (\textit{Ban ngày}/725380), cả đầu và cuối (\textit{Ban ngày}/855731).}
  \label{fig:eda_solid_examples}
\end{figure}

Hiện tượng này ảnh hưởng trực tiếp đến cách lấy mẫu: khi lấy $T=16$ khung đều nhau bằng \texttt{linspace} trên 565 video nói trên, trung bình có 6,4/16 khung được chọn là khung đơn sắc, và ở 224 video có từ 8/16 khung trở lên (ít nhất một nửa chuỗi đầu vào) là khung đơn sắc. Ngược lại, mỗi tập test chỉ có 1 video chứa đúng 1 khung đơn sắc. Như vậy khung đơn sắc là đặc điểm \emph{riêng của tập train}: nếu không xử lý, mô hình vừa mất thông tin ký hiệu, vừa có thể học các tín hiệu giả không tồn tại ở dữ liệu kiểm tra.

Sáu hàng đầu của Hình~\ref{fig:eda_suspicious} là các video có tỷ lệ khung đơn sắc cao nhất. Hai video được xác nhận bằng mắt là \textbf{rỗng hoàn toàn}: \texttt{Biết/214576.mp4} (2,5~KB, 16 khung, toàn màu xanh lục) và \texttt{Khai báo/885339.mp4} (2,7~KB, 23 khung, toàn màu xanh lam, kèm 2 bản sao \texttt{\_1}, \texttt{\_2}); cả hai không có người ký. Bốn video tiếp theo có 20--21/30 khung đơn sắc: \texttt{Nôn ói/915809}, \texttt{Thất lạc/708519}, \texttt{Chân/984642} và \texttt{Khai báo/362529}. Người ký chỉ xuất hiện trong khoảng 10 khung, ví dụ \texttt{Khai báo/362529} có 20 khung đỏ đầu tiên ($f_0$--$f_{19}$).

\paragraph{Giá trị ngoại lai} Nhóm dùng quy tắc Tukey (ngoài khoảng $[Q_1-1{,}5\,\mathrm{IQR},\,Q_3+1{,}5\,\mathrm{IQR}]$ tính trên tập train) cho thời lượng, dung lượng và FPS. Số lượng ngoại lai khá nhiều nhưng phần lớn chỉ lệch nhẹ (Mục~\ref{subsec:eda_uni}). Đáng chú ý là các ngoại lai cực đoan: thời lượng trên 5~s có 10 video train, 1 video public test và 3 video private test (lớn nhất 35,3~s ở train, 28,1~s ở private test); FPS dưới 5 có 9 video train và 2 video private test; dung lượng dưới 10~KB có 4 tệp train, đều là các video rỗng nói trên. Các ngoại lai này không nhất thiết là lỗi mà phản ánh việc video được lấy mẫu lại về $\approx$30 khung từ các clip gốc có độ dài rất khác nhau.

\paragraph{Nhiễu nhãn} Khi gom nhóm theo MD5, có hai nhóm nội dung giống hệt nhau nhưng nằm ở \emph{hai lớp khác nhau}: \texttt{Chậm lại/284430.mp4} trùng với \texttt{Chào/392683.mp4} và \texttt{Chào/468203.mp4}; \texttt{Chậm lại/254739.mp4} trùng với \texttt{Chào/792094.mp4}. Nhóm đã xem trực tiếp (Hình~\ref{fig:eda_suspicious}, năm hàng cuối) và xác nhận chúng có cùng người ký, cùng chuỗi động tác. Một video không thể đồng thời mang hai nghĩa nên ít nhất một trong hai nhãn là sai; tuy nhiên việc xác định nhãn đúng cần người am hiểu ngôn ngữ ký hiệu Việt Nam, nên báo cáo chỉ ghi nhận đây là \emph{xung đột nhãn}.

\begin{figure}[H]
  \centering
  \includegraphics[width=0.9\textwidth,height=0.72\textheight,keepaspectratio]{flagged_videos_check.png}
  \caption{Kiểm tra bằng mắt các video do bước kiểm tra chất lượng tự đánh dấu. Sáu hàng đầu là các video có tỷ lệ khung đơn sắc cao nhất, gồm hai video rỗng hoàn toàn (\textit{Biết}/214576, \textit{Khai báo}/885339). Năm hàng cuối là hai nhóm video trùng nội dung nhưng khác nhãn (\textit{Chậm lại} $\leftrightarrow$ \textit{Chào}).}
  \label{fig:eda_suspicious}
\end{figure}

\paragraph{Trùng lặp giữa các tập} Public test có 8 bản trùng nội bộ (1.622 video duy nhất) và \textbf{157 video giống hệt video trong train}; private test có 2.782 video duy nhất, không trùng với train, nhưng có 896 video trùng với public test. Ngoài ra, 2.806 tên tệp trong private test trùng với tên tệp trong train nhưng nội dung hoàn toàn khác (MD5 khác, chênh lệch điểm ảnh lớn), nên đây chỉ là trùng cách đặt mã số; nguyên nhân chưa rõ và tên tệp \emph{không} được dùng để suy luận nhãn.

\noindent\textit{Hệ quả:} dữ liệu train chưa sạch; trước khi huấn luyện cần khử trùng lặp, loại video rỗng, xử lý khung đơn sắc và rà soát các nhãn xung đột. Điểm trên public test có thể bị ảnh hưởng một phần bởi 157 video đã có trong train.

% =====================================================================
\subsubsection{Phân tích đơn biến}
\label{subsec:eda_uni}

\paragraph{Script và cách thực hiện}
\begin{itemize}
  \item \textbf{Script:} \edafile{03\_univariate}{01\_numeric\_attributes.py}.
  \item \textbf{Cách khai thác dữ liệu:} với mỗi thuộc tính số (\texttt{frames}, \texttt{fps}, \texttt{dur}, \texttt{size\_kb}), script tính \texttt{groupby(\textquotesingle split\textquotesingle).describe()} (trung bình, độ lệch chuẩn, min, các tứ phân vị, max), hệ số lệch (\textit{skewness}) và các giá trị phổ biến nhất. Ngoại lai được xác định theo quy tắc Tukey với ngưỡng tính trên train rồi áp lên cả ba tập. Script vẽ histogram của từng thuộc tính cho ba tập trên cùng trục; histogram FPS và dung lượng dùng mật độ (\texttt{density=True}) để so sánh được các tập có kích thước khác nhau.
  \item \textbf{Đầu ra:} \texttt{desc\_by\_split.csv} (Bảng~\ref{tab:eda_desc}), \texttt{univariate\_summary.txt}, \texttt{frames\_hist\_by\_split.png}, \texttt{fps\_hist\_by\_split.png}, \texttt{filesize\_hist\_by\_split.png}.
\end{itemize}

Mục này khảo sát phân phối của từng thuộc tính số: số khung hình, FPS, thời lượng và dung lượng, tính trên từng tập (Bảng~\ref{tab:eda_desc}). Độ phân giải đã được trình bày ở Bảng~\ref{tab:eda_res}.

\begin{table}[H]
  \centering
  \caption{Thống kê mô tả các thuộc tính số theo từng tập.}
  \label{tab:eda_desc}
  \small
  \begin{tabular}{@{}llrrrrrrr@{}}
    \toprule
    \textbf{Thuộc tính} & \textbf{Tập} & \textbf{Mean} & \textbf{Std} & \textbf{Min} & \textbf{Q1} & \textbf{Median} & \textbf{Q3} & \textbf{Max} \\
    \midrule
    \multirow{3}{*}{Số khung}
      & train   & 29,88 & 1,26 & 16 & 30 & 30 & 30 & 30 \\
      & public  & 29,96 & 0,66 & 16 & 30 & 30 & 30 & 30 \\
      & private & 29,98 & 0,44 & 21 & 30 & 30 & 30 & 30 \\
    \midrule
    \multirow{3}{*}{FPS}
      & train   & 12,94 & 3,11 & 0,85 & 10,98 & 12,34 & 14,06 & 30,02 \\
      & public  & 12,97 & 3,00 & 5,93 & 10,98 & 12,50 & 13,85 & 30,00 \\
      & private & 13,32 & 3,69 & 1,07 & 10,84 & 12,50 & 14,06 & 30,02 \\
    \midrule
    \multirow{3}{*}{Thời lượng (s)}
      & train   & 2,48 & 1,36 & 0,70 & 2,13 & 2,43 & 2,73 & 35,29 \\
      & public  & 2,42 & 0,51 & 0,70 & 2,17 & 2,40 & 2,73 & 5,06 \\
      & private & 2,42 & 0,82 & 0,70 & 2,13 & 2,40 & 2,77 & 28,14 \\
    \midrule
    \multirow{3}{*}{Dung lượng (KB)}
      & train   & 85,5  & 38,5 & 2,5  & 60,2 & 80,6 & 103,0 & 263,8 \\
      & public  & 113,6 & 77,5 & 26,2 & 76,1 & 91,0 & 125,8 & 483,2 \\
      & private & 116,3 & 77,7 & 42,2 & 77,7 & 94,6 & 126,0 & 488,0 \\
    \bottomrule
  \end{tabular}
\end{table}

\paragraph{Số khung hình} Hầu hết video có đúng 30 khung: 3.832/3.875 video train (98,9\%), 1.623/1.630 public test và 2.850/2.859 private test; ở train có 26 video chỉ có 16 khung. Không video nào ngắn hơn 16 khung (Hình~\ref{fig:eda_frames_hist}), nên không cần đệm khung thực sự khi lấy $T=16$.

\begin{figure}[H]
  \centering
  \includegraphics[width=\textwidth]{frames_hist_by_split.png}
  \caption{Histogram số khung hình mỗi video theo tập (trục tung thang log). Đường đứt nét đánh dấu $T=16$.}
  \label{fig:eda_frames_hist}
\end{figure}

\paragraph{FPS và thời lượng} Dù số khung gần như cố định, FPS biến thiên mạnh: train có 961 giá trị FPS khác nhau, từ 0,85 đến 30,02 (trung vị 12,3), phân phối lệch phải với đuôi dài (hệ số lệch 1,01 ở train; Hình~\ref{fig:eda_fps_hist}). Vì thời lượng $=\texttt{frames}/\texttt{fps}$ nên thời lượng thực của cùng 30 khung dao động lớn: trung vị khoảng 2,4~s nhưng có video tới 35,3~s ở train và 28,1~s ở private test, trong khi 99\% video đều ngắn hơn khoảng 4~s. Điều này cho thấy video đã được lấy mẫu lại về khoảng 30 khung bất kể độ dài clip gốc.

\begin{figure}[H]
  \centering
  \includegraphics[width=0.85\textwidth]{fps_hist_by_split.png}
  \caption{Phân phối FPS theo từng tập.}
  \label{fig:eda_fps_hist}
\end{figure}

\paragraph{Dung lượng tệp} Dung lượng train có trung vị 80,6~KB (trung bình 85,5; từ 2,5 đến 263,8~KB). Hai tập test có phân phối lệch phải hơn hẳn: trung bình 113,6--116,3~KB, lớn nhất tới 483--488~KB (Hình~\ref{fig:eda_size_hist}); khác biệt này được phân tích ở Mục~\ref{subsec:eda_shift}.

\begin{figure}[H]
  \centering
  \includegraphics[width=0.85\textwidth]{filesize_hist_by_split.png}
  \caption{Phân phối dung lượng tệp (KB) theo từng tập.}
  \label{fig:eda_size_hist}
\end{figure}

\noindent\textit{Hệ quả:} (i) không cần đổi kích thước khung hình, tuy kiến trúc CRNN với \texttt{Adaptive\allowbreak AvgPool2d} vẫn đảm bảo tính tổng quát nếu gặp video khác độ phân giải; (ii) vì mọi video có ít nhất 16 khung, lấy mẫu $T=16$ khung đều nhau bằng \texttt{linspace} (\textit{Frame Downsampling}) là đủ, lặp khung (\textit{repeat padding}) chỉ giữ làm phương án dự phòng; (iii) phân phối FPS và dung lượng lệch phải với nhiều ngoại lai nên nên dùng trung vị thay vì trung bình khi mô tả.

% =====================================================================
\subsubsection{Phân tích biến mục tiêu}
\label{subsec:eda_class}

\paragraph{Script và cách thực hiện}
\begin{itemize}
  \item \textbf{Script:} \edafile{04\_target\_analysis}{01\_class\_distribution.py}.
  \item \textbf{Cách khai thác dữ liệu:} lọc các dòng của tập train trong \texttt{video\_metadata.csv}, đếm số video mỗi nhãn bằng \texttt{value\_counts()}, sau đó tính trung bình, trung vị, độ lệch chuẩn, tỷ lệ lớp lớn nhất/nhỏ nhất và số lớp có không quá 25 hoặc từ 50 video trở lên. Số video mỗi lớp sau khử trùng lặp lấy từ \texttt{class\_counts\_train\_unique\_vs\_all.csv} của bước \texttt{02\_data\_quality}.
  \item \textbf{Đầu ra:} \texttt{class\_summary.txt}, \texttt{class\_counts\_train.csv}, \texttt{class\_distribution\_all.png}, \texttt{class\_top\_bottom10.png}, \texttt{samples\_per\_class\_hist.png}.
\end{itemize}

Nhóm dùng \texttt{value\_counts()} trên cột nhãn của tập \texttt{train} để đếm số video mỗi lớp, rồi vẽ biểu đồ cho cả 100 lớp (Hình~\ref{fig:eda_class_dist}), 10 lớp nhiều nhất/ít nhất (Hình~\ref{fig:eda_top_bottom}) và histogram số video mỗi lớp (Hình~\ref{fig:eda_spc_hist}).

\begin{figure}[H]
  \centering
  \includegraphics[width=\textwidth]{class_distribution_all.png}
  \caption{Phân phối số video theo lớp trong tập huấn luyện (100 lớp, sắp xếp giảm dần).}
  \label{fig:eda_class_dist}
\end{figure}

Tập huấn luyện gồm 3.875 video nhưng phân bố không đồng đều: trung bình 38,75 video/lớp, trung vị 45,5, độ lệch chuẩn 22,78. Lớp nhiều nhất (\textit{Trường học}) có 74 video, trong khi mười lớp chỉ có 6--7 video; tỷ lệ chênh giữa lớp lớn nhất và nhỏ nhất xấp xỉ $12{,}3$ lần (Bảng~\ref{tab:eda_top_bottom}). Có 36/100 lớp có không quá 25 video và 42/100 lớp có từ 50 video trở lên, tức phân phối có dạng hai đỉnh: nhóm lớp ``giàu'' (khoảng 45--74 video) và nhóm lớp ``nghèo'' (khoảng 6--11 video). Nhiều lớp ít mẫu thuộc nhóm từ vựng y tế -- dịch bệnh.

\begin{table}[H]
  \centering
  \caption{Mười lớp nhiều mẫu nhất và mười lớp ít mẫu nhất trong tập huấn luyện.}
  \label{tab:eda_top_bottom}
  \begin{tabular}{@{}lr@{\hspace{2.5em}}lr@{}}
    \toprule
    \multicolumn{2}{c}{\textbf{Nhiều nhất}} & \multicolumn{2}{c}{\textbf{Ít nhất}} \\
    \cmidrule(r){1-2}\cmidrule(l){3-4}
    \textbf{Lớp} & \textbf{Số video} & \textbf{Lớp} & \textbf{Số video} \\
    \midrule
    Trường học & 74 & Nhầm      & 7 \\
    Rau        & 72 & Virus     & 7 \\
    Nhà        & 71 & Bệnh nhân & 6 \\
    Áp dụng    & 71 & Cách ly   & 6 \\
    Chậm lại   & 69 & Cơ thể    & 6 \\
    Cám dỗ     & 68 & Ho        & 6 \\
    Nghe       & 68 & Phục hồi  & 6 \\
    Nói xấu    & 68 & Sốt       & 6 \\
    Biết       & 67 & Xe máy    & 6 \\
    Chạy       & 67 & Ủng hộ    & 6 \\
    \bottomrule
  \end{tabular}
\end{table}

\begin{figure}[H]
  \centering
  \begin{minipage}{0.56\textwidth}
    \centering
    \includegraphics[width=\textwidth]{class_top_bottom10.png}
    \caption{Mười lớp nhiều nhất (xanh) và ít nhất (đỏ).}
    \label{fig:eda_top_bottom}
  \end{minipage}\hfill
  \begin{minipage}{0.42\textwidth}
    \centering
    \includegraphics[width=\textwidth]{samples_per_class_hist.png}
    \caption{Histogram số video mỗi lớp.}
    \label{fig:eda_spc_hist}
  \end{minipage}
\end{figure}

Các con số trên tính trên tệp gốc, \emph{chưa} khử trùng lặp.

\noindent\textit{Hệ quả:} dữ liệu bị mất cân bằng lớp (\textit{Class Imbalance}) rõ rệt. Nếu đánh giá bằng Accuracy, mô hình có thể thiên về lớp phổ biến mà vẫn đạt điểm cao; vì vậy dùng Macro F1 (cũng là thước đo chính thức) và cần biện pháp xử lý mất cân bằng ở bước tiền xử lý/huấn luyện.

% =====================================================================
\subsubsection{Phân tích hai biến và đa biến}
\label{subsec:eda_multi}

\paragraph{Script và cách thực hiện}
\begin{itemize}
  \item \textbf{Script:} \edafile{05\_bivariate\_multivariate}{01\_correlation\_and\_label\_association.py}.
  \item \textbf{Cách khai thác dữ liệu:} ghép \texttt{video\_metadata.csv} với \texttt{video\_quality\_flags.csv} theo khoá (tập, lớp, tên tệp); bỏ các bản sao \texttt{\_N} để chỉ giữ 2.776 video duy nhất của train. Tính ma trận tương quan Pearson và Spearman (tương quan hạng, ít nhạy với ngoại lai) giữa các thuộc tính số; với từng thuộc tính, chia giá trị theo 100 lớp rồi chạy kiểm định Kruskal--Wallis (\texttt{scipy.stats.kruskal}) và tính $\eta^2 = SS_{\text{giữa lớp}}/SS_{\text{tổng}}$. Trung vị và tứ phân vị FPS của từng lớp được sắp xếp tăng dần để vẽ.
  \item \textbf{Đầu ra:} \texttt{multivariate\_summary.txt}, \texttt{corr\_heatmap.png}, \texttt{fps\_by\_class.png}, \texttt{size\_vs\_solid.png}.
\end{itemize}

Vì dữ liệu video không có đặc trưng dạng bảng, phân tích hai biến và đa biến tập trung vào mối quan hệ giữa các thuộc tính metadata với nhau và với nhãn lớp. Để tránh bản sao làm lệch kết quả, các phép tính dưới đây dùng 2.776 video duy nhất của tập train.

\paragraph{Tương quan giữa các thuộc tính} Hình~\ref{fig:eda_corr} cho thấy hệ số tương quan Pearson giữa số khung, FPS, dung lượng, tỷ lệ khung đơn sắc và độ chuyển động. (Thời lượng không đưa vào vì $\texttt{dur}=\texttt{frames}/\texttt{fps}$ nên phụ thuộc hoàn toàn vào hai biến còn lại.) Hầu hết các cặp có tương quan yếu ($|r|<0{,}35$); hai quan hệ rõ nhất là dung lượng -- tỷ lệ khung đơn sắc ($r\approx-0{,}45$; Spearman $\rho\approx-0{,}51$) và FPS -- chuyển động ($r\approx-0{,}34$; $\rho\approx-0{,}30$). Tương quan Spearman cho cùng chiều và độ lớn tương tự, nên các quan hệ này không phải do ngoại lai tạo ra. Quan hệ đầu là hợp lý: khung đơn sắc nén rất nhỏ nên video chứa nhiều khung đơn sắc có dung lượng thấp (Hình~\ref{fig:eda_size_solid}). Quan hệ sau cho thấy video có FPS thấp (clip gốc dài hơn) có độ chuyển động giữa các khung liên tiếp lớn hơn. Số khung gần như không liên quan tới các biến khác vì hầu như không đổi.

\begin{figure}[H]
  \centering
  \begin{minipage}{0.48\textwidth}
    \centering
    \includegraphics[width=\textwidth]{corr_heatmap.png}
    \caption{Ma trận tương quan Pearson giữa các thuộc tính (train, video duy nhất).}
    \label{fig:eda_corr}
  \end{minipage}\hfill
  \begin{minipage}{0.48\textwidth}
    \centering
    \includegraphics[width=\textwidth]{size_vs_solid.png}
    \caption{Dung lượng tệp theo tỷ lệ khung đơn sắc.}
    \label{fig:eda_size_solid}
  \end{minipage}
\end{figure}

\paragraph{Quan hệ giữa thuộc tính và nhãn lớp} Nhóm so sánh phân phối của từng thuộc tính giữa 100 lớp bằng kiểm định Kruskal--Wallis và hệ số $\eta^2$ (tỷ lệ phương sai giải thích bởi nhãn). Kết quả: FPS có $\eta^2\approx0{,}28$, dung lượng $\eta^2\approx0{,}15$, số khung $\eta^2\approx0{,}08$; cả ba đều có $p<10^{-10}$. Nói cách khác, FPS của video khác nhau rõ rệt giữa các lớp: trung vị FPS của lớp dao động từ khoảng 9,4 (\textit{Ngón tay}) đến khoảng 19,2 (\textit{Ho}), trong khi trung vị toàn tập là 12,4 (Hình~\ref{fig:eda_fps_class}). Điều này gợi ý rằng thời lượng clip gốc (và do đó FPS sau khi lấy mẫu lại) có liên hệ với từ được ký -- có thể vì các từ khác nhau có độ dài tự nhiên khác nhau -- hoặc vì cách thu thập dữ liệu theo từng lớp. Nhóm chưa xác định được nguyên nhân.

\begin{figure}[H]
  \centering
  \includegraphics[width=0.95\textwidth]{fps_by_class.png}
  \caption{Trung vị và khoảng tứ phân vị của FPS theo lớp (train, video duy nhất), các lớp sắp xếp theo trung vị.}
  \label{fig:eda_fps_class}
\end{figure}

\paragraph{Khung đơn sắc theo lớp} Khung đơn sắc không phân bố đều: 78/100 lớp có ít nhất một video bị ảnh hưởng, và có lớp có tới 30--40\% video duy nhất chứa khung đơn sắc (\textit{Có thể}, \textit{Xin lỗi}, \textit{Thất lạc}, \textit{Mua}). Do đó việc loại bỏ khung đơn sắc sẽ tác động không đồng đều lên các lớp.

\noindent\textit{Hệ quả:} (i) các thuộc tính metadata phần lớn độc lập nhau nên không có cặp đặc trưng dư thừa cần loại bỏ; (ii) vì FPS và dung lượng có quan hệ thống kê với nhãn, chúng có thể trở thành ``đường tắt'' (\textit{shortcut}) mà mô hình học được thay vì học ký hiệu, nên không nên đưa FPS hoặc dung lượng làm đặc trưng đầu vào và cần theo dõi nguy cơ này khi đánh giá; (iii) cần kiểm tra lại phân phối nhãn sau khi loại khung đơn sắc.

% =====================================================================
\subsubsection{So sánh phân phối giữa các tập}
\label{subsec:eda_shift}

\paragraph{Script và cách thực hiện}
\begin{itemize}
  \item \textbf{Script:} \edafile{06\_split\_comparison}{01\_split\_shift.py}.
  \item \textbf{Cách khai thác dữ liệu:} với train chỉ giữ video duy nhất, để các bản sao không làm lệch phân phối. Với mỗi thuộc tính (\texttt{frames}, \texttt{fps}, \texttt{dur}, \texttt{size\_kb}, \texttt{motion}), script tính trung vị từng tập và chạy kiểm định Kolmogorov--Smirnov hai mẫu (\texttt{scipy.stats.ks\_2samp}) cho từng cặp tập. Thống kê $D$ là khoảng cách lớn nhất giữa hai hàm phân phối tích luỹ thực nghiệm (ECDF); $D$ càng lớn thì hai phân phối càng khác nhau. Script cũng tính tỷ lệ video có khung đơn sắc ở từng tập và số video trùng nội dung giữa các tập (theo \texttt{file\_hashes.csv}). Ngưỡng ngoại lai trên của dung lượng (167~KB, tính từ train) lấy từ bước đơn biến.
  \item \textbf{Đầu ra:} \texttt{split\_shift\_summary.txt}, \texttt{ecdf\_by\_split.png}.
\end{itemize}

Nhóm so sánh phân phối các thuộc tính giữa ba tập trên cùng trục (Bảng~\ref{tab:eda_desc}, Bảng~\ref{tab:eda_ks}, Hình~\ref{fig:eda_ecdf}) để tìm dấu hiệu khác biệt miền (\textit{Domain Shift}).

\begin{table}[H]
  \centering
  \caption{Kiểm định Kolmogorov--Smirnov hai mẫu giữa các tập (train: video duy nhất). Giá trị là thống kê $D$, trong ngoặc là $p$-value.}
  \label{tab:eda_ks}
  \small
  \begin{tabular}{@{}lccc@{}}
    \toprule
    \textbf{Thuộc tính} & \textbf{train -- public} & \textbf{train -- private} & \textbf{public -- private} \\
    \midrule
    Số khung        & 0,005 (1)                 & 0,007 (1)                 & 0,002 (1) \\
    FPS             & 0,075 ($1{,}7\cdot10^{-5}$) & 0,069 ($2{,}9\cdot10^{-6}$) & 0,094 ($1{,}7\cdot10^{-8}$) \\
    Thời lượng      & 0,075 ($1{,}7\cdot10^{-5}$) & 0,068 ($4{,}5\cdot10^{-6}$) & 0,094 ($1{,}7\cdot10^{-8}$) \\
    Dung lượng      & \textbf{0,230} ($2\cdot10^{-48}$) & \textbf{0,237} ($4\cdot10^{-70}$) & 0,059 (0,0013) \\
    Chuyển động     & 0,059 (0,0015)            & 0,079 ($4{,}8\cdot10^{-8}$) & 0,035 (0,15) \\
    \bottomrule
  \end{tabular}
\end{table}

\begin{figure}[H]
  \centering
  \includegraphics[width=\textwidth]{ecdf_by_split.png}
  \caption{Hàm phân phối tích luỹ thực nghiệm của FPS, dung lượng và thời lượng theo tập (train: video duy nhất). Đường dung lượng của train tách hẳn khỏi hai tập test.}
  \label{fig:eda_ecdf}
\end{figure}

Vì cỡ mẫu lớn nên mọi khác biệt dù nhỏ đều có $p$ rất bé; nhóm đánh giá chủ yếu theo độ lớn của $D$. Số khung hình giống nhau ở cả ba tập ($D<0{,}01$), độ phân giải cũng vậy. FPS, thời lượng và độ chuyển động khác biệt nhỏ ($D<0{,}1$). Riêng dung lượng khác biệt rõ giữa train và hai tập test ($D\approx0{,}23$--$0{,}24$), trong khi hai tập test gần nhau ($D\approx0{,}06$). Cụ thể:

\begin{itemize}
  \item \textbf{Dung lượng:} video test nặng hơn đáng kể (trung bình 113,6--116,3~KB so với 85,5~KB; trung vị 91,0--94,6~KB so với 80,6~KB). Tỷ lệ video vượt 167~KB (ngưỡng ngoại lai của train) là 3,1\% ở train nhưng 6,0\% ở public và 8,2\% ở private. Một phần chênh lệch có thể giải thích bởi khung đơn sắc (nén rất nhỏ) chỉ có ở train: trung vị dung lượng của video train có khung đơn sắc là 49,0~KB, còn của video train không có khung đơn sắc là 87,6~KB, gần mức của test. Tuy nhiên đuôi dài tới 488~KB ở test không xuất hiện ở train (lớn nhất 263,8~KB), cho thấy test có thể chứa các video chi tiết hoặc nhiễu hơn.
  \item \textbf{FPS:} trung vị gần như nhau (12,4 ở train, 12,5 ở hai tập test), nhưng private test có độ lệch chuẩn lớn hơn (3,69 so với 3,11 ở train) và đuôi dài hơn.
  \item \textbf{Khung đơn sắc:} khoảng 20\% video train nhưng gần như không có ở test (Bảng~\ref{tab:eda_solid}).
  \item \textbf{Người ký, nền, ánh sáng:} có thể khác giữa train và test; nhóm chưa định lượng được yếu tố này.
\end{itemize}

\noindent\textit{Hệ quả:} cần các phép tăng cường giúp mô hình bền vững với thay đổi màu sắc, độ sáng và khung cắt (\textit{ColorJitter}, \textit{Random Resized Crop}); việc loại bỏ khung đơn sắc ở train đồng thời giúp thu hẹp khoảng cách giữa phân phối train và test.

% =====================================================================
\subsubsection{Kết luận EDA và khuyến nghị cho bước tiền xử lý}
\label{subsec:eda_conclusion}

Bảng~\ref{tab:eda_reco} tổng hợp các vấn đề phát hiện được, bằng chứng và khuyến nghị cho bước tiền xử lý (Mục~\ref{sec:preprocess}) cũng như huấn luyện.

\begin{table}[H]
  \centering
  \caption{Tổng hợp vấn đề, bằng chứng và khuyến nghị cho bước tiền xử lý.}
  \label{tab:eda_reco}
  \small
  \begin{tabular}{@{}p{0.2\textwidth}p{0.36\textwidth}p{0.38\textwidth}@{}}
    \toprule
    \textbf{Vấn đề} & \textbf{Bằng chứng} & \textbf{Khuyến nghị} \\
    \midrule
    Tên lớp khác dạng Unicode & Khoá trong \texttt{label\_mapping.pkl} ở dạng NFD, thư mục lớp ở dạng NFC; 91/100 lớp không khớp nếu so nguyên văn & Chuẩn hoá NFC cả hai phía trước khi tra nhãn \\
    \addlinespace
    Video trùng lặp (\texttt{\_N}) & 1.069 tệp \texttt{\_N} giống hệt bản gốc; chỉ còn 2.776 video duy nhất & Khử trùng lặp theo MD5; chia train/validation theo ID gốc (bỏ hậu tố \texttt{\_N}) bằng \texttt{StratifiedGroupKFold} để tránh rò rỉ \\
    \addlinespace
    Mất cân bằng lớp & 6--74 video/lớp (sau khử trùng lặp: 5--51), 36/100 lớp có $\leq$25 video & Dùng Macro F1; trọng số lớp hoặc \texttt{WeightedRandomSampler}; oversample kèm tăng cường cho lớp ít mẫu; có thể dùng label smoothing \\
    \addlinespace
    Khung đơn sắc & 565/2.776 video duy nhất (20,4\%) chứa trung bình 12,2 khung đơn sắc; test gần như không có & Loại bỏ khung đơn sắc (độ lệch chuẩn điểm ảnh $<3$) \emph{trước} khi lấy mẫu $T=16$; nếu còn ít hơn 16 khung hợp lệ thì lặp khung \\
    \addlinespace
    Video rỗng & \texttt{Biết/214576}, \texttt{Khai báo/885339} (cùng 2 bản sao) đơn sắc toàn bộ & Loại khỏi tập huấn luyện \\
    \addlinespace
    Xung đột nhãn & 2 nhóm nội dung giống hệt nhưng thuộc \textit{Chậm lại} và \textit{Chào} & Xác minh nhãn đúng; nếu không xác minh được thì loại bỏ \\
    \addlinespace
    Độ dài không đồng nhất & 98,9\% video có 30 khung, tối thiểu 16 khung; thời lượng tới 35,3~s & Lấy mẫu $T=16$ khung đều bằng \texttt{linspace} \\
    \addlinespace
    FPS/dung lượng liên hệ với nhãn & $\eta^2\approx0{,}28$ (FPS), $0{,}15$ (dung lượng) giữa các lớp & Không dùng FPS, dung lượng làm đặc trưng; theo dõi nguy cơ shortcut \\
    \addlinespace
    Khác biệt miền train--test & Test nặng hơn (trung bình 114--116~KB so với 85,5~KB; KS $D\approx0{,}23$); khung đơn sắc chỉ có ở train & ColorJitter, Random Resized Crop, chuẩn hoá ImageNet; \emph{không} lật ngang \\
    \addlinespace
    Trùng train--public test & 157 video public giống hệt video train & Lưu ý khi diễn giải điểm public test \\
    \bottomrule
  \end{tabular}
\end{table}

Tóm lại, bộ dữ liệu có định dạng rất đồng nhất (100\% video $224\times224$, gần như đều 30 khung), nên khâu đọc và lấy mẫu khung đơn giản. Tuy nhiên dữ liệu huấn luyện có ba vấn đề cần xử lý trước khi huấn luyện: mất cân bằng lớp, khoảng 28\% tệp là bản sao gây nguy cơ rò rỉ dữ liệu, và khoảng 20\% video chứa khung đơn sắc vốn không xuất hiện ở dữ liệu kiểm tra. Đây là những yếu tố có thể khiến điểm validation cao hơn đáng kể so với điểm thực tế trên tập kiểm tra nếu không được xử lý.

\paragraph{Hạn chế của phân tích} FPS và thời lượng đọc từ metadata của container nên có thể không phản ánh tốc độ quay thực. Ngưỡng độ lệch chuẩn $<3$ chỉ phát hiện khung đơn sắc hoàn toàn, chưa phát hiện khung mờ hoặc chỉ có nền. Nhóm chưa định lượng sự khác biệt về người ký và bối cảnh giữa train và test, chưa xác định được nguyên nhân FPS liên hệ với nhãn, và chưa xác định được nhãn đúng cho các video xung đột.
]
% Cấu trúc theo khung EDA chuẩn:
%   3.1.1 Mục tiêu và tổng quan dữ liệu
%   3.1.2 Cấu trúc dữ liệu
%   3.1.3 Chất lượng dữ liệu
%   3.1.4 Phân tích đơn biến
%   3.1.5 Phân tích biến mục tiêu
%   3.1.6 Phân tích hai biến và đa biến
%   3.1.7 So sánh phân phối giữa các tập
%   3.1.8 Kết luận và khuyến nghị
% Gói cần: graphicx, booktabs, float, caption, amsmath
% =====================================================================

\subsection{Phân tích dữ liệu (EDA)}
\label{sec:eda}

% \edafile{thư mục}{tệp}: in đường dẫn script, cho phép xuống dòng sau dấu /
\newcommand{\edafile}[2]{\texttt{#1/}\allowbreak\texttt{#2}}
% Cho phép xuống dòng sau dấu gạch dưới trong tên tệp dài (vd. class\_counts\_train\_unique\_vs\_all.csv)
\let\edaOrigUS\_
\renewcommand{\_}{\edaOrigUS\allowbreak}
\sloppy

Mục này trình bày kết quả phân tích khám phá dữ liệu (\textit{Exploratory Data Analysis -- EDA}) trên bộ dữ liệu video ngôn ngữ ký hiệu. EDA chỉ nhằm \emph{quan sát} xu hướng, phân phối và chất lượng dữ liệu. Nhóm sử dụng Python làm công cụ để thực hiện EDA (\texttt{os}, \texttt{pickle}, \texttt{pandas}, \texttt{OpenCV}, \texttt{numpy}, \texttt{hashlib}, \texttt{matplotlib}).

Phần phân tích được tổ chức theo trình tự quy trình EDA: tổng quan và cấu trúc dữ liệu (Mục~\ref{subsec:eda_overview}--\ref{subsec:eda_structure}), chất lượng dữ liệu (Mục~\ref{subsec:eda_quality}), phân tích đơn biến (Mục~\ref{subsec:eda_uni}), phân tích biến mục tiêu (Mục~\ref{subsec:eda_class}), phân tích hai biến và đa biến (Mục~\ref{subsec:eda_multi}), so sánh giữa các tập (Mục~\ref{subsec:eda_shift}) và cuối cùng là kết luận kèm khuyến nghị (Mục~\ref{subsec:eda_conclusion}).

\begin{table}[H]
  \centering
  \caption{Các script EDA, thứ tự chạy và mục sử dụng kết quả.}
  \label{tab:eda_scripts}
  \small
  \begin{tabular}{@{}p{0.17\textwidth}p{0.26\textwidth}p{0.27\textwidth}p{0.18\textwidth}@{}}
    \toprule
    \textbf{Bước (thư mục)} & \textbf{Script -- câu hỏi} & \textbf{Kết quả chính} & \textbf{Dùng ở} \\
    \midrule
    \texttt{01\_overview} & \texttt{01\_scan\_metadata.py} -- dữ liệu gồm những gì? & \texttt{video\_metadata.csv} & 3.1.1 \\
      & \texttt{02\_dataset\_structure.py} -- chia tập, định dạng, tệp nhãn & \texttt{structure\_summary.txt} \newline \texttt{split\_sizes.png} & 3.1.2 \\
    \midrule
    \texttt{02\_data\_quality} & \texttt{01\_integrity\_missing.py} -- giá trị thiếu, tệp lỗi & \texttt{integrity\_summary.txt} & 3.1.3 \\
      & \texttt{02\_duplicates.py} -- trùng lặp, xung đột nhãn, trùng giữa các tập & \texttt{file\_hashes.csv} \newline \texttt{dups\_summary.txt} \newline \texttt{dedup\_summary.txt} \newline \texttt{class\_counts\_train\_unique\_vs\_all.csv} & 3.1.3, 3.1.5 \\
      & \texttt{03\_frame\_scan.py} -- quét nội dung từng khung & \texttt{video\_quality\_flags.csv} & 3.1.3, 3.1.6, 3.1.7 \\
      & \texttt{04\_frame\_quality\_summary.py} -- tổng hợp khung đơn sắc & \texttt{solid\_frame\_summary.txt} \newline \texttt{class\_solid\_frame\_ratio.csv} & 3.1.3, 3.1.6 \\
    \midrule
    \texttt{03\_univariate} & \texttt{01\_numeric\_attributes.py} -- phân phối, ngoại lai & \texttt{desc\_by\_split.csv} \newline \texttt{univariate\_summary.txt} \newline 3 histogram & 3.1.3, 3.1.4 \\
    \midrule
    \texttt{04\_target\_analysis} & \texttt{01\_class\_distribution.py} -- phân phối nhãn & \texttt{class\_summary.txt} \newline \texttt{class\_counts\_train.csv} \newline 3 hình & 3.1.5 \\
    \midrule
    \texttt{05\_bivariate\_multivariate} & \texttt{01\_correlation\_and\_label\_association.py} & \texttt{multivariate\_summary.txt} \newline 3 hình & 3.1.6 \\
    \midrule
    \texttt{06\_split\_comparison} & \texttt{01\_split\_shift.py} -- khác biệt phân phối giữa các tập & \texttt{split\_shift\_summary.txt} \newline \texttt{ecdf\_by\_split.png} & 3.1.7 \\
    \midrule
    \texttt{07\_visual\_inspection} & \texttt{01\_sample\_frames.py} -- khung hình mẫu & 2 lưới khung hình & 3.1.2 \\
      & \texttt{02\_flagged\_videos.py} -- xem lại video bị đánh dấu & \texttt{flagged\_videos\_info.txt} \newline 2 hình & 3.1.3 \\
    \bottomrule
  \end{tabular}
\end{table}

% =====================================================================
\subsubsection{Mục tiêu và tổng quan dữ liệu}
\label{subsec:eda_overview}

\paragraph{Bài toán và thước đo} Bài toán là phân loại video ngôn ngữ ký hiệu (\textit{Sign Language Video Classification}) thành $N=100$ lớp từ vựng tiếng Việt. Đầu vào là một video ghi cử chỉ của một người, được tiền xử lý thành chuỗi $T=16$ khung hình RGB; đầu ra là nhãn dự đoán. Thước đo chính thức là Macro F1-Score. Biến mục tiêu là nhãn lớp (biến phân loại, 100 giá trị), còn các thuộc tính mô tả video (số khung hình, FPS, dung lượng\ldots) chỉ đóng vai trò thông tin phụ trợ cho phân tích.

\paragraph{Nguồn dữ liệu} Bộ dữ liệu là đề chính thức của cuộc thi \textit{Olympic Trí tuệ Nhân tạo Sinh viên 2025 -- Vòng Sơ loại} do Khoa Trí tuệ Nhân tạo, Học viện Công nghệ Bưu chính Viễn thông (PTIT) tổ chức. Các lớp thuộc nhiều nhóm chủ đề: y tế -- dịch bệnh (\textit{Bệnh viện}, \textit{Cách ly}, \textit{Khẩu trang}, \textit{Virus}, \textit{Sốt}, \textit{Ho}\ldots), cảm xúc (\textit{Lo lắng}, \textit{Xúc động}, \textit{Thương}, \textit{Ghét}\ldots), hành động thường ngày (\textit{Ăn}, \textit{Uống}, \textit{Chạy}, \textit{Đi}, \textit{Mua}\ldots), thời gian -- vị trí (\textit{Ban ngày}, \textit{Trưa}, \textit{Hôm nay}, \textit{Phía sau}, \textit{Rẽ trái}\ldots) và phương tiện (\textit{Ô tô}, \textit{Xe máy}, \textit{Xe đạp}).

\paragraph{Từ điển dữ liệu} Vì dữ liệu là video nên không có bảng thuộc tính sẵn. Nhóm duyệt cây thư mục, mở từng video bằng \texttt{cv2.VideoCapture} và gom thông tin vào một bảng metadata duy nhất (\texttt{01\_overview/video\_metadata.csv}, 8.364 dòng). Mã băm MD5 được lưu ở \texttt{02\_data\_quality/file\_hashes.csv}, còn các thuộc tính suy ra từ nội dung khung hình (khung đơn sắc, chuyển động) được lưu ở \texttt{02\_data\_quality/video\_quality\_flags.csv}. Ba bảng có chung khoá (tập, lớp, tên tệp). Bảng~\ref{tab:eda_dict} liệt kê ý nghĩa từng thuộc tính; mọi thống kê phía sau chỉ là phép ghép và \texttt{groupby} trên ba bảng này.

\begin{table}[H]
  \centering
  \caption{Từ điển các thuộc tính dùng trong phân tích.}
  \label{tab:eda_dict}
  \small
  \begin{tabular}{@{}p{0.22\textwidth}p{0.17\textwidth}p{0.5\textwidth}@{}}
    \toprule
    \textbf{Thuộc tính} & \textbf{Kiểu} & \textbf{Ý nghĩa / cách tính} \\
    \midrule
    \texttt{split}        & phân loại & Tập chứa video: \texttt{train}, \texttt{public\_test}, \texttt{private\_test} \\
    \texttt{cls}          & phân loại & Nhãn lớp, lấy từ tên thư mục con của \texttt{train} (không có ở hai tập test) \\
    \texttt{file}         & chuỗi     & Tên tệp \texttt{.mp4}; hậu tố \texttt{\_N} biểu thị bản sao \\
    \texttt{frames}       & số nguyên & Số khung hình (\texttt{CAP\_\allowbreak PROP\_\allowbreak FRAME\_\allowbreak COUNT}) \\
    \texttt{fps}          & số thực   & Tốc độ khung hình ghi trong container \\
    \texttt{w}, \texttt{h}& số nguyên & Chiều rộng, chiều cao khung hình (pixel) \\
    \texttt{dur}          & số thực   & Thời lượng (giây) $=\texttt{frames}/\texttt{fps}$ \\
    \texttt{size\_kb}     & số thực   & Dung lượng tệp (KB) \\
    \texttt{md5}          & chuỗi     & Mã băm nội dung tệp, dùng phát hiện video trùng (lưu trong \texttt{file\_hashes.csv}) \\
    \texttt{solid\_ratio} & số thực   & Tỷ lệ khung đơn sắc trong video (khung có độ lệch chuẩn điểm ảnh $<3$) \\
    \texttt{motion}       & số thực   & Độ chuyển động trung bình $\overline{|f_t-f_{t-1}|}$ giữa các khung liên tiếp (ảnh xám) \\
    \bottomrule
  \end{tabular}
\end{table}

% =====================================================================
\subsubsection{Cấu trúc dữ liệu}
\label{subsec:eda_structure}

\paragraph{Quy mô và cách chia} Sau khi giải nén \texttt{dataset.zip} (khoảng 864~MB), dữ liệu gồm ba tập video và một tệp ánh xạ nhãn (Bảng~\ref{tab:eda_splits}, Hình~\ref{fig:eda_split_sizes}). Tệp \texttt{label\_mapping.pkl} là một \texttt{dict[str, int]} chứa đúng 100 lớp, đánh chỉ số 0--99 theo thứ tự bảng chữ cái tiếng Việt (\textit{An ủi}~$=0$, \textit{Áp dụng}~$=1$, \textit{Ăn}~$=2$, \ldots, \textit{Đồng ý}~$=99$). Khi đối chiếu với tên thư mục lớp, nhóm phát hiện hai bên dùng \emph{hai dạng chuẩn hoá Unicode khác nhau}: khoá trong \texttt{label\_mapping.pkl} ở dạng NFD (dấu tách rời khỏi chữ cái), còn tên thư mục ở dạng NFC (dấu dựng sẵn). Vì vậy 91/100 tên lớp (mọi lớp có dấu) không khớp nếu so nguyên văn, dù hiển thị giống hệt nhau; sau khi chuẩn hoá cả hai về NFC thì 100 lớp khớp hoàn toàn và chỉ số liên tục 0--99 không trùng.

\begin{table}[H]
  \centering
  \caption{Cấu trúc và quy mô các tập dữ liệu.}
  \label{tab:eda_splits}
  \begin{tabular}{@{}lp{0.42\textwidth}rr@{}}
    \toprule
    \textbf{Thành phần} & \textbf{Nội dung} & \textbf{Số video} & \textbf{Dung lượng} \\
    \midrule
    \texttt{train/}         & 100 thư mục con, mỗi thư mục là một lớp (có nhãn) & 3.875 & 323~MB \\
    \texttt{public\_test/}  & Video không nhãn, dùng cho bảng xếp hạng công khai & 1.630 & 181~MB \\
    \texttt{private\_test/} & Video không nhãn, dùng cho xếp hạng cuối cùng      & 2.859 & 325~MB \\
    \texttt{label\_mapping.pkl} & \texttt{dict[str, int]}: tên lớp (Unicode NFD) $\rightarrow$ chỉ số 0--99 & -- & -- \\
    \midrule
    \textbf{Tổng} & & \textbf{8.364} & $\approx$\textbf{829~MB} \\
    \bottomrule
  \end{tabular}
\end{table}

\begin{figure}[H]
  \centering
  \includegraphics[width=0.6\textwidth]{split_sizes.png}
  \caption{Số lượng video trong ba tập \texttt{train}, \texttt{public\_test} và \texttt{private\_test}.}
  \label{fig:eda_split_sizes}
\end{figure}

\paragraph{Định dạng video} Toàn bộ 8.364 tệp đều là \texttt{.mp4} và đọc được; không có video nào có 0 khung hình. Độ phân giải đồng nhất tuyệt đối: 100\% video ở cả ba tập đều là $224\times224$ (Bảng~\ref{tab:eda_res}). Số khung hình gần như cố định ở mức 30 (chi tiết ở Mục~\ref{subsec:eda_uni}).

\begin{table}[H]
  \centering
  \caption{Thống kê độ phân giải video.}
  \label{tab:eda_res}
  \begin{tabular}{@{}lrrr@{}}
    \toprule
    \textbf{Tập} & \textbf{Độ phân giải} & \textbf{Số video} & \textbf{Tỷ lệ} \\
    \midrule
    train         & $224\times224$ & 3.875 & 100\% \\
    public\_test  & $224\times224$ & 1.630 & 100\% \\
    private\_test & $224\times224$ & 2.859 & 100\% \\
    \bottomrule
  \end{tabular}
\end{table}

\paragraph{Quan sát trực quan} Để kiểm chứng bằng mắt, nhóm trích 6--8 khung hình phân bố đều trong video (\texttt{cap.set(cv2.CAP\_\allowbreak PROP\_\allowbreak POS\_\allowbreak FRAMES, i)} với $i$ lấy từ \texttt{np.linspace}) và ghép thành lưới ảnh cho một số lớp (Hình~\ref{fig:eda_samples}). Các video được quay chính diện, người ký ngồi hoặc đứng giữa khung hình, thấy rõ phần thân trên, hai tay và khuôn mặt; nền chủ yếu là tường trơn với ánh sáng trong nhà. Sự đa dạng nằm ở người ký (giới tính, trang phục, kính), màu nền (be, xanh nhạt, trắng) và vị trí tay lúc bắt đầu. Ký hiệu thường chỉ chiếm một phần thời lượng: các khung đầu và cuối thường là tư thế nghỉ. Ở một số video xuất hiện các khung \emph{đơn sắc hoàn toàn} (hàng \textit{Phỏng vấn}, \textit{Uống}); hiện tượng này được định lượng ở Mục~\ref{subsec:eda_quality}.

\begin{figure}[H]
  \centering
  \includegraphics[width=0.85\textwidth]{sample_frames_6classes.png}
  \caption{Khung hình mẫu (lấy đều theo thời gian) của sáu lớp. Hàng \textit{Phỏng vấn} và \textit{Uống} có các khung đơn sắc ở cuối video.}
  \label{fig:eda_samples}
\end{figure}

\noindent\textit{Hệ quả:} khi gán nhãn số cho video phải chuẩn hoá tên lớp bằng \texttt{unicodedata.normalize(\textquotesingle NFC\textquotesingle, ...)} trước khi tra \texttt{label\_mapping.pkl}, nếu không phần lớn video sẽ không tìm được nhãn. Chỉ có tập \texttt{train} mang nhãn nên tập kiểm định (\textit{validation}) phải tách ra từ chính tập này; việc tách cần tính đến trùng lặp (Mục~\ref{subsec:eda_quality}). Thông tin phân biệt nằm ở hình dạng và chuyển động của bàn tay, khẩu hình và tư thế, nên các phép tăng cường không được làm đổi ngữ nghĩa không gian; đặc biệt \emph{không} dùng lật ngang (\textit{Horizontal Flip}) vì có thể đổi nghĩa tay trái/phải.

% =====================================================================
\subsubsection{Chất lượng dữ liệu}
\label{subsec:eda_quality}

\paragraph{Script và cách thực hiện}
\begin{itemize}
  \item \textbf{Script:} bốn script trong \texttt{02\_data\_quality/}, chạy theo thứ tự: \edafile{02\_data\_quality}{01\_integrity\_missing.py}, \edafile{02\_data\_quality}{02\_duplicates.py}, \edafile{02\_data\_quality}{03\_frame\_scan.py}, \edafile{02\_data\_quality}{04\_frame\_quality\_summary.py}. Kiểm tra bằng mắt dùng \edafile{07\_visual\_inspection}{02\_flagged\_videos.py}.
  \item \textbf{Cách khai thác dữ liệu:} 
  
  (i) \emph{Giá trị thiếu và lỗi đọc:} \texttt{01\_integrity\_missing.py} đếm ô trống (\texttt{isna}) theo cột và theo tập, đếm video có 0 khung, FPS hoặc kích thước bằng 0, rồi liệt kê các tệp có giá trị phi lý theo ngưỡng tuyệt đối (dưới 10~KB, dưới 16 khung, FPS dưới 5). 
  
  (ii) \emph{Trùng lặp:} \texttt{02\_duplicates.py} đọc toàn bộ byte của từng tệp, tính \texttt{hashlib.md5} một lần và lưu vào \texttt{file\_hashes.csv}; hai tệp cùng mã băm là cùng nội dung. Gom nhóm theo mã băm để đếm bản sao, tìm nhóm có nhiều hơn một nhãn (xung đột nhãn), lấy giao các tập mã băm giữa train, public test, private test, rồi giữ một tệp cho mỗi mã băm và đếm lại số video mỗi lớp. 
  
  (iii) \emph{Khung đơn sắc:} \texttt{03\_frame\_scan.py} giải mã \emph{từng} khung bằng \texttt{cap.read()}, tính độ lệch chuẩn điểm ảnh của từng kênh màu và coi khung là đơn sắc khi giá trị lớn nhất nhỏ hơn 3; đồng thời đếm số khung đơn sắc liên tiếp ở đầu video và tính độ chuyển động là trung bình $|f_t-f_{t-1}|$ trên ảnh xám. \texttt{04\_frame\_quality\_summary.py} đối chiếu số khung đọc được với metadata, bỏ các bản sao \texttt{\_N}, phân loại vị trí khung đơn sắc (đầu, cuối, giữa, toàn bộ) và mô phỏng việc lấy $T=16$ khung bằng \texttt{np.linspace} để đếm số khung đơn sắc bị lấy trúng. 
  
  (iv) \emph{Kiểm tra bằng mắt:} \texttt{02\_flagged\_videos.py} \emph{tự chọn} video cần xem từ kết quả hai bước trên, không dùng danh sách chọn tay. Tiêu chí chọn gồm: 6 video có tỷ lệ khung đơn sắc cao nhất trong số các video có từ 50\% khung đơn sắc trở lên; mọi video thuộc nhóm MD5 nằm ở hơn một lớp; một ví dụ cho mỗi dạng xuất hiện khung đơn sắc. Ngoại lai theo quy tắc Tukey được tính ở bước đơn biến (Mục~\ref{subsec:eda_uni}).
  \item \textbf{Đầu ra:} \texttt{integrity\_summary.txt}, \texttt{file\_hashes.csv}, \texttt{dups\_summary.txt}, \texttt{dedup\_summary.txt}, \texttt{class\_counts\_train\_unique\_vs\_all.csv}, \texttt{video\_quality\_flags.csv}, \texttt{solid\_frame\_summary.txt}, \texttt{class\_solid\_frame\_ratio.csv}; ở bước kiểm tra bằng mắt là \texttt{flagged\_videos\_info.txt}, \texttt{flagged\_videos\_check.png}, \texttt{solid\_frame\_patterns.png}.
\end{itemize}

Nhóm kiểm tra chất lượng dữ liệu theo các khía cạnh: giá trị thiếu, bản ghi trùng lặp, khung hình đơn sắc/video rỗng, giá trị ngoại lai, nhiễu nhãn và trùng lặp giữa các tập. Bảng~\ref{tab:eda_quality} tóm tắt các vấn đề phát hiện được.

\begin{table}[H]
  \centering
  \caption{Tóm tắt các vấn đề chất lượng dữ liệu.}
  \label{tab:eda_quality}
  \small
  \begin{tabular}{@{}p{0.26\textwidth}p{0.36\textwidth}p{0.30\textwidth}@{}}
    \toprule
    \textbf{Vấn đề} & \textbf{Số lượng} & \textbf{Ví dụ} \\
    \midrule
    Bản sao hậu tố \texttt{\_N} trong train & 1.069 tệp giống hệt bản gốc; còn 2.776 video duy nhất & \texttt{127824.mp4}, \texttt{127824\_1.mp4}, \texttt{127824\_2.mp4} \\
    Video có khung đơn sắc (train) & 758/3.875 tệp; 565/2.776 video duy nhất (20,4\%), ở 78/100 lớp & \texttt{Thất lạc/708519}, \texttt{Bộ y tế/523389} \\
    Video rỗng hoàn toàn & 2 video gốc (4 tệp kể cả bản sao) & \texttt{Biết/214576}, \texttt{Khai báo/885339} \\
    Nhiễu nhãn (cùng nội dung, khác lớp) & 2 nhóm nội dung, 5 tệp & \textit{Chậm lại} $\leftrightarrow$ \textit{Chào} \\
    Trùng train $\rightarrow$ public test & 157 video public giống hệt video train & -- \\
    Trùng public $\leftrightarrow$ private & 896 video & -- \\
    \bottomrule
  \end{tabular}
\end{table}

\paragraph{Giá trị thiếu và lỗi đọc} Bảng metadata không có giá trị thiếu ở bất kỳ thuộc tính nào (nhãn \texttt{cls} chỉ trống ở hai tập test vì các tập này vốn không có nhãn). Mọi video đều mở và giải mã được, và ở cả 8.364 video, số khung giải mã thực tế khớp với số khung ghi trong metadata. Như vậy không có vấn đề về tệp hỏng hay nhãn bị thiếu ở tập train.

\paragraph{Trùng lặp trong tập train} Nhóm tính mã băm MD5 nội dung từng tệp (\texttt{hashlib.md5}) rồi gom nhóm theo mã băm. Có 1.069/3.875 tệp mang hậu tố \texttt{\_1}, \texttt{\_2}, \texttt{\_3}\ldots{} So sánh từng cặp cho thấy cả 1.069 tệp này đều có bản gốc tồn tại, cùng lớp và \emph{giống hệt bản gốc từng byte} (cùng số khung, FPS, dung lượng): đây là bản sao, không phải biến thể tăng cường. Có 176 video gốc có 1 bản sao, 169 video có 2 và 185 video có 3. Tính cả các bản sao khác, tập train có 1.099 tệp dư thừa và chỉ còn \textbf{2.776 video duy nhất}. Các bản sao tập trung ở những lớp vốn đã nhiều mẫu (\textit{Chậm lại} và \textit{Trường học} cùng thêm 25 tệp, \textit{Rau} và \textit{Nhìn} thêm 23, \textit{Khóc} thêm 22\ldots), nên sau khi khử trùng lặp số video mỗi lớp còn từ 5 đến 51 (trung bình 27,8; tỷ lệ chênh $\approx$10,2 lần). Bảng trước/sau khử trùng lặp của từng lớp lưu ở \texttt{class\_counts\_train\_unique\_vs\_all.csv}. Hệ quả là nguy cơ rò rỉ dữ liệu (\textit{Data Leakage}): nếu chia train/validation ngẫu nhiên theo từng tệp, các bản sao của cùng một video rơi vào cả hai phần, làm điểm validation bị thổi phồng.

\paragraph{Khung hình đơn sắc và video rỗng} Nhóm đọc \emph{toàn bộ} khung hình của 8.364 video và tính độ lệch chuẩn giá trị điểm ảnh trên từng kênh màu; khung có độ lệch chuẩn lớn nhất giữa ba kênh nhỏ hơn 3 được coi là \emph{khung đơn sắc}. Kết quả tổng hợp ở Bảng~\ref{tab:eda_solid}.

\begin{table}[H]
  \centering
  \caption{Số video chứa khung hình đơn sắc theo từng tập.}
  \label{tab:eda_solid}
  \begin{tabular}{@{}lrrrr@{}}
    \toprule
    \textbf{Tập} & \textbf{Tổng video} & \textbf{Có $\geq$1 khung đơn sắc} & \textbf{Đơn sắc toàn bộ} & \textbf{Tỷ lệ} \\
    \midrule
    train (tất cả tệp)        & 3.875 & 758 & 4 & 19,6\% \\
    train (video duy nhất)    & 2.776 & 565 & 2 & 20,4\% \\
    public\_test              & 1.630 & 1   & 0 & 0,06\% \\
    private\_test             & 2.859 & 1   & 0 & 0,03\% \\
    \bottomrule
  \end{tabular}
\end{table}

Đây là phát hiện đáng chú ý nhất của EDA. Khoảng một phần năm số video duy nhất trong train (565/2.776, thuộc 78/100 lớp) chứa khung đơn sắc, và số khung này không nhỏ: trung bình 12,2 khung trên khoảng 30 khung mỗi video (trung vị 12, nhiều nhất 23). Các khung có màu thuần (đỏ, xanh lục, xanh lam, xanh lục lam, tím, vàng, trắng, đen, xám) và xuất hiện theo các dạng: chỉ ở đầu video (231 video), chỉ ở cuối (161), cả đầu và cuối (116), xen ở giữa (55) hoặc toàn bộ video (2) (Hình~\ref{fig:eda_solid_examples}). Lớp có tỷ lệ video bị ảnh hưởng cao nhất là \textit{Có thể} (6/15 video duy nhất), \textit{Xin lỗi} (5/14) và \textit{Thất lạc} (7/20) (chi tiết ở \texttt{class\_solid\_frame\_ratio.csv}).

\begin{figure}[H]
  \centering
  \includegraphics[width=\textwidth]{solid_frame_patterns.png}
  \caption{Bốn dạng khung đơn sắc trong tập train (khung chẵn $f_0$--$f_{28}$), mỗi dạng một video được chọn tự động: chỉ ở đầu (\textit{An ủi}/133326), chỉ ở cuối (\textit{Ban ngày}/363626), xen ở giữa (\textit{Ban ngày}/725380), cả đầu và cuối (\textit{Ban ngày}/855731).}
  \label{fig:eda_solid_examples}
\end{figure}

Hiện tượng này ảnh hưởng trực tiếp đến cách lấy mẫu: khi lấy $T=16$ khung đều nhau bằng \texttt{linspace} trên 565 video nói trên, trung bình có 6,4/16 khung được chọn là khung đơn sắc, và ở 224 video có từ 8/16 khung trở lên (ít nhất một nửa chuỗi đầu vào) là khung đơn sắc. Ngược lại, mỗi tập test chỉ có 1 video chứa đúng 1 khung đơn sắc. Như vậy khung đơn sắc là đặc điểm \emph{riêng của tập train}: nếu không xử lý, mô hình vừa mất thông tin ký hiệu, vừa có thể học các tín hiệu giả không tồn tại ở dữ liệu kiểm tra.

Sáu hàng đầu của Hình~\ref{fig:eda_suspicious} là các video có tỷ lệ khung đơn sắc cao nhất. Hai video được xác nhận bằng mắt là \textbf{rỗng hoàn toàn}: \texttt{Biết/214576.mp4} (2,5~KB, 16 khung, toàn màu xanh lục) và \texttt{Khai báo/885339.mp4} (2,7~KB, 23 khung, toàn màu xanh lam, kèm 2 bản sao \texttt{\_1}, \texttt{\_2}); cả hai không có người ký. Bốn video tiếp theo có 20--21/30 khung đơn sắc: \texttt{Nôn ói/915809}, \texttt{Thất lạc/708519}, \texttt{Chân/984642} và \texttt{Khai báo/362529}. Người ký chỉ xuất hiện trong khoảng 10 khung, ví dụ \texttt{Khai báo/362529} có 20 khung đỏ đầu tiên ($f_0$--$f_{19}$).

\paragraph{Giá trị ngoại lai} Nhóm dùng quy tắc Tukey (ngoài khoảng $[Q_1-1{,}5\,\mathrm{IQR},\,Q_3+1{,}5\,\mathrm{IQR}]$ tính trên tập train) cho thời lượng, dung lượng và FPS. Số lượng ngoại lai khá nhiều nhưng phần lớn chỉ lệch nhẹ (Mục~\ref{subsec:eda_uni}). Đáng chú ý là các ngoại lai cực đoan: thời lượng trên 5~s có 10 video train, 1 video public test và 3 video private test (lớn nhất 35,3~s ở train, 28,1~s ở private test); FPS dưới 5 có 9 video train và 2 video private test; dung lượng dưới 10~KB có 4 tệp train, đều là các video rỗng nói trên. Các ngoại lai này không nhất thiết là lỗi mà phản ánh việc video được lấy mẫu lại về $\approx$30 khung từ các clip gốc có độ dài rất khác nhau.

\paragraph{Nhiễu nhãn} Khi gom nhóm theo MD5, có hai nhóm nội dung giống hệt nhau nhưng nằm ở \emph{hai lớp khác nhau}: \texttt{Chậm lại/284430.mp4} trùng với \texttt{Chào/392683.mp4} và \texttt{Chào/468203.mp4}; \texttt{Chậm lại/254739.mp4} trùng với \texttt{Chào/792094.mp4}. Nhóm đã xem trực tiếp (Hình~\ref{fig:eda_suspicious}, năm hàng cuối) và xác nhận chúng có cùng người ký, cùng chuỗi động tác. Một video không thể đồng thời mang hai nghĩa nên ít nhất một trong hai nhãn là sai; tuy nhiên việc xác định nhãn đúng cần người am hiểu ngôn ngữ ký hiệu Việt Nam, nên báo cáo chỉ ghi nhận đây là \emph{xung đột nhãn}.

\begin{figure}[H]
  \centering
  \includegraphics[width=0.9\textwidth,height=0.72\textheight,keepaspectratio]{flagged_videos_check.png}
  \caption{Kiểm tra bằng mắt các video do bước kiểm tra chất lượng tự đánh dấu. Sáu hàng đầu là các video có tỷ lệ khung đơn sắc cao nhất, gồm hai video rỗng hoàn toàn (\textit{Biết}/214576, \textit{Khai báo}/885339). Năm hàng cuối là hai nhóm video trùng nội dung nhưng khác nhãn (\textit{Chậm lại} $\leftrightarrow$ \textit{Chào}).}
  \label{fig:eda_suspicious}
\end{figure}

\paragraph{Trùng lặp giữa các tập} Public test có 8 bản trùng nội bộ (1.622 video duy nhất) và \textbf{157 video giống hệt video trong train}; private test có 2.782 video duy nhất, không trùng với train, nhưng có 896 video trùng với public test. Ngoài ra, 2.806 tên tệp trong private test trùng với tên tệp trong train nhưng nội dung hoàn toàn khác (MD5 khác, chênh lệch điểm ảnh lớn), nên đây chỉ là trùng cách đặt mã số; nguyên nhân chưa rõ và tên tệp \emph{không} được dùng để suy luận nhãn.

\noindent\textit{Hệ quả:} dữ liệu train chưa sạch; trước khi huấn luyện cần khử trùng lặp, loại video rỗng, xử lý khung đơn sắc và rà soát các nhãn xung đột. Điểm trên public test có thể bị ảnh hưởng một phần bởi 157 video đã có trong train.

% =====================================================================
\subsubsection{Phân tích đơn biến}
\label{subsec:eda_uni}

\paragraph{Script và cách thực hiện}
\begin{itemize}
  \item \textbf{Script:} \edafile{03\_univariate}{01\_numeric\_attributes.py}.
  \item \textbf{Cách khai thác dữ liệu:} với mỗi thuộc tính số (\texttt{frames}, \texttt{fps}, \texttt{dur}, \texttt{size\_kb}), script tính \texttt{groupby(\textquotesingle split\textquotesingle).describe()} (trung bình, độ lệch chuẩn, min, các tứ phân vị, max), hệ số lệch (\textit{skewness}) và các giá trị phổ biến nhất. Ngoại lai được xác định theo quy tắc Tukey với ngưỡng tính trên train rồi áp lên cả ba tập. Script vẽ histogram của từng thuộc tính cho ba tập trên cùng trục; histogram FPS và dung lượng dùng mật độ (\texttt{density=True}) để so sánh được các tập có kích thước khác nhau.
  \item \textbf{Đầu ra:} \texttt{desc\_by\_split.csv} (Bảng~\ref{tab:eda_desc}), \texttt{univariate\_summary.txt}, \texttt{frames\_hist\_by\_split.png}, \texttt{fps\_hist\_by\_split.png}, \texttt{filesize\_hist\_by\_split.png}.
\end{itemize}

Mục này khảo sát phân phối của từng thuộc tính số: số khung hình, FPS, thời lượng và dung lượng, tính trên từng tập (Bảng~\ref{tab:eda_desc}). Độ phân giải đã được trình bày ở Bảng~\ref{tab:eda_res}.

\begin{table}[H]
  \centering
  \caption{Thống kê mô tả các thuộc tính số theo từng tập.}
  \label{tab:eda_desc}
  \small
  \begin{tabular}{@{}llrrrrrrr@{}}
    \toprule
    \textbf{Thuộc tính} & \textbf{Tập} & \textbf{Mean} & \textbf{Std} & \textbf{Min} & \textbf{Q1} & \textbf{Median} & \textbf{Q3} & \textbf{Max} \\
    \midrule
    \multirow{3}{*}{Số khung}
      & train   & 29,88 & 1,26 & 16 & 30 & 30 & 30 & 30 \\
      & public  & 29,96 & 0,66 & 16 & 30 & 30 & 30 & 30 \\
      & private & 29,98 & 0,44 & 21 & 30 & 30 & 30 & 30 \\
    \midrule
    \multirow{3}{*}{FPS}
      & train   & 12,94 & 3,11 & 0,85 & 10,98 & 12,34 & 14,06 & 30,02 \\
      & public  & 12,97 & 3,00 & 5,93 & 10,98 & 12,50 & 13,85 & 30,00 \\
      & private & 13,32 & 3,69 & 1,07 & 10,84 & 12,50 & 14,06 & 30,02 \\
    \midrule
    \multirow{3}{*}{Thời lượng (s)}
      & train   & 2,48 & 1,36 & 0,70 & 2,13 & 2,43 & 2,73 & 35,29 \\
      & public  & 2,42 & 0,51 & 0,70 & 2,17 & 2,40 & 2,73 & 5,06 \\
      & private & 2,42 & 0,82 & 0,70 & 2,13 & 2,40 & 2,77 & 28,14 \\
    \midrule
    \multirow{3}{*}{Dung lượng (KB)}
      & train   & 85,5  & 38,5 & 2,5  & 60,2 & 80,6 & 103,0 & 263,8 \\
      & public  & 113,6 & 77,5 & 26,2 & 76,1 & 91,0 & 125,8 & 483,2 \\
      & private & 116,3 & 77,7 & 42,2 & 77,7 & 94,6 & 126,0 & 488,0 \\
    \bottomrule
  \end{tabular}
\end{table}

\paragraph{Số khung hình} Hầu hết video có đúng 30 khung: 3.832/3.875 video train (98,9\%), 1.623/1.630 public test và 2.850/2.859 private test; ở train có 26 video chỉ có 16 khung. Không video nào ngắn hơn 16 khung (Hình~\ref{fig:eda_frames_hist}), nên không cần đệm khung thực sự khi lấy $T=16$.

\begin{figure}[H]
  \centering
  \includegraphics[width=\textwidth]{frames_hist_by_split.png}
  \caption{Histogram số khung hình mỗi video theo tập (trục tung thang log). Đường đứt nét đánh dấu $T=16$.}
  \label{fig:eda_frames_hist}
\end{figure}

\paragraph{FPS và thời lượng} Dù số khung gần như cố định, FPS biến thiên mạnh: train có 961 giá trị FPS khác nhau, từ 0,85 đến 30,02 (trung vị 12,3), phân phối lệch phải với đuôi dài (hệ số lệch 1,01 ở train; Hình~\ref{fig:eda_fps_hist}). Vì thời lượng $=\texttt{frames}/\texttt{fps}$ nên thời lượng thực của cùng 30 khung dao động lớn: trung vị khoảng 2,4~s nhưng có video tới 35,3~s ở train và 28,1~s ở private test, trong khi 99\% video đều ngắn hơn khoảng 4~s. Điều này cho thấy video đã được lấy mẫu lại về khoảng 30 khung bất kể độ dài clip gốc.

\begin{figure}[H]
  \centering
  \includegraphics[width=0.85\textwidth]{fps_hist_by_split.png}
  \caption{Phân phối FPS theo từng tập.}
  \label{fig:eda_fps_hist}
\end{figure}

\paragraph{Dung lượng tệp} Dung lượng train có trung vị 80,6~KB (trung bình 85,5; từ 2,5 đến 263,8~KB). Hai tập test có phân phối lệch phải hơn hẳn: trung bình 113,6--116,3~KB, lớn nhất tới 483--488~KB (Hình~\ref{fig:eda_size_hist}); khác biệt này được phân tích ở Mục~\ref{subsec:eda_shift}.

\begin{figure}[H]
  \centering
  \includegraphics[width=0.85\textwidth]{filesize_hist_by_split.png}
  \caption{Phân phối dung lượng tệp (KB) theo từng tập.}
  \label{fig:eda_size_hist}
\end{figure}

\noindent\textit{Hệ quả:} (i) không cần đổi kích thước khung hình, tuy kiến trúc CRNN với \texttt{Adaptive\allowbreak AvgPool2d} vẫn đảm bảo tính tổng quát nếu gặp video khác độ phân giải; (ii) vì mọi video có ít nhất 16 khung, lấy mẫu $T=16$ khung đều nhau bằng \texttt{linspace} (\textit{Frame Downsampling}) là đủ, lặp khung (\textit{repeat padding}) chỉ giữ làm phương án dự phòng; (iii) phân phối FPS và dung lượng lệch phải với nhiều ngoại lai nên nên dùng trung vị thay vì trung bình khi mô tả.

% =====================================================================
\subsubsection{Phân tích biến mục tiêu}
\label{subsec:eda_class}

\paragraph{Script và cách thực hiện}
\begin{itemize}
  \item \textbf{Script:} \edafile{04\_target\_analysis}{01\_class\_distribution.py}.
  \item \textbf{Cách khai thác dữ liệu:} lọc các dòng của tập train trong \texttt{video\_metadata.csv}, đếm số video mỗi nhãn bằng \texttt{value\_counts()}, sau đó tính trung bình, trung vị, độ lệch chuẩn, tỷ lệ lớp lớn nhất/nhỏ nhất và số lớp có không quá 25 hoặc từ 50 video trở lên. Số video mỗi lớp sau khử trùng lặp lấy từ \texttt{class\_counts\_train\_unique\_vs\_all.csv} của bước \texttt{02\_data\_quality}.
  \item \textbf{Đầu ra:} \texttt{class\_summary.txt}, \texttt{class\_counts\_train.csv}, \texttt{class\_distribution\_all.png}, \texttt{class\_top\_bottom10.png}, \texttt{samples\_per\_class\_hist.png}.
\end{itemize}

Nhóm dùng \texttt{value\_counts()} trên cột nhãn của tập \texttt{train} để đếm số video mỗi lớp, rồi vẽ biểu đồ cho cả 100 lớp (Hình~\ref{fig:eda_class_dist}), 10 lớp nhiều nhất/ít nhất (Hình~\ref{fig:eda_top_bottom}) và histogram số video mỗi lớp (Hình~\ref{fig:eda_spc_hist}).

\begin{figure}[H]
  \centering
  \includegraphics[width=\textwidth]{class_distribution_all.png}
  \caption{Phân phối số video theo lớp trong tập huấn luyện (100 lớp, sắp xếp giảm dần).}
  \label{fig:eda_class_dist}
\end{figure}

Tập huấn luyện gồm 3.875 video nhưng phân bố không đồng đều: trung bình 38,75 video/lớp, trung vị 45,5, độ lệch chuẩn 22,78. Lớp nhiều nhất (\textit{Trường học}) có 74 video, trong khi mười lớp chỉ có 6--7 video; tỷ lệ chênh giữa lớp lớn nhất và nhỏ nhất xấp xỉ $12{,}3$ lần (Bảng~\ref{tab:eda_top_bottom}). Có 36/100 lớp có không quá 25 video và 42/100 lớp có từ 50 video trở lên, tức phân phối có dạng hai đỉnh: nhóm lớp ``giàu'' (khoảng 45--74 video) và nhóm lớp ``nghèo'' (khoảng 6--11 video). Nhiều lớp ít mẫu thuộc nhóm từ vựng y tế -- dịch bệnh.

\begin{table}[H]
  \centering
  \caption{Mười lớp nhiều mẫu nhất và mười lớp ít mẫu nhất trong tập huấn luyện.}
  \label{tab:eda_top_bottom}
  \begin{tabular}{@{}lr@{\hspace{2.5em}}lr@{}}
    \toprule
    \multicolumn{2}{c}{\textbf{Nhiều nhất}} & \multicolumn{2}{c}{\textbf{Ít nhất}} \\
    \cmidrule(r){1-2}\cmidrule(l){3-4}
    \textbf{Lớp} & \textbf{Số video} & \textbf{Lớp} & \textbf{Số video} \\
    \midrule
    Trường học & 74 & Nhầm      & 7 \\
    Rau        & 72 & Virus     & 7 \\
    Nhà        & 71 & Bệnh nhân & 6 \\
    Áp dụng    & 71 & Cách ly   & 6 \\
    Chậm lại   & 69 & Cơ thể    & 6 \\
    Cám dỗ     & 68 & Ho        & 6 \\
    Nghe       & 68 & Phục hồi  & 6 \\
    Nói xấu    & 68 & Sốt       & 6 \\
    Biết       & 67 & Xe máy    & 6 \\
    Chạy       & 67 & Ủng hộ    & 6 \\
    \bottomrule
  \end{tabular}
\end{table}

\begin{figure}[H]
  \centering
  \begin{minipage}{0.56\textwidth}
    \centering
    \includegraphics[width=\textwidth]{class_top_bottom10.png}
    \caption{Mười lớp nhiều nhất (xanh) và ít nhất (đỏ).}
    \label{fig:eda_top_bottom}
  \end{minipage}\hfill
  \begin{minipage}{0.42\textwidth}
    \centering
    \includegraphics[width=\textwidth]{samples_per_class_hist.png}
    \caption{Histogram số video mỗi lớp.}
    \label{fig:eda_spc_hist}
  \end{minipage}
\end{figure}

Các con số trên tính trên tệp gốc, \emph{chưa} khử trùng lặp.

\noindent\textit{Hệ quả:} dữ liệu bị mất cân bằng lớp (\textit{Class Imbalance}) rõ rệt. Nếu đánh giá bằng Accuracy, mô hình có thể thiên về lớp phổ biến mà vẫn đạt điểm cao; vì vậy dùng Macro F1 (cũng là thước đo chính thức) và cần biện pháp xử lý mất cân bằng ở bước tiền xử lý/huấn luyện.

% =====================================================================
\subsubsection{Phân tích hai biến và đa biến}
\label{subsec:eda_multi}

\paragraph{Script và cách thực hiện}
\begin{itemize}
  \item \textbf{Script:} \edafile{05\_bivariate\_multivariate}{01\_correlation\_and\_label\_association.py}.
  \item \textbf{Cách khai thác dữ liệu:} ghép \texttt{video\_metadata.csv} với \texttt{video\_quality\_flags.csv} theo khoá (tập, lớp, tên tệp); bỏ các bản sao \texttt{\_N} để chỉ giữ 2.776 video duy nhất của train. Tính ma trận tương quan Pearson và Spearman (tương quan hạng, ít nhạy với ngoại lai) giữa các thuộc tính số; với từng thuộc tính, chia giá trị theo 100 lớp rồi chạy kiểm định Kruskal--Wallis (\texttt{scipy.stats.kruskal}) và tính $\eta^2 = SS_{\text{giữa lớp}}/SS_{\text{tổng}}$. Trung vị và tứ phân vị FPS của từng lớp được sắp xếp tăng dần để vẽ.
  \item \textbf{Đầu ra:} \texttt{multivariate\_summary.txt}, \texttt{corr\_heatmap.png}, \texttt{fps\_by\_class.png}, \texttt{size\_vs\_solid.png}.
\end{itemize}

Vì dữ liệu video không có đặc trưng dạng bảng, phân tích hai biến và đa biến tập trung vào mối quan hệ giữa các thuộc tính metadata với nhau và với nhãn lớp. Để tránh bản sao làm lệch kết quả, các phép tính dưới đây dùng 2.776 video duy nhất của tập train.

\paragraph{Tương quan giữa các thuộc tính} Hình~\ref{fig:eda_corr} cho thấy hệ số tương quan Pearson giữa số khung, FPS, dung lượng, tỷ lệ khung đơn sắc và độ chuyển động. (Thời lượng không đưa vào vì $\texttt{dur}=\texttt{frames}/\texttt{fps}$ nên phụ thuộc hoàn toàn vào hai biến còn lại.) Hầu hết các cặp có tương quan yếu ($|r|<0{,}35$); hai quan hệ rõ nhất là dung lượng -- tỷ lệ khung đơn sắc ($r\approx-0{,}45$; Spearman $\rho\approx-0{,}51$) và FPS -- chuyển động ($r\approx-0{,}34$; $\rho\approx-0{,}30$). Tương quan Spearman cho cùng chiều và độ lớn tương tự, nên các quan hệ này không phải do ngoại lai tạo ra. Quan hệ đầu là hợp lý: khung đơn sắc nén rất nhỏ nên video chứa nhiều khung đơn sắc có dung lượng thấp (Hình~\ref{fig:eda_size_solid}). Quan hệ sau cho thấy video có FPS thấp (clip gốc dài hơn) có độ chuyển động giữa các khung liên tiếp lớn hơn. Số khung gần như không liên quan tới các biến khác vì hầu như không đổi.

\begin{figure}[H]
  \centering
  \begin{minipage}{0.48\textwidth}
    \centering
    \includegraphics[width=\textwidth]{corr_heatmap.png}
    \caption{Ma trận tương quan Pearson giữa các thuộc tính (train, video duy nhất).}
    \label{fig:eda_corr}
  \end{minipage}\hfill
  \begin{minipage}{0.48\textwidth}
    \centering
    \includegraphics[width=\textwidth]{size_vs_solid.png}
    \caption{Dung lượng tệp theo tỷ lệ khung đơn sắc.}
    \label{fig:eda_size_solid}
  \end{minipage}
\end{figure}

\paragraph{Quan hệ giữa thuộc tính và nhãn lớp} Nhóm so sánh phân phối của từng thuộc tính giữa 100 lớp bằng kiểm định Kruskal--Wallis và hệ số $\eta^2$ (tỷ lệ phương sai giải thích bởi nhãn). Kết quả: FPS có $\eta^2\approx0{,}28$, dung lượng $\eta^2\approx0{,}15$, số khung $\eta^2\approx0{,}08$; cả ba đều có $p<10^{-10}$. Nói cách khác, FPS của video khác nhau rõ rệt giữa các lớp: trung vị FPS của lớp dao động từ khoảng 9,4 (\textit{Ngón tay}) đến khoảng 19,2 (\textit{Ho}), trong khi trung vị toàn tập là 12,4 (Hình~\ref{fig:eda_fps_class}). Điều này gợi ý rằng thời lượng clip gốc (và do đó FPS sau khi lấy mẫu lại) có liên hệ với từ được ký -- có thể vì các từ khác nhau có độ dài tự nhiên khác nhau -- hoặc vì cách thu thập dữ liệu theo từng lớp. Nhóm chưa xác định được nguyên nhân.

\begin{figure}[H]
  \centering
  \includegraphics[width=0.95\textwidth]{fps_by_class.png}
  \caption{Trung vị và khoảng tứ phân vị của FPS theo lớp (train, video duy nhất), các lớp sắp xếp theo trung vị.}
  \label{fig:eda_fps_class}
\end{figure}

\paragraph{Khung đơn sắc theo lớp} Khung đơn sắc không phân bố đều: 78/100 lớp có ít nhất một video bị ảnh hưởng, và có lớp có tới 30--40\% video duy nhất chứa khung đơn sắc (\textit{Có thể}, \textit{Xin lỗi}, \textit{Thất lạc}, \textit{Mua}). Do đó việc loại bỏ khung đơn sắc sẽ tác động không đồng đều lên các lớp.

\noindent\textit{Hệ quả:} (i) các thuộc tính metadata phần lớn độc lập nhau nên không có cặp đặc trưng dư thừa cần loại bỏ; (ii) vì FPS và dung lượng có quan hệ thống kê với nhãn, chúng có thể trở thành ``đường tắt'' (\textit{shortcut}) mà mô hình học được thay vì học ký hiệu, nên không nên đưa FPS hoặc dung lượng làm đặc trưng đầu vào và cần theo dõi nguy cơ này khi đánh giá; (iii) cần kiểm tra lại phân phối nhãn sau khi loại khung đơn sắc.

% =====================================================================
\subsubsection{So sánh phân phối giữa các tập}
\label{subsec:eda_shift}

\paragraph{Script và cách thực hiện}
\begin{itemize}
  \item \textbf{Script:} \edafile{06\_split\_comparison}{01\_split\_shift.py}.
  \item \textbf{Cách khai thác dữ liệu:} với train chỉ giữ video duy nhất, để các bản sao không làm lệch phân phối. Với mỗi thuộc tính (\texttt{frames}, \texttt{fps}, \texttt{dur}, \texttt{size\_kb}, \texttt{motion}), script tính trung vị từng tập và chạy kiểm định Kolmogorov--Smirnov hai mẫu (\texttt{scipy.stats.ks\_2samp}) cho từng cặp tập. Thống kê $D$ là khoảng cách lớn nhất giữa hai hàm phân phối tích luỹ thực nghiệm (ECDF); $D$ càng lớn thì hai phân phối càng khác nhau. Script cũng tính tỷ lệ video có khung đơn sắc ở từng tập và số video trùng nội dung giữa các tập (theo \texttt{file\_hashes.csv}). Ngưỡng ngoại lai trên của dung lượng (167~KB, tính từ train) lấy từ bước đơn biến.
  \item \textbf{Đầu ra:} \texttt{split\_shift\_summary.txt}, \texttt{ecdf\_by\_split.png}.
\end{itemize}

Nhóm so sánh phân phối các thuộc tính giữa ba tập trên cùng trục (Bảng~\ref{tab:eda_desc}, Bảng~\ref{tab:eda_ks}, Hình~\ref{fig:eda_ecdf}) để tìm dấu hiệu khác biệt miền (\textit{Domain Shift}).

\begin{table}[H]
  \centering
  \caption{Kiểm định Kolmogorov--Smirnov hai mẫu giữa các tập (train: video duy nhất). Giá trị là thống kê $D$, trong ngoặc là $p$-value.}
  \label{tab:eda_ks}
  \small
  \begin{tabular}{@{}lccc@{}}
    \toprule
    \textbf{Thuộc tính} & \textbf{train -- public} & \textbf{train -- private} & \textbf{public -- private} \\
    \midrule
    Số khung        & 0,005 (1)                 & 0,007 (1)                 & 0,002 (1) \\
    FPS             & 0,075 ($1{,}7\cdot10^{-5}$) & 0,069 ($2{,}9\cdot10^{-6}$) & 0,094 ($1{,}7\cdot10^{-8}$) \\
    Thời lượng      & 0,075 ($1{,}7\cdot10^{-5}$) & 0,068 ($4{,}5\cdot10^{-6}$) & 0,094 ($1{,}7\cdot10^{-8}$) \\
    Dung lượng      & \textbf{0,230} ($2\cdot10^{-48}$) & \textbf{0,237} ($4\cdot10^{-70}$) & 0,059 (0,0013) \\
    Chuyển động     & 0,059 (0,0015)            & 0,079 ($4{,}8\cdot10^{-8}$) & 0,035 (0,15) \\
    \bottomrule
  \end{tabular}
\end{table}

\begin{figure}[H]
  \centering
  \includegraphics[width=\textwidth]{ecdf_by_split.png}
  \caption{Hàm phân phối tích luỹ thực nghiệm của FPS, dung lượng và thời lượng theo tập (train: video duy nhất). Đường dung lượng của train tách hẳn khỏi hai tập test.}
  \label{fig:eda_ecdf}
\end{figure}

Vì cỡ mẫu lớn nên mọi khác biệt dù nhỏ đều có $p$ rất bé; nhóm đánh giá chủ yếu theo độ lớn của $D$. Số khung hình giống nhau ở cả ba tập ($D<0{,}01$), độ phân giải cũng vậy. FPS, thời lượng và độ chuyển động khác biệt nhỏ ($D<0{,}1$). Riêng dung lượng khác biệt rõ giữa train và hai tập test ($D\approx0{,}23$--$0{,}24$), trong khi hai tập test gần nhau ($D\approx0{,}06$). Cụ thể:

\begin{itemize}
  \item \textbf{Dung lượng:} video test nặng hơn đáng kể (trung bình 113,6--116,3~KB so với 85,5~KB; trung vị 91,0--94,6~KB so với 80,6~KB). Tỷ lệ video vượt 167~KB (ngưỡng ngoại lai của train) là 3,1\% ở train nhưng 6,0\% ở public và 8,2\% ở private. Một phần chênh lệch có thể giải thích bởi khung đơn sắc (nén rất nhỏ) chỉ có ở train: trung vị dung lượng của video train có khung đơn sắc là 49,0~KB, còn của video train không có khung đơn sắc là 87,6~KB, gần mức của test. Tuy nhiên đuôi dài tới 488~KB ở test không xuất hiện ở train (lớn nhất 263,8~KB), cho thấy test có thể chứa các video chi tiết hoặc nhiễu hơn.
  \item \textbf{FPS:} trung vị gần như nhau (12,4 ở train, 12,5 ở hai tập test), nhưng private test có độ lệch chuẩn lớn hơn (3,69 so với 3,11 ở train) và đuôi dài hơn.
  \item \textbf{Khung đơn sắc:} khoảng 20\% video train nhưng gần như không có ở test (Bảng~\ref{tab:eda_solid}).
  \item \textbf{Người ký, nền, ánh sáng:} có thể khác giữa train và test; nhóm chưa định lượng được yếu tố này.
\end{itemize}

\noindent\textit{Hệ quả:} cần các phép tăng cường giúp mô hình bền vững với thay đổi màu sắc, độ sáng và khung cắt (\textit{ColorJitter}, \textit{Random Resized Crop}); việc loại bỏ khung đơn sắc ở train đồng thời giúp thu hẹp khoảng cách giữa phân phối train và test.

% =====================================================================
\subsubsection{Kết luận EDA và khuyến nghị cho bước tiền xử lý}
\label{subsec:eda_conclusion}

Bảng~\ref{tab:eda_reco} tổng hợp các vấn đề phát hiện được, bằng chứng và khuyến nghị cho bước tiền xử lý (Mục~\ref{sec:preprocess}) cũng như huấn luyện.

\begin{table}[H]
  \centering
  \caption{Tổng hợp vấn đề, bằng chứng và khuyến nghị cho bước tiền xử lý.}
  \label{tab:eda_reco}
  \small
  \begin{tabular}{@{}p{0.2\textwidth}p{0.36\textwidth}p{0.38\textwidth}@{}}
    \toprule
    \textbf{Vấn đề} & \textbf{Bằng chứng} & \textbf{Khuyến nghị} \\
    \midrule
    Tên lớp khác dạng Unicode & Khoá trong \texttt{label\_mapping.pkl} ở dạng NFD, thư mục lớp ở dạng NFC; 91/100 lớp không khớp nếu so nguyên văn & Chuẩn hoá NFC cả hai phía trước khi tra nhãn \\
    \addlinespace
    Video trùng lặp (\texttt{\_N}) & 1.069 tệp \texttt{\_N} giống hệt bản gốc; chỉ còn 2.776 video duy nhất & Khử trùng lặp theo MD5; chia train/validation theo ID gốc (bỏ hậu tố \texttt{\_N}) bằng \texttt{StratifiedGroupKFold} để tránh rò rỉ \\
    \addlinespace
    Mất cân bằng lớp & 6--74 video/lớp (sau khử trùng lặp: 5--51), 36/100 lớp có $\leq$25 video & Dùng Macro F1; trọng số lớp hoặc \texttt{WeightedRandomSampler}; oversample kèm tăng cường cho lớp ít mẫu; có thể dùng label smoothing \\
    \addlinespace
    Khung đơn sắc & 565/2.776 video duy nhất (20,4\%) chứa trung bình 12,2 khung đơn sắc; test gần như không có & Loại bỏ khung đơn sắc (độ lệch chuẩn điểm ảnh $<3$) \emph{trước} khi lấy mẫu $T=16$; nếu còn ít hơn 16 khung hợp lệ thì lặp khung \\
    \addlinespace
    Video rỗng & \texttt{Biết/214576}, \texttt{Khai báo/885339} (cùng 2 bản sao) đơn sắc toàn bộ & Loại khỏi tập huấn luyện \\
    \addlinespace
    Xung đột nhãn & 2 nhóm nội dung giống hệt nhưng thuộc \textit{Chậm lại} và \textit{Chào} & Xác minh nhãn đúng; nếu không xác minh được thì loại bỏ \\
    \addlinespace
    Độ dài không đồng nhất & 98,9\% video có 30 khung, tối thiểu 16 khung; thời lượng tới 35,3~s & Lấy mẫu $T=16$ khung đều bằng \texttt{linspace} \\
    \addlinespace
    FPS/dung lượng liên hệ với nhãn & $\eta^2\approx0{,}28$ (FPS), $0{,}15$ (dung lượng) giữa các lớp & Không dùng FPS, dung lượng làm đặc trưng; theo dõi nguy cơ shortcut \\
    \addlinespace
    Khác biệt miền train--test & Test nặng hơn (trung bình 114--116~KB so với 85,5~KB; KS $D\approx0{,}23$); khung đơn sắc chỉ có ở train & ColorJitter, Random Resized Crop, chuẩn hoá ImageNet; \emph{không} lật ngang \\
    \addlinespace
    Trùng train--public test & 157 video public giống hệt video train & Lưu ý khi diễn giải điểm public test \\
    \bottomrule
  \end{tabular}
\end{table}

Tóm lại, bộ dữ liệu có định dạng rất đồng nhất (100\% video $224\times224$, gần như đều 30 khung), nên khâu đọc và lấy mẫu khung đơn giản. Tuy nhiên dữ liệu huấn luyện có ba vấn đề cần xử lý trước khi huấn luyện: mất cân bằng lớp, khoảng 28\% tệp là bản sao gây nguy cơ rò rỉ dữ liệu, và khoảng 20\% video chứa khung đơn sắc vốn không xuất hiện ở dữ liệu kiểm tra. Đây là những yếu tố có thể khiến điểm validation cao hơn đáng kể so với điểm thực tế trên tập kiểm tra nếu không được xử lý.

\paragraph{Hạn chế của phân tích} FPS và thời lượng đọc từ metadata của container nên có thể không phản ánh tốc độ quay thực. Ngưỡng độ lệch chuẩn $<3$ chỉ phát hiện khung đơn sắc hoàn toàn, chưa phát hiện khung mờ hoặc chỉ có nền. Nhóm chưa định lượng sự khác biệt về người ký và bối cảnh giữa train và test, chưa xác định được nguyên nhân FPS liên hệ với nhãn, và chưa xác định được nhãn đúng cho các video xung đột.
]
% Cấu trúc theo khung EDA chuẩn:
%   3.1.1 Mục tiêu và tổng quan dữ liệu
%   3.1.2 Cấu trúc dữ liệu
%   3.1.3 Chất lượng dữ liệu
%   3.1.4 Phân tích đơn biến
%   3.1.5 Phân tích biến mục tiêu
%   3.1.6 Phân tích hai biến và đa biến
%   3.1.7 So sánh phân phối giữa các tập
%   3.1.8 Kết luận và khuyến nghị
% Gói cần: graphicx, booktabs, float, caption, amsmath
% =====================================================================

\subsection{Phân tích dữ liệu (EDA)}
\label{sec:eda}

% \edafile{thư mục}{tệp}: in đường dẫn script, cho phép xuống dòng sau dấu /
\newcommand{\edafile}[2]{\texttt{#1/}\allowbreak\texttt{#2}}
% Cho phép xuống dòng sau dấu gạch dưới trong tên tệp dài (vd. class\_counts\_train\_unique\_vs\_all.csv)
\let\edaOrigUS\_
\renewcommand{\_}{\edaOrigUS\allowbreak}
\sloppy

Mục này trình bày kết quả phân tích khám phá dữ liệu (\textit{Exploratory Data Analysis -- EDA}) trên bộ dữ liệu video ngôn ngữ ký hiệu. EDA chỉ nhằm \emph{quan sát} xu hướng, phân phối và chất lượng dữ liệu. Nhóm sử dụng Python làm công cụ để thực hiện EDA (\texttt{os}, \texttt{pickle}, \texttt{pandas}, \texttt{OpenCV}, \texttt{numpy}, \texttt{hashlib}, \texttt{matplotlib}).

Phần phân tích được tổ chức theo trình tự quy trình EDA: tổng quan và cấu trúc dữ liệu (Mục~\ref{subsec:eda_overview}--\ref{subsec:eda_structure}), chất lượng dữ liệu (Mục~\ref{subsec:eda_quality}), phân tích đơn biến (Mục~\ref{subsec:eda_uni}), phân tích biến mục tiêu (Mục~\ref{subsec:eda_class}), phân tích hai biến và đa biến (Mục~\ref{subsec:eda_multi}), so sánh giữa các tập (Mục~\ref{subsec:eda_shift}) và cuối cùng là kết luận kèm khuyến nghị (Mục~\ref{subsec:eda_conclusion}).

\begin{table}[H]
  \centering
  \caption{Các script EDA, thứ tự chạy và mục sử dụng kết quả.}
  \label{tab:eda_scripts}
  \small
  \begin{tabular}{@{}p{0.17\textwidth}p{0.26\textwidth}p{0.27\textwidth}p{0.18\textwidth}@{}}
    \toprule
    \textbf{Bước (thư mục)} & \textbf{Script -- câu hỏi} & \textbf{Kết quả chính} & \textbf{Dùng ở} \\
    \midrule
    \texttt{01\_overview} & \texttt{01\_scan\_metadata.py} -- dữ liệu gồm những gì? & \texttt{video\_metadata.csv} & 3.1.1 \\
      & \texttt{02\_dataset\_structure.py} -- chia tập, định dạng, tệp nhãn & \texttt{structure\_summary.txt} \newline \texttt{split\_sizes.png} & 3.1.2 \\
    \midrule
    \texttt{02\_data\_quality} & \texttt{01\_integrity\_missing.py} -- giá trị thiếu, tệp lỗi & \texttt{integrity\_summary.txt} & 3.1.3 \\
      & \texttt{02\_duplicates.py} -- trùng lặp, xung đột nhãn, trùng giữa các tập & \texttt{file\_hashes.csv} \newline \texttt{dups\_summary.txt} \newline \texttt{dedup\_summary.txt} \newline \texttt{class\_counts\_train\_unique\_vs\_all.csv} & 3.1.3, 3.1.5 \\
      & \texttt{03\_frame\_scan.py} -- quét nội dung từng khung & \texttt{video\_quality\_flags.csv} & 3.1.3, 3.1.6, 3.1.7 \\
      & \texttt{04\_frame\_quality\_summary.py} -- tổng hợp khung đơn sắc & \texttt{solid\_frame\_summary.txt} \newline \texttt{class\_solid\_frame\_ratio.csv} & 3.1.3, 3.1.6 \\
    \midrule
    \texttt{03\_univariate} & \texttt{01\_numeric\_attributes.py} -- phân phối, ngoại lai & \texttt{desc\_by\_split.csv} \newline \texttt{univariate\_summary.txt} \newline 3 histogram & 3.1.3, 3.1.4 \\
    \midrule
    \texttt{04\_target\_analysis} & \texttt{01\_class\_distribution.py} -- phân phối nhãn & \texttt{class\_summary.txt} \newline \texttt{class\_counts\_train.csv} \newline 3 hình & 3.1.5 \\
    \midrule
    \texttt{05\_bivariate\_multivariate} & \texttt{01\_correlation\_and\_label\_association.py} & \texttt{multivariate\_summary.txt} \newline 3 hình & 3.1.6 \\
    \midrule
    \texttt{06\_split\_comparison} & \texttt{01\_split\_shift.py} -- khác biệt phân phối giữa các tập & \texttt{split\_shift\_summary.txt} \newline \texttt{ecdf\_by\_split.png} & 3.1.7 \\
    \midrule
    \texttt{07\_visual\_inspection} & \texttt{01\_sample\_frames.py} -- khung hình mẫu & 2 lưới khung hình & 3.1.2 \\
      & \texttt{02\_flagged\_videos.py} -- xem lại video bị đánh dấu & \texttt{flagged\_videos\_info.txt} \newline 2 hình & 3.1.3 \\
    \bottomrule
  \end{tabular}
\end{table}

% =====================================================================
\subsubsection{Mục tiêu và tổng quan dữ liệu}
\label{subsec:eda_overview}

\paragraph{Bài toán và thước đo} Bài toán là phân loại video ngôn ngữ ký hiệu (\textit{Sign Language Video Classification}) thành $N=100$ lớp từ vựng tiếng Việt. Đầu vào là một video ghi cử chỉ của một người, được tiền xử lý thành chuỗi $T=16$ khung hình RGB; đầu ra là nhãn dự đoán. Thước đo chính thức là Macro F1-Score. Biến mục tiêu là nhãn lớp (biến phân loại, 100 giá trị), còn các thuộc tính mô tả video (số khung hình, FPS, dung lượng\ldots) chỉ đóng vai trò thông tin phụ trợ cho phân tích.

\paragraph{Nguồn dữ liệu} Bộ dữ liệu là đề chính thức của cuộc thi \textit{Olympic Trí tuệ Nhân tạo Sinh viên 2025 -- Vòng Sơ loại} do Khoa Trí tuệ Nhân tạo, Học viện Công nghệ Bưu chính Viễn thông (PTIT) tổ chức. Các lớp thuộc nhiều nhóm chủ đề: y tế -- dịch bệnh (\textit{Bệnh viện}, \textit{Cách ly}, \textit{Khẩu trang}, \textit{Virus}, \textit{Sốt}, \textit{Ho}\ldots), cảm xúc (\textit{Lo lắng}, \textit{Xúc động}, \textit{Thương}, \textit{Ghét}\ldots), hành động thường ngày (\textit{Ăn}, \textit{Uống}, \textit{Chạy}, \textit{Đi}, \textit{Mua}\ldots), thời gian -- vị trí (\textit{Ban ngày}, \textit{Trưa}, \textit{Hôm nay}, \textit{Phía sau}, \textit{Rẽ trái}\ldots) và phương tiện (\textit{Ô tô}, \textit{Xe máy}, \textit{Xe đạp}).

\paragraph{Từ điển dữ liệu} Vì dữ liệu là video nên không có bảng thuộc tính sẵn. Nhóm duyệt cây thư mục, mở từng video bằng \texttt{cv2.VideoCapture} và gom thông tin vào một bảng metadata duy nhất (\texttt{01\_overview/video\_metadata.csv}, 8.364 dòng). Mã băm MD5 được lưu ở \texttt{02\_data\_quality/file\_hashes.csv}, còn các thuộc tính suy ra từ nội dung khung hình (khung đơn sắc, chuyển động) được lưu ở \texttt{02\_data\_quality/video\_quality\_flags.csv}. Ba bảng có chung khoá (tập, lớp, tên tệp). Bảng~\ref{tab:eda_dict} liệt kê ý nghĩa từng thuộc tính; mọi thống kê phía sau chỉ là phép ghép và \texttt{groupby} trên ba bảng này.

\begin{table}[H]
  \centering
  \caption{Từ điển các thuộc tính dùng trong phân tích.}
  \label{tab:eda_dict}
  \small
  \begin{tabular}{@{}p{0.22\textwidth}p{0.17\textwidth}p{0.5\textwidth}@{}}
    \toprule
    \textbf{Thuộc tính} & \textbf{Kiểu} & \textbf{Ý nghĩa / cách tính} \\
    \midrule
    \texttt{split}        & phân loại & Tập chứa video: \texttt{train}, \texttt{public\_test}, \texttt{private\_test} \\
    \texttt{cls}          & phân loại & Nhãn lớp, lấy từ tên thư mục con của \texttt{train} (không có ở hai tập test) \\
    \texttt{file}         & chuỗi     & Tên tệp \texttt{.mp4}; hậu tố \texttt{\_N} biểu thị bản sao \\
    \texttt{frames}       & số nguyên & Số khung hình (\texttt{CAP\_\allowbreak PROP\_\allowbreak FRAME\_\allowbreak COUNT}) \\
    \texttt{fps}          & số thực   & Tốc độ khung hình ghi trong container \\
    \texttt{w}, \texttt{h}& số nguyên & Chiều rộng, chiều cao khung hình (pixel) \\
    \texttt{dur}          & số thực   & Thời lượng (giây) $=\texttt{frames}/\texttt{fps}$ \\
    \texttt{size\_kb}     & số thực   & Dung lượng tệp (KB) \\
    \texttt{md5}          & chuỗi     & Mã băm nội dung tệp, dùng phát hiện video trùng (lưu trong \texttt{file\_hashes.csv}) \\
    \texttt{solid\_ratio} & số thực   & Tỷ lệ khung đơn sắc trong video (khung có độ lệch chuẩn điểm ảnh $<3$) \\
    \texttt{motion}       & số thực   & Độ chuyển động trung bình $\overline{|f_t-f_{t-1}|}$ giữa các khung liên tiếp (ảnh xám) \\
    \bottomrule
  \end{tabular}
\end{table}

% =====================================================================
\subsubsection{Cấu trúc dữ liệu}
\label{subsec:eda_structure}

\paragraph{Quy mô và cách chia} Sau khi giải nén \texttt{dataset.zip} (khoảng 864~MB), dữ liệu gồm ba tập video và một tệp ánh xạ nhãn (Bảng~\ref{tab:eda_splits}, Hình~\ref{fig:eda_split_sizes}). Tệp \texttt{label\_mapping.pkl} là một \texttt{dict[str, int]} chứa đúng 100 lớp, đánh chỉ số 0--99 theo thứ tự bảng chữ cái tiếng Việt (\textit{An ủi}~$=0$, \textit{Áp dụng}~$=1$, \textit{Ăn}~$=2$, \ldots, \textit{Đồng ý}~$=99$). Khi đối chiếu với tên thư mục lớp, nhóm phát hiện hai bên dùng \emph{hai dạng chuẩn hoá Unicode khác nhau}: khoá trong \texttt{label\_mapping.pkl} ở dạng NFD (dấu tách rời khỏi chữ cái), còn tên thư mục ở dạng NFC (dấu dựng sẵn). Vì vậy 91/100 tên lớp (mọi lớp có dấu) không khớp nếu so nguyên văn, dù hiển thị giống hệt nhau; sau khi chuẩn hoá cả hai về NFC thì 100 lớp khớp hoàn toàn và chỉ số liên tục 0--99 không trùng.

\begin{table}[H]
  \centering
  \caption{Cấu trúc và quy mô các tập dữ liệu.}
  \label{tab:eda_splits}
  \begin{tabular}{@{}lp{0.42\textwidth}rr@{}}
    \toprule
    \textbf{Thành phần} & \textbf{Nội dung} & \textbf{Số video} & \textbf{Dung lượng} \\
    \midrule
    \texttt{train/}         & 100 thư mục con, mỗi thư mục là một lớp (có nhãn) & 3.875 & 323~MB \\
    \texttt{public\_test/}  & Video không nhãn, dùng cho bảng xếp hạng công khai & 1.630 & 181~MB \\
    \texttt{private\_test/} & Video không nhãn, dùng cho xếp hạng cuối cùng      & 2.859 & 325~MB \\
    \texttt{label\_mapping.pkl} & \texttt{dict[str, int]}: tên lớp (Unicode NFD) $\rightarrow$ chỉ số 0--99 & -- & -- \\
    \midrule
    \textbf{Tổng} & & \textbf{8.364} & $\approx$\textbf{829~MB} \\
    \bottomrule
  \end{tabular}
\end{table}

\begin{figure}[H]
  \centering
  \includegraphics[width=0.6\textwidth]{split_sizes.png}
  \caption{Số lượng video trong ba tập \texttt{train}, \texttt{public\_test} và \texttt{private\_test}.}
  \label{fig:eda_split_sizes}
\end{figure}

\paragraph{Định dạng video} Toàn bộ 8.364 tệp đều là \texttt{.mp4} và đọc được; không có video nào có 0 khung hình. Độ phân giải đồng nhất tuyệt đối: 100\% video ở cả ba tập đều là $224\times224$ (Bảng~\ref{tab:eda_res}). Số khung hình gần như cố định ở mức 30 (chi tiết ở Mục~\ref{subsec:eda_uni}).

\begin{table}[H]
  \centering
  \caption{Thống kê độ phân giải video.}
  \label{tab:eda_res}
  \begin{tabular}{@{}lrrr@{}}
    \toprule
    \textbf{Tập} & \textbf{Độ phân giải} & \textbf{Số video} & \textbf{Tỷ lệ} \\
    \midrule
    train         & $224\times224$ & 3.875 & 100\% \\
    public\_test  & $224\times224$ & 1.630 & 100\% \\
    private\_test & $224\times224$ & 2.859 & 100\% \\
    \bottomrule
  \end{tabular}
\end{table}

\paragraph{Quan sát trực quan} Để kiểm chứng bằng mắt, nhóm trích 6--8 khung hình phân bố đều trong video (\texttt{cap.set(cv2.CAP\_\allowbreak PROP\_\allowbreak POS\_\allowbreak FRAMES, i)} với $i$ lấy từ \texttt{np.linspace}) và ghép thành lưới ảnh cho một số lớp (Hình~\ref{fig:eda_samples}). Các video được quay chính diện, người ký ngồi hoặc đứng giữa khung hình, thấy rõ phần thân trên, hai tay và khuôn mặt; nền chủ yếu là tường trơn với ánh sáng trong nhà. Sự đa dạng nằm ở người ký (giới tính, trang phục, kính), màu nền (be, xanh nhạt, trắng) và vị trí tay lúc bắt đầu. Ký hiệu thường chỉ chiếm một phần thời lượng: các khung đầu và cuối thường là tư thế nghỉ. Ở một số video xuất hiện các khung \emph{đơn sắc hoàn toàn} (hàng \textit{Phỏng vấn}, \textit{Uống}); hiện tượng này được định lượng ở Mục~\ref{subsec:eda_quality}.

\begin{figure}[H]
  \centering
  \includegraphics[width=0.85\textwidth]{sample_frames_6classes.png}
  \caption{Khung hình mẫu (lấy đều theo thời gian) của sáu lớp. Hàng \textit{Phỏng vấn} và \textit{Uống} có các khung đơn sắc ở cuối video.}
  \label{fig:eda_samples}
\end{figure}

\noindent\textit{Hệ quả:} khi gán nhãn số cho video phải chuẩn hoá tên lớp bằng \texttt{unicodedata.normalize(\textquotesingle NFC\textquotesingle, ...)} trước khi tra \texttt{label\_mapping.pkl}, nếu không phần lớn video sẽ không tìm được nhãn. Chỉ có tập \texttt{train} mang nhãn nên tập kiểm định (\textit{validation}) phải tách ra từ chính tập này; việc tách cần tính đến trùng lặp (Mục~\ref{subsec:eda_quality}). Thông tin phân biệt nằm ở hình dạng và chuyển động của bàn tay, khẩu hình và tư thế, nên các phép tăng cường không được làm đổi ngữ nghĩa không gian; đặc biệt \emph{không} dùng lật ngang (\textit{Horizontal Flip}) vì có thể đổi nghĩa tay trái/phải.

% =====================================================================
\subsubsection{Chất lượng dữ liệu}
\label{subsec:eda_quality}

\paragraph{Script và cách thực hiện}
\begin{itemize}
  \item \textbf{Script:} bốn script trong \texttt{02\_data\_quality/}, chạy theo thứ tự: \edafile{02\_data\_quality}{01\_integrity\_missing.py}, \edafile{02\_data\_quality}{02\_duplicates.py}, \edafile{02\_data\_quality}{03\_frame\_scan.py}, \edafile{02\_data\_quality}{04\_frame\_quality\_summary.py}. Kiểm tra bằng mắt dùng \edafile{07\_visual\_inspection}{02\_flagged\_videos.py}.
  \item \textbf{Cách khai thác dữ liệu:} 
  
  (i) \emph{Giá trị thiếu và lỗi đọc:} \texttt{01\_integrity\_missing.py} đếm ô trống (\texttt{isna}) theo cột và theo tập, đếm video có 0 khung, FPS hoặc kích thước bằng 0, rồi liệt kê các tệp có giá trị phi lý theo ngưỡng tuyệt đối (dưới 10~KB, dưới 16 khung, FPS dưới 5). 
  
  (ii) \emph{Trùng lặp:} \texttt{02\_duplicates.py} đọc toàn bộ byte của từng tệp, tính \texttt{hashlib.md5} một lần và lưu vào \texttt{file\_hashes.csv}; hai tệp cùng mã băm là cùng nội dung. Gom nhóm theo mã băm để đếm bản sao, tìm nhóm có nhiều hơn một nhãn (xung đột nhãn), lấy giao các tập mã băm giữa train, public test, private test, rồi giữ một tệp cho mỗi mã băm và đếm lại số video mỗi lớp. 
  
  (iii) \emph{Khung đơn sắc:} \texttt{03\_frame\_scan.py} giải mã \emph{từng} khung bằng \texttt{cap.read()}, tính độ lệch chuẩn điểm ảnh của từng kênh màu và coi khung là đơn sắc khi giá trị lớn nhất nhỏ hơn 3; đồng thời đếm số khung đơn sắc liên tiếp ở đầu video và tính độ chuyển động là trung bình $|f_t-f_{t-1}|$ trên ảnh xám. \texttt{04\_frame\_quality\_summary.py} đối chiếu số khung đọc được với metadata, bỏ các bản sao \texttt{\_N}, phân loại vị trí khung đơn sắc (đầu, cuối, giữa, toàn bộ) và mô phỏng việc lấy $T=16$ khung bằng \texttt{np.linspace} để đếm số khung đơn sắc bị lấy trúng. 
  
  (iv) \emph{Kiểm tra bằng mắt:} \texttt{02\_flagged\_videos.py} \emph{tự chọn} video cần xem từ kết quả hai bước trên, không dùng danh sách chọn tay. Tiêu chí chọn gồm: 6 video có tỷ lệ khung đơn sắc cao nhất trong số các video có từ 50\% khung đơn sắc trở lên; mọi video thuộc nhóm MD5 nằm ở hơn một lớp; một ví dụ cho mỗi dạng xuất hiện khung đơn sắc. Ngoại lai theo quy tắc Tukey được tính ở bước đơn biến (Mục~\ref{subsec:eda_uni}).
  \item \textbf{Đầu ra:} \texttt{integrity\_summary.txt}, \texttt{file\_hashes.csv}, \texttt{dups\_summary.txt}, \texttt{dedup\_summary.txt}, \texttt{class\_counts\_train\_unique\_vs\_all.csv}, \texttt{video\_quality\_flags.csv}, \texttt{solid\_frame\_summary.txt}, \texttt{class\_solid\_frame\_ratio.csv}; ở bước kiểm tra bằng mắt là \texttt{flagged\_videos\_info.txt}, \texttt{flagged\_videos\_check.png}, \texttt{solid\_frame\_patterns.png}.
\end{itemize}

Nhóm kiểm tra chất lượng dữ liệu theo các khía cạnh: giá trị thiếu, bản ghi trùng lặp, khung hình đơn sắc/video rỗng, giá trị ngoại lai, nhiễu nhãn và trùng lặp giữa các tập. Bảng~\ref{tab:eda_quality} tóm tắt các vấn đề phát hiện được.

\begin{table}[H]
  \centering
  \caption{Tóm tắt các vấn đề chất lượng dữ liệu.}
  \label{tab:eda_quality}
  \small
  \begin{tabular}{@{}p{0.26\textwidth}p{0.36\textwidth}p{0.30\textwidth}@{}}
    \toprule
    \textbf{Vấn đề} & \textbf{Số lượng} & \textbf{Ví dụ} \\
    \midrule
    Bản sao hậu tố \texttt{\_N} trong train & 1.069 tệp giống hệt bản gốc; còn 2.776 video duy nhất & \texttt{127824.mp4}, \texttt{127824\_1.mp4}, \texttt{127824\_2.mp4} \\
    Video có khung đơn sắc (train) & 758/3.875 tệp; 565/2.776 video duy nhất (20,4\%), ở 78/100 lớp & \texttt{Thất lạc/708519}, \texttt{Bộ y tế/523389} \\
    Video rỗng hoàn toàn & 2 video gốc (4 tệp kể cả bản sao) & \texttt{Biết/214576}, \texttt{Khai báo/885339} \\
    Nhiễu nhãn (cùng nội dung, khác lớp) & 2 nhóm nội dung, 5 tệp & \textit{Chậm lại} $\leftrightarrow$ \textit{Chào} \\
    Trùng train $\rightarrow$ public test & 157 video public giống hệt video train & -- \\
    Trùng public $\leftrightarrow$ private & 896 video & -- \\
    \bottomrule
  \end{tabular}
\end{table}

\paragraph{Giá trị thiếu và lỗi đọc} Bảng metadata không có giá trị thiếu ở bất kỳ thuộc tính nào (nhãn \texttt{cls} chỉ trống ở hai tập test vì các tập này vốn không có nhãn). Mọi video đều mở và giải mã được, và ở cả 8.364 video, số khung giải mã thực tế khớp với số khung ghi trong metadata. Như vậy không có vấn đề về tệp hỏng hay nhãn bị thiếu ở tập train.

\paragraph{Trùng lặp trong tập train} Nhóm tính mã băm MD5 nội dung từng tệp (\texttt{hashlib.md5}) rồi gom nhóm theo mã băm. Có 1.069/3.875 tệp mang hậu tố \texttt{\_1}, \texttt{\_2}, \texttt{\_3}\ldots{} So sánh từng cặp cho thấy cả 1.069 tệp này đều có bản gốc tồn tại, cùng lớp và \emph{giống hệt bản gốc từng byte} (cùng số khung, FPS, dung lượng): đây là bản sao, không phải biến thể tăng cường. Có 176 video gốc có 1 bản sao, 169 video có 2 và 185 video có 3. Tính cả các bản sao khác, tập train có 1.099 tệp dư thừa và chỉ còn \textbf{2.776 video duy nhất}. Các bản sao tập trung ở những lớp vốn đã nhiều mẫu (\textit{Chậm lại} và \textit{Trường học} cùng thêm 25 tệp, \textit{Rau} và \textit{Nhìn} thêm 23, \textit{Khóc} thêm 22\ldots), nên sau khi khử trùng lặp số video mỗi lớp còn từ 5 đến 51 (trung bình 27,8; tỷ lệ chênh $\approx$10,2 lần). Bảng trước/sau khử trùng lặp của từng lớp lưu ở \texttt{class\_counts\_train\_unique\_vs\_all.csv}. Hệ quả là nguy cơ rò rỉ dữ liệu (\textit{Data Leakage}): nếu chia train/validation ngẫu nhiên theo từng tệp, các bản sao của cùng một video rơi vào cả hai phần, làm điểm validation bị thổi phồng.

\paragraph{Khung hình đơn sắc và video rỗng} Nhóm đọc \emph{toàn bộ} khung hình của 8.364 video và tính độ lệch chuẩn giá trị điểm ảnh trên từng kênh màu; khung có độ lệch chuẩn lớn nhất giữa ba kênh nhỏ hơn 3 được coi là \emph{khung đơn sắc}. Kết quả tổng hợp ở Bảng~\ref{tab:eda_solid}.

\begin{table}[H]
  \centering
  \caption{Số video chứa khung hình đơn sắc theo từng tập.}
  \label{tab:eda_solid}
  \begin{tabular}{@{}lrrrr@{}}
    \toprule
    \textbf{Tập} & \textbf{Tổng video} & \textbf{Có $\geq$1 khung đơn sắc} & \textbf{Đơn sắc toàn bộ} & \textbf{Tỷ lệ} \\
    \midrule
    train (tất cả tệp)        & 3.875 & 758 & 4 & 19,6\% \\
    train (video duy nhất)    & 2.776 & 565 & 2 & 20,4\% \\
    public\_test              & 1.630 & 1   & 0 & 0,06\% \\
    private\_test             & 2.859 & 1   & 0 & 0,03\% \\
    \bottomrule
  \end{tabular}
\end{table}

Đây là phát hiện đáng chú ý nhất của EDA. Khoảng một phần năm số video duy nhất trong train (565/2.776, thuộc 78/100 lớp) chứa khung đơn sắc, và số khung này không nhỏ: trung bình 12,2 khung trên khoảng 30 khung mỗi video (trung vị 12, nhiều nhất 23). Các khung có màu thuần (đỏ, xanh lục, xanh lam, xanh lục lam, tím, vàng, trắng, đen, xám) và xuất hiện theo các dạng: chỉ ở đầu video (231 video), chỉ ở cuối (161), cả đầu và cuối (116), xen ở giữa (55) hoặc toàn bộ video (2) (Hình~\ref{fig:eda_solid_examples}). Lớp có tỷ lệ video bị ảnh hưởng cao nhất là \textit{Có thể} (6/15 video duy nhất), \textit{Xin lỗi} (5/14) và \textit{Thất lạc} (7/20) (chi tiết ở \texttt{class\_solid\_frame\_ratio.csv}).

\begin{figure}[H]
  \centering
  \includegraphics[width=\textwidth]{solid_frame_patterns.png}
  \caption{Bốn dạng khung đơn sắc trong tập train (khung chẵn $f_0$--$f_{28}$), mỗi dạng một video được chọn tự động: chỉ ở đầu (\textit{An ủi}/133326), chỉ ở cuối (\textit{Ban ngày}/363626), xen ở giữa (\textit{Ban ngày}/725380), cả đầu và cuối (\textit{Ban ngày}/855731).}
  \label{fig:eda_solid_examples}
\end{figure}

Hiện tượng này ảnh hưởng trực tiếp đến cách lấy mẫu: khi lấy $T=16$ khung đều nhau bằng \texttt{linspace} trên 565 video nói trên, trung bình có 6,4/16 khung được chọn là khung đơn sắc, và ở 224 video có từ 8/16 khung trở lên (ít nhất một nửa chuỗi đầu vào) là khung đơn sắc. Ngược lại, mỗi tập test chỉ có 1 video chứa đúng 1 khung đơn sắc. Như vậy khung đơn sắc là đặc điểm \emph{riêng của tập train}: nếu không xử lý, mô hình vừa mất thông tin ký hiệu, vừa có thể học các tín hiệu giả không tồn tại ở dữ liệu kiểm tra.

Sáu hàng đầu của Hình~\ref{fig:eda_suspicious} là các video có tỷ lệ khung đơn sắc cao nhất. Hai video được xác nhận bằng mắt là \textbf{rỗng hoàn toàn}: \texttt{Biết/214576.mp4} (2,5~KB, 16 khung, toàn màu xanh lục) và \texttt{Khai báo/885339.mp4} (2,7~KB, 23 khung, toàn màu xanh lam, kèm 2 bản sao \texttt{\_1}, \texttt{\_2}); cả hai không có người ký. Bốn video tiếp theo có 20--21/30 khung đơn sắc: \texttt{Nôn ói/915809}, \texttt{Thất lạc/708519}, \texttt{Chân/984642} và \texttt{Khai báo/362529}. Người ký chỉ xuất hiện trong khoảng 10 khung, ví dụ \texttt{Khai báo/362529} có 20 khung đỏ đầu tiên ($f_0$--$f_{19}$).

\paragraph{Giá trị ngoại lai} Nhóm dùng quy tắc Tukey (ngoài khoảng $[Q_1-1{,}5\,\mathrm{IQR},\,Q_3+1{,}5\,\mathrm{IQR}]$ tính trên tập train) cho thời lượng, dung lượng và FPS. Số lượng ngoại lai khá nhiều nhưng phần lớn chỉ lệch nhẹ (Mục~\ref{subsec:eda_uni}). Đáng chú ý là các ngoại lai cực đoan: thời lượng trên 5~s có 10 video train, 1 video public test và 3 video private test (lớn nhất 35,3~s ở train, 28,1~s ở private test); FPS dưới 5 có 9 video train và 2 video private test; dung lượng dưới 10~KB có 4 tệp train, đều là các video rỗng nói trên. Các ngoại lai này không nhất thiết là lỗi mà phản ánh việc video được lấy mẫu lại về $\approx$30 khung từ các clip gốc có độ dài rất khác nhau.

\paragraph{Nhiễu nhãn} Khi gom nhóm theo MD5, có hai nhóm nội dung giống hệt nhau nhưng nằm ở \emph{hai lớp khác nhau}: \texttt{Chậm lại/284430.mp4} trùng với \texttt{Chào/392683.mp4} và \texttt{Chào/468203.mp4}; \texttt{Chậm lại/254739.mp4} trùng với \texttt{Chào/792094.mp4}. Nhóm đã xem trực tiếp (Hình~\ref{fig:eda_suspicious}, năm hàng cuối) và xác nhận chúng có cùng người ký, cùng chuỗi động tác. Một video không thể đồng thời mang hai nghĩa nên ít nhất một trong hai nhãn là sai; tuy nhiên việc xác định nhãn đúng cần người am hiểu ngôn ngữ ký hiệu Việt Nam, nên báo cáo chỉ ghi nhận đây là \emph{xung đột nhãn}.

\begin{figure}[H]
  \centering
  \includegraphics[width=0.9\textwidth,height=0.72\textheight,keepaspectratio]{flagged_videos_check.png}
  \caption{Kiểm tra bằng mắt các video do bước kiểm tra chất lượng tự đánh dấu. Sáu hàng đầu là các video có tỷ lệ khung đơn sắc cao nhất, gồm hai video rỗng hoàn toàn (\textit{Biết}/214576, \textit{Khai báo}/885339). Năm hàng cuối là hai nhóm video trùng nội dung nhưng khác nhãn (\textit{Chậm lại} $\leftrightarrow$ \textit{Chào}).}
  \label{fig:eda_suspicious}
\end{figure}

\paragraph{Trùng lặp giữa các tập} Public test có 8 bản trùng nội bộ (1.622 video duy nhất) và \textbf{157 video giống hệt video trong train}; private test có 2.782 video duy nhất, không trùng với train, nhưng có 896 video trùng với public test. Ngoài ra, 2.806 tên tệp trong private test trùng với tên tệp trong train nhưng nội dung hoàn toàn khác (MD5 khác, chênh lệch điểm ảnh lớn), nên đây chỉ là trùng cách đặt mã số; nguyên nhân chưa rõ và tên tệp \emph{không} được dùng để suy luận nhãn.

\noindent\textit{Hệ quả:} dữ liệu train chưa sạch; trước khi huấn luyện cần khử trùng lặp, loại video rỗng, xử lý khung đơn sắc và rà soát các nhãn xung đột. Điểm trên public test có thể bị ảnh hưởng một phần bởi 157 video đã có trong train.

% =====================================================================
\subsubsection{Phân tích đơn biến}
\label{subsec:eda_uni}

\paragraph{Script và cách thực hiện}
\begin{itemize}
  \item \textbf{Script:} \edafile{03\_univariate}{01\_numeric\_attributes.py}.
  \item \textbf{Cách khai thác dữ liệu:} với mỗi thuộc tính số (\texttt{frames}, \texttt{fps}, \texttt{dur}, \texttt{size\_kb}), script tính \texttt{groupby(\textquotesingle split\textquotesingle).describe()} (trung bình, độ lệch chuẩn, min, các tứ phân vị, max), hệ số lệch (\textit{skewness}) và các giá trị phổ biến nhất. Ngoại lai được xác định theo quy tắc Tukey với ngưỡng tính trên train rồi áp lên cả ba tập. Script vẽ histogram của từng thuộc tính cho ba tập trên cùng trục; histogram FPS và dung lượng dùng mật độ (\texttt{density=True}) để so sánh được các tập có kích thước khác nhau.
  \item \textbf{Đầu ra:} \texttt{desc\_by\_split.csv} (Bảng~\ref{tab:eda_desc}), \texttt{univariate\_summary.txt}, \texttt{frames\_hist\_by\_split.png}, \texttt{fps\_hist\_by\_split.png}, \texttt{filesize\_hist\_by\_split.png}.
\end{itemize}

Mục này khảo sát phân phối của từng thuộc tính số: số khung hình, FPS, thời lượng và dung lượng, tính trên từng tập (Bảng~\ref{tab:eda_desc}). Độ phân giải đã được trình bày ở Bảng~\ref{tab:eda_res}.

\begin{table}[H]
  \centering
  \caption{Thống kê mô tả các thuộc tính số theo từng tập.}
  \label{tab:eda_desc}
  \small
  \begin{tabular}{@{}llrrrrrrr@{}}
    \toprule
    \textbf{Thuộc tính} & \textbf{Tập} & \textbf{Mean} & \textbf{Std} & \textbf{Min} & \textbf{Q1} & \textbf{Median} & \textbf{Q3} & \textbf{Max} \\
    \midrule
    \multirow{3}{*}{Số khung}
      & train   & 29,88 & 1,26 & 16 & 30 & 30 & 30 & 30 \\
      & public  & 29,96 & 0,66 & 16 & 30 & 30 & 30 & 30 \\
      & private & 29,98 & 0,44 & 21 & 30 & 30 & 30 & 30 \\
    \midrule
    \multirow{3}{*}{FPS}
      & train   & 12,94 & 3,11 & 0,85 & 10,98 & 12,34 & 14,06 & 30,02 \\
      & public  & 12,97 & 3,00 & 5,93 & 10,98 & 12,50 & 13,85 & 30,00 \\
      & private & 13,32 & 3,69 & 1,07 & 10,84 & 12,50 & 14,06 & 30,02 \\
    \midrule
    \multirow{3}{*}{Thời lượng (s)}
      & train   & 2,48 & 1,36 & 0,70 & 2,13 & 2,43 & 2,73 & 35,29 \\
      & public  & 2,42 & 0,51 & 0,70 & 2,17 & 2,40 & 2,73 & 5,06 \\
      & private & 2,42 & 0,82 & 0,70 & 2,13 & 2,40 & 2,77 & 28,14 \\
    \midrule
    \multirow{3}{*}{Dung lượng (KB)}
      & train   & 85,5  & 38,5 & 2,5  & 60,2 & 80,6 & 103,0 & 263,8 \\
      & public  & 113,6 & 77,5 & 26,2 & 76,1 & 91,0 & 125,8 & 483,2 \\
      & private & 116,3 & 77,7 & 42,2 & 77,7 & 94,6 & 126,0 & 488,0 \\
    \bottomrule
  \end{tabular}
\end{table}

\paragraph{Số khung hình} Hầu hết video có đúng 30 khung: 3.832/3.875 video train (98,9\%), 1.623/1.630 public test và 2.850/2.859 private test; ở train có 26 video chỉ có 16 khung. Không video nào ngắn hơn 16 khung (Hình~\ref{fig:eda_frames_hist}), nên không cần đệm khung thực sự khi lấy $T=16$.

\begin{figure}[H]
  \centering
  \includegraphics[width=\textwidth]{frames_hist_by_split.png}
  \caption{Histogram số khung hình mỗi video theo tập (trục tung thang log). Đường đứt nét đánh dấu $T=16$.}
  \label{fig:eda_frames_hist}
\end{figure}

\paragraph{FPS và thời lượng} Dù số khung gần như cố định, FPS biến thiên mạnh: train có 961 giá trị FPS khác nhau, từ 0,85 đến 30,02 (trung vị 12,3), phân phối lệch phải với đuôi dài (hệ số lệch 1,01 ở train; Hình~\ref{fig:eda_fps_hist}). Vì thời lượng $=\texttt{frames}/\texttt{fps}$ nên thời lượng thực của cùng 30 khung dao động lớn: trung vị khoảng 2,4~s nhưng có video tới 35,3~s ở train và 28,1~s ở private test, trong khi 99\% video đều ngắn hơn khoảng 4~s. Điều này cho thấy video đã được lấy mẫu lại về khoảng 30 khung bất kể độ dài clip gốc.

\begin{figure}[H]
  \centering
  \includegraphics[width=0.85\textwidth]{fps_hist_by_split.png}
  \caption{Phân phối FPS theo từng tập.}
  \label{fig:eda_fps_hist}
\end{figure}

\paragraph{Dung lượng tệp} Dung lượng train có trung vị 80,6~KB (trung bình 85,5; từ 2,5 đến 263,8~KB). Hai tập test có phân phối lệch phải hơn hẳn: trung bình 113,6--116,3~KB, lớn nhất tới 483--488~KB (Hình~\ref{fig:eda_size_hist}); khác biệt này được phân tích ở Mục~\ref{subsec:eda_shift}.

\begin{figure}[H]
  \centering
  \includegraphics[width=0.85\textwidth]{filesize_hist_by_split.png}
  \caption{Phân phối dung lượng tệp (KB) theo từng tập.}
  \label{fig:eda_size_hist}
\end{figure}

\noindent\textit{Hệ quả:} (i) không cần đổi kích thước khung hình, tuy kiến trúc CRNN với \texttt{Adaptive\allowbreak AvgPool2d} vẫn đảm bảo tính tổng quát nếu gặp video khác độ phân giải; (ii) vì mọi video có ít nhất 16 khung, lấy mẫu $T=16$ khung đều nhau bằng \texttt{linspace} (\textit{Frame Downsampling}) là đủ, lặp khung (\textit{repeat padding}) chỉ giữ làm phương án dự phòng; (iii) phân phối FPS và dung lượng lệch phải với nhiều ngoại lai nên nên dùng trung vị thay vì trung bình khi mô tả.

% =====================================================================
\subsubsection{Phân tích biến mục tiêu}
\label{subsec:eda_class}

\paragraph{Script và cách thực hiện}
\begin{itemize}
  \item \textbf{Script:} \edafile{04\_target\_analysis}{01\_class\_distribution.py}.
  \item \textbf{Cách khai thác dữ liệu:} lọc các dòng của tập train trong \texttt{video\_metadata.csv}, đếm số video mỗi nhãn bằng \texttt{value\_counts()}, sau đó tính trung bình, trung vị, độ lệch chuẩn, tỷ lệ lớp lớn nhất/nhỏ nhất và số lớp có không quá 25 hoặc từ 50 video trở lên. Số video mỗi lớp sau khử trùng lặp lấy từ \texttt{class\_counts\_train\_unique\_vs\_all.csv} của bước \texttt{02\_data\_quality}.
  \item \textbf{Đầu ra:} \texttt{class\_summary.txt}, \texttt{class\_counts\_train.csv}, \texttt{class\_distribution\_all.png}, \texttt{class\_top\_bottom10.png}, \texttt{samples\_per\_class\_hist.png}.
\end{itemize}

Nhóm dùng \texttt{value\_counts()} trên cột nhãn của tập \texttt{train} để đếm số video mỗi lớp, rồi vẽ biểu đồ cho cả 100 lớp (Hình~\ref{fig:eda_class_dist}), 10 lớp nhiều nhất/ít nhất (Hình~\ref{fig:eda_top_bottom}) và histogram số video mỗi lớp (Hình~\ref{fig:eda_spc_hist}).

\begin{figure}[H]
  \centering
  \includegraphics[width=\textwidth]{class_distribution_all.png}
  \caption{Phân phối số video theo lớp trong tập huấn luyện (100 lớp, sắp xếp giảm dần).}
  \label{fig:eda_class_dist}
\end{figure}

Tập huấn luyện gồm 3.875 video nhưng phân bố không đồng đều: trung bình 38,75 video/lớp, trung vị 45,5, độ lệch chuẩn 22,78. Lớp nhiều nhất (\textit{Trường học}) có 74 video, trong khi mười lớp chỉ có 6--7 video; tỷ lệ chênh giữa lớp lớn nhất và nhỏ nhất xấp xỉ $12{,}3$ lần (Bảng~\ref{tab:eda_top_bottom}). Có 36/100 lớp có không quá 25 video và 42/100 lớp có từ 50 video trở lên, tức phân phối có dạng hai đỉnh: nhóm lớp ``giàu'' (khoảng 45--74 video) và nhóm lớp ``nghèo'' (khoảng 6--11 video). Nhiều lớp ít mẫu thuộc nhóm từ vựng y tế -- dịch bệnh.

\begin{table}[H]
  \centering
  \caption{Mười lớp nhiều mẫu nhất và mười lớp ít mẫu nhất trong tập huấn luyện.}
  \label{tab:eda_top_bottom}
  \begin{tabular}{@{}lr@{\hspace{2.5em}}lr@{}}
    \toprule
    \multicolumn{2}{c}{\textbf{Nhiều nhất}} & \multicolumn{2}{c}{\textbf{Ít nhất}} \\
    \cmidrule(r){1-2}\cmidrule(l){3-4}
    \textbf{Lớp} & \textbf{Số video} & \textbf{Lớp} & \textbf{Số video} \\
    \midrule
    Trường học & 74 & Nhầm      & 7 \\
    Rau        & 72 & Virus     & 7 \\
    Nhà        & 71 & Bệnh nhân & 6 \\
    Áp dụng    & 71 & Cách ly   & 6 \\
    Chậm lại   & 69 & Cơ thể    & 6 \\
    Cám dỗ     & 68 & Ho        & 6 \\
    Nghe       & 68 & Phục hồi  & 6 \\
    Nói xấu    & 68 & Sốt       & 6 \\
    Biết       & 67 & Xe máy    & 6 \\
    Chạy       & 67 & Ủng hộ    & 6 \\
    \bottomrule
  \end{tabular}
\end{table}

\begin{figure}[H]
  \centering
  \begin{minipage}{0.56\textwidth}
    \centering
    \includegraphics[width=\textwidth]{class_top_bottom10.png}
    \caption{Mười lớp nhiều nhất (xanh) và ít nhất (đỏ).}
    \label{fig:eda_top_bottom}
  \end{minipage}\hfill
  \begin{minipage}{0.42\textwidth}
    \centering
    \includegraphics[width=\textwidth]{samples_per_class_hist.png}
    \caption{Histogram số video mỗi lớp.}
    \label{fig:eda_spc_hist}
  \end{minipage}
\end{figure}

Các con số trên tính trên tệp gốc, \emph{chưa} khử trùng lặp.

\noindent\textit{Hệ quả:} dữ liệu bị mất cân bằng lớp (\textit{Class Imbalance}) rõ rệt. Nếu đánh giá bằng Accuracy, mô hình có thể thiên về lớp phổ biến mà vẫn đạt điểm cao; vì vậy dùng Macro F1 (cũng là thước đo chính thức) và cần biện pháp xử lý mất cân bằng ở bước tiền xử lý/huấn luyện.

% =====================================================================
\subsubsection{Phân tích hai biến và đa biến}
\label{subsec:eda_multi}

\paragraph{Script và cách thực hiện}
\begin{itemize}
  \item \textbf{Script:} \edafile{05\_bivariate\_multivariate}{01\_correlation\_and\_label\_association.py}.
  \item \textbf{Cách khai thác dữ liệu:} ghép \texttt{video\_metadata.csv} với \texttt{video\_quality\_flags.csv} theo khoá (tập, lớp, tên tệp); bỏ các bản sao \texttt{\_N} để chỉ giữ 2.776 video duy nhất của train. Tính ma trận tương quan Pearson và Spearman (tương quan hạng, ít nhạy với ngoại lai) giữa các thuộc tính số; với từng thuộc tính, chia giá trị theo 100 lớp rồi chạy kiểm định Kruskal--Wallis (\texttt{scipy.stats.kruskal}) và tính $\eta^2 = SS_{\text{giữa lớp}}/SS_{\text{tổng}}$. Trung vị và tứ phân vị FPS của từng lớp được sắp xếp tăng dần để vẽ.
  \item \textbf{Đầu ra:} \texttt{multivariate\_summary.txt}, \texttt{corr\_heatmap.png}, \texttt{fps\_by\_class.png}, \texttt{size\_vs\_solid.png}.
\end{itemize}

Vì dữ liệu video không có đặc trưng dạng bảng, phân tích hai biến và đa biến tập trung vào mối quan hệ giữa các thuộc tính metadata với nhau và với nhãn lớp. Để tránh bản sao làm lệch kết quả, các phép tính dưới đây dùng 2.776 video duy nhất của tập train.

\paragraph{Tương quan giữa các thuộc tính} Hình~\ref{fig:eda_corr} cho thấy hệ số tương quan Pearson giữa số khung, FPS, dung lượng, tỷ lệ khung đơn sắc và độ chuyển động. (Thời lượng không đưa vào vì $\texttt{dur}=\texttt{frames}/\texttt{fps}$ nên phụ thuộc hoàn toàn vào hai biến còn lại.) Hầu hết các cặp có tương quan yếu ($|r|<0{,}35$); hai quan hệ rõ nhất là dung lượng -- tỷ lệ khung đơn sắc ($r\approx-0{,}45$; Spearman $\rho\approx-0{,}51$) và FPS -- chuyển động ($r\approx-0{,}34$; $\rho\approx-0{,}30$). Tương quan Spearman cho cùng chiều và độ lớn tương tự, nên các quan hệ này không phải do ngoại lai tạo ra. Quan hệ đầu là hợp lý: khung đơn sắc nén rất nhỏ nên video chứa nhiều khung đơn sắc có dung lượng thấp (Hình~\ref{fig:eda_size_solid}). Quan hệ sau cho thấy video có FPS thấp (clip gốc dài hơn) có độ chuyển động giữa các khung liên tiếp lớn hơn. Số khung gần như không liên quan tới các biến khác vì hầu như không đổi.

\begin{figure}[H]
  \centering
  \begin{minipage}{0.48\textwidth}
    \centering
    \includegraphics[width=\textwidth]{corr_heatmap.png}
    \caption{Ma trận tương quan Pearson giữa các thuộc tính (train, video duy nhất).}
    \label{fig:eda_corr}
  \end{minipage}\hfill
  \begin{minipage}{0.48\textwidth}
    \centering
    \includegraphics[width=\textwidth]{size_vs_solid.png}
    \caption{Dung lượng tệp theo tỷ lệ khung đơn sắc.}
    \label{fig:eda_size_solid}
  \end{minipage}
\end{figure}

\paragraph{Quan hệ giữa thuộc tính và nhãn lớp} Nhóm so sánh phân phối của từng thuộc tính giữa 100 lớp bằng kiểm định Kruskal--Wallis và hệ số $\eta^2$ (tỷ lệ phương sai giải thích bởi nhãn). Kết quả: FPS có $\eta^2\approx0{,}28$, dung lượng $\eta^2\approx0{,}15$, số khung $\eta^2\approx0{,}08$; cả ba đều có $p<10^{-10}$. Nói cách khác, FPS của video khác nhau rõ rệt giữa các lớp: trung vị FPS của lớp dao động từ khoảng 9,4 (\textit{Ngón tay}) đến khoảng 19,2 (\textit{Ho}), trong khi trung vị toàn tập là 12,4 (Hình~\ref{fig:eda_fps_class}). Điều này gợi ý rằng thời lượng clip gốc (và do đó FPS sau khi lấy mẫu lại) có liên hệ với từ được ký -- có thể vì các từ khác nhau có độ dài tự nhiên khác nhau -- hoặc vì cách thu thập dữ liệu theo từng lớp. Nhóm chưa xác định được nguyên nhân.

\begin{figure}[H]
  \centering
  \includegraphics[width=0.95\textwidth]{fps_by_class.png}
  \caption{Trung vị và khoảng tứ phân vị của FPS theo lớp (train, video duy nhất), các lớp sắp xếp theo trung vị.}
  \label{fig:eda_fps_class}
\end{figure}

\paragraph{Khung đơn sắc theo lớp} Khung đơn sắc không phân bố đều: 78/100 lớp có ít nhất một video bị ảnh hưởng, và có lớp có tới 30--40\% video duy nhất chứa khung đơn sắc (\textit{Có thể}, \textit{Xin lỗi}, \textit{Thất lạc}, \textit{Mua}). Do đó việc loại bỏ khung đơn sắc sẽ tác động không đồng đều lên các lớp.

\noindent\textit{Hệ quả:} (i) các thuộc tính metadata phần lớn độc lập nhau nên không có cặp đặc trưng dư thừa cần loại bỏ; (ii) vì FPS và dung lượng có quan hệ thống kê với nhãn, chúng có thể trở thành ``đường tắt'' (\textit{shortcut}) mà mô hình học được thay vì học ký hiệu, nên không nên đưa FPS hoặc dung lượng làm đặc trưng đầu vào và cần theo dõi nguy cơ này khi đánh giá; (iii) cần kiểm tra lại phân phối nhãn sau khi loại khung đơn sắc.

% =====================================================================
\subsubsection{So sánh phân phối giữa các tập}
\label{subsec:eda_shift}

\paragraph{Script và cách thực hiện}
\begin{itemize}
  \item \textbf{Script:} \edafile{06\_split\_comparison}{01\_split\_shift.py}.
  \item \textbf{Cách khai thác dữ liệu:} với train chỉ giữ video duy nhất, để các bản sao không làm lệch phân phối. Với mỗi thuộc tính (\texttt{frames}, \texttt{fps}, \texttt{dur}, \texttt{size\_kb}, \texttt{motion}), script tính trung vị từng tập và chạy kiểm định Kolmogorov--Smirnov hai mẫu (\texttt{scipy.stats.ks\_2samp}) cho từng cặp tập. Thống kê $D$ là khoảng cách lớn nhất giữa hai hàm phân phối tích luỹ thực nghiệm (ECDF); $D$ càng lớn thì hai phân phối càng khác nhau. Script cũng tính tỷ lệ video có khung đơn sắc ở từng tập và số video trùng nội dung giữa các tập (theo \texttt{file\_hashes.csv}). Ngưỡng ngoại lai trên của dung lượng (167~KB, tính từ train) lấy từ bước đơn biến.
  \item \textbf{Đầu ra:} \texttt{split\_shift\_summary.txt}, \texttt{ecdf\_by\_split.png}.
\end{itemize}

Nhóm so sánh phân phối các thuộc tính giữa ba tập trên cùng trục (Bảng~\ref{tab:eda_desc}, Bảng~\ref{tab:eda_ks}, Hình~\ref{fig:eda_ecdf}) để tìm dấu hiệu khác biệt miền (\textit{Domain Shift}).

\begin{table}[H]
  \centering
  \caption{Kiểm định Kolmogorov--Smirnov hai mẫu giữa các tập (train: video duy nhất). Giá trị là thống kê $D$, trong ngoặc là $p$-value.}
  \label{tab:eda_ks}
  \small
  \begin{tabular}{@{}lccc@{}}
    \toprule
    \textbf{Thuộc tính} & \textbf{train -- public} & \textbf{train -- private} & \textbf{public -- private} \\
    \midrule
    Số khung        & 0,005 (1)                 & 0,007 (1)                 & 0,002 (1) \\
    FPS             & 0,075 ($1{,}7\cdot10^{-5}$) & 0,069 ($2{,}9\cdot10^{-6}$) & 0,094 ($1{,}7\cdot10^{-8}$) \\
    Thời lượng      & 0,075 ($1{,}7\cdot10^{-5}$) & 0,068 ($4{,}5\cdot10^{-6}$) & 0,094 ($1{,}7\cdot10^{-8}$) \\
    Dung lượng      & \textbf{0,230} ($2\cdot10^{-48}$) & \textbf{0,237} ($4\cdot10^{-70}$) & 0,059 (0,0013) \\
    Chuyển động     & 0,059 (0,0015)            & 0,079 ($4{,}8\cdot10^{-8}$) & 0,035 (0,15) \\
    \bottomrule
  \end{tabular}
\end{table}

\begin{figure}[H]
  \centering
  \includegraphics[width=\textwidth]{ecdf_by_split.png}
  \caption{Hàm phân phối tích luỹ thực nghiệm của FPS, dung lượng và thời lượng theo tập (train: video duy nhất). Đường dung lượng của train tách hẳn khỏi hai tập test.}
  \label{fig:eda_ecdf}
\end{figure}

Vì cỡ mẫu lớn nên mọi khác biệt dù nhỏ đều có $p$ rất bé; nhóm đánh giá chủ yếu theo độ lớn của $D$. Số khung hình giống nhau ở cả ba tập ($D<0{,}01$), độ phân giải cũng vậy. FPS, thời lượng và độ chuyển động khác biệt nhỏ ($D<0{,}1$). Riêng dung lượng khác biệt rõ giữa train và hai tập test ($D\approx0{,}23$--$0{,}24$), trong khi hai tập test gần nhau ($D\approx0{,}06$). Cụ thể:

\begin{itemize}
  \item \textbf{Dung lượng:} video test nặng hơn đáng kể (trung bình 113,6--116,3~KB so với 85,5~KB; trung vị 91,0--94,6~KB so với 80,6~KB). Tỷ lệ video vượt 167~KB (ngưỡng ngoại lai của train) là 3,1\% ở train nhưng 6,0\% ở public và 8,2\% ở private. Một phần chênh lệch có thể giải thích bởi khung đơn sắc (nén rất nhỏ) chỉ có ở train: trung vị dung lượng của video train có khung đơn sắc là 49,0~KB, còn của video train không có khung đơn sắc là 87,6~KB, gần mức của test. Tuy nhiên đuôi dài tới 488~KB ở test không xuất hiện ở train (lớn nhất 263,8~KB), cho thấy test có thể chứa các video chi tiết hoặc nhiễu hơn.
  \item \textbf{FPS:} trung vị gần như nhau (12,4 ở train, 12,5 ở hai tập test), nhưng private test có độ lệch chuẩn lớn hơn (3,69 so với 3,11 ở train) và đuôi dài hơn.
  \item \textbf{Khung đơn sắc:} khoảng 20\% video train nhưng gần như không có ở test (Bảng~\ref{tab:eda_solid}).
  \item \textbf{Người ký, nền, ánh sáng:} có thể khác giữa train và test; nhóm chưa định lượng được yếu tố này.
\end{itemize}

\noindent\textit{Hệ quả:} cần các phép tăng cường giúp mô hình bền vững với thay đổi màu sắc, độ sáng và khung cắt (\textit{ColorJitter}, \textit{Random Resized Crop}); việc loại bỏ khung đơn sắc ở train đồng thời giúp thu hẹp khoảng cách giữa phân phối train và test.

% =====================================================================
\subsubsection{Kết luận EDA và khuyến nghị cho bước tiền xử lý}
\label{subsec:eda_conclusion}

Bảng~\ref{tab:eda_reco} tổng hợp các vấn đề phát hiện được, bằng chứng và khuyến nghị cho bước tiền xử lý (Mục~\ref{sec:preprocess}) cũng như huấn luyện.

\begin{table}[H]
  \centering
  \caption{Tổng hợp vấn đề, bằng chứng và khuyến nghị cho bước tiền xử lý.}
  \label{tab:eda_reco}
  \small
  \begin{tabular}{@{}p{0.2\textwidth}p{0.36\textwidth}p{0.38\textwidth}@{}}
    \toprule
    \textbf{Vấn đề} & \textbf{Bằng chứng} & \textbf{Khuyến nghị} \\
    \midrule
    Tên lớp khác dạng Unicode & Khoá trong \texttt{label\_mapping.pkl} ở dạng NFD, thư mục lớp ở dạng NFC; 91/100 lớp không khớp nếu so nguyên văn & Chuẩn hoá NFC cả hai phía trước khi tra nhãn \\
    \addlinespace
    Video trùng lặp (\texttt{\_N}) & 1.069 tệp \texttt{\_N} giống hệt bản gốc; chỉ còn 2.776 video duy nhất & Khử trùng lặp theo MD5; chia train/validation theo ID gốc (bỏ hậu tố \texttt{\_N}) bằng \texttt{StratifiedGroupKFold} để tránh rò rỉ \\
    \addlinespace
    Mất cân bằng lớp & 6--74 video/lớp (sau khử trùng lặp: 5--51), 36/100 lớp có $\leq$25 video & Dùng Macro F1; trọng số lớp hoặc \texttt{WeightedRandomSampler}; oversample kèm tăng cường cho lớp ít mẫu; có thể dùng label smoothing \\
    \addlinespace
    Khung đơn sắc & 565/2.776 video duy nhất (20,4\%) chứa trung bình 12,2 khung đơn sắc; test gần như không có & Loại bỏ khung đơn sắc (độ lệch chuẩn điểm ảnh $<3$) \emph{trước} khi lấy mẫu $T=16$; nếu còn ít hơn 16 khung hợp lệ thì lặp khung \\
    \addlinespace
    Video rỗng & \texttt{Biết/214576}, \texttt{Khai báo/885339} (cùng 2 bản sao) đơn sắc toàn bộ & Loại khỏi tập huấn luyện \\
    \addlinespace
    Xung đột nhãn & 2 nhóm nội dung giống hệt nhưng thuộc \textit{Chậm lại} và \textit{Chào} & Xác minh nhãn đúng; nếu không xác minh được thì loại bỏ \\
    \addlinespace
    Độ dài không đồng nhất & 98,9\% video có 30 khung, tối thiểu 16 khung; thời lượng tới 35,3~s & Lấy mẫu $T=16$ khung đều bằng \texttt{linspace} \\
    \addlinespace
    FPS/dung lượng liên hệ với nhãn & $\eta^2\approx0{,}28$ (FPS), $0{,}15$ (dung lượng) giữa các lớp & Không dùng FPS, dung lượng làm đặc trưng; theo dõi nguy cơ shortcut \\
    \addlinespace
    Khác biệt miền train--test & Test nặng hơn (trung bình 114--116~KB so với 85,5~KB; KS $D\approx0{,}23$); khung đơn sắc chỉ có ở train & ColorJitter, Random Resized Crop, chuẩn hoá ImageNet; \emph{không} lật ngang \\
    \addlinespace
    Trùng train--public test & 157 video public giống hệt video train & Lưu ý khi diễn giải điểm public test \\
    \bottomrule
  \end{tabular}
\end{table}

Tóm lại, bộ dữ liệu có định dạng rất đồng nhất (100\% video $224\times224$, gần như đều 30 khung), nên khâu đọc và lấy mẫu khung đơn giản. Tuy nhiên dữ liệu huấn luyện có ba vấn đề cần xử lý trước khi huấn luyện: mất cân bằng lớp, khoảng 28\% tệp là bản sao gây nguy cơ rò rỉ dữ liệu, và khoảng 20\% video chứa khung đơn sắc vốn không xuất hiện ở dữ liệu kiểm tra. Đây là những yếu tố có thể khiến điểm validation cao hơn đáng kể so với điểm thực tế trên tập kiểm tra nếu không được xử lý.

\paragraph{Hạn chế của phân tích} FPS và thời lượng đọc từ metadata của container nên có thể không phản ánh tốc độ quay thực. Ngưỡng độ lệch chuẩn $<3$ chỉ phát hiện khung đơn sắc hoàn toàn, chưa phát hiện khung mờ hoặc chỉ có nền. Nhóm chưa định lượng sự khác biệt về người ký và bối cảnh giữa train và test, chưa xác định được nguyên nhân FPS liên hệ với nhãn, và chưa xác định được nhãn đúng cho các video xung đột.


% 3.4 – Tiền xử lý dữ liệu
% =====================================================================
% Mục 3.2 – Tiền xử lý dữ liệu  [chèn bằng % =====================================================================
% Mục 3.2 – Tiền xử lý dữ liệu  [chèn bằng % =====================================================================
% Mục 3.2 – Tiền xử lý dữ liệu  [chèn bằng \input{chapters/preprocess}]
% Mã nguồn: processing/data_processing.ipynb, processing/preprocess.py
% Hình: image/processing/*.png (sinh từ processing/output/figures/)
%   3.2.1 Tổng quan quy trình
%   3.2.2 Làm sạch dữ liệu
%   3.2.3 Xử lý khung nhiễu và lấy mẫu theo thời gian
%   3.2.4 Chuẩn hoá
%   3.2.5 Chia tập huấn luyện và kiểm định
%   3.2.6 Xử lý mất cân bằng lớp
%   3.2.7 Tăng cường dữ liệu
%   3.2.8 Dữ liệu kiểm tra
%   3.2.9 Kết quả và đầu ra
% =====================================================================

\subsection{Tiền xử lý dữ liệu}
\label{sec:preprocess}

Mục này trình bày các bước biến dữ liệu thô thành đầu vào sạch và thống nhất cho mọi mô hình ở Chương~\ref{sec:model}. Mỗi bước xuất phát từ một vấn đề đã phát hiện ở bước EDA (Bảng~\ref{tab:eda_reco}). Nhóm dùng Python với \texttt{pandas}, \texttt{numpy}, \texttt{OpenCV} và \texttt{scikit-learn}; mã nguồn gồm hai phần:

\begin{itemize}
  \item \textbf{Notebook làm sạch} \texttt{processing/data\_processing.ipynb}: chạy \emph{một lần} trên toàn bộ tập train để kiểm tra, loại bỏ và đánh dấu dữ liệu lỗi, chia tập kiểm định và tính trọng số lớp. Kết quả được lưu thành một bảng chỉ mục (\textit{manifest}) dùng chung cho cả nhóm.
  \item \textbf{Module dùng chung} \texttt{processing/preprocess.py}: các phép biến đổi áp dụng \emph{mỗi lần} mô hình đọc một video (bỏ khung đơn sắc, lấy mẫu khung, tăng cường, chuẩn hoá). Mô hình SVM, CNN-LSTM, Vision Transformer và ứng dụng demo cùng gọi module này, nên dữ liệu ở huấn luyện, kiểm định, kiểm tra và demo được xử lý giống hệt nhau.
\end{itemize}

Nhóm không lưu sẵn khung hình đã xử lý ra đĩa vì hai lý do: lưu 16 khung $224\times224\times3$ cho 2.772 video tốn khoảng 6,7~GB, và phép tăng cường cần tạo ra phiên bản khác nhau ở mỗi lần đọc. Mọi bước có yếu tố ngẫu nhiên đều cố định hạt giống \texttt{SEED = 42} để kết quả tái lập được.

% =====================================================================
\subsubsection{Tổng quan quy trình}
\label{subsec:pp_overview}

Bảng~\ref{tab:pp_overview} liệt kê mười bước tiền xử lý theo thứ tự thực hiện và xếp mỗi bước vào nhóm tác vụ tiền xử lý dữ liệu tương ứng: làm sạch (\textit{data cleaning}), tích hợp (\textit{data integration}), thu gọn (\textit{data reduction}) và biến đổi dữ liệu (\textit{data transformation})~\cite{han2011}.

\begin{table}[H]
  \centering
  \caption{Các bước tiền xử lý, nhóm tác vụ và nơi thực hiện.}
  \label{tab:pp_overview}
  \small
  \begin{tabular}{@{}cp{0.33\textwidth}p{0.3\textwidth}p{0.2\textwidth}@{}}
    \toprule
    \textbf{Bước} & \textbf{Việc} & \textbf{Nhóm tác vụ} & \textbf{Thực hiện ở} \\
    \midrule
    1  & Kiểm tra toàn vẹn, chuẩn hoá tên lớp & Làm sạch: dữ liệu thiếu, lỗi & Notebook \\
    2  & Khử trùng lặp theo MD5 & Tích hợp: dư thừa dữ liệu & Notebook \\
    3  & Loại video xung đột nhãn & Làm sạch: dữ liệu không nhất quán & Notebook \\
    4  & Loại video rỗng & Làm sạch: dữ liệu thiếu & Notebook \\
    5  & Đánh dấu khung đơn sắc & Làm sạch: dữ liệu nhiễu & Notebook, \texttt{preprocess.py} \\
    6  & Lấy đều $T=16$ khung trên các khung hợp lệ & Thu gọn: lấy mẫu & \texttt{preprocess.py} \\
    7  & Chuẩn hoá z-score theo kênh màu & Biến đổi: chuẩn hoá & \texttt{preprocess.py} \\
    8  & Chia train/validation phân tầng & Đánh giá mô hình & Notebook \\
    9  & Trọng số lớp & Xử lý mất cân bằng lớp & Notebook \\
    10 & Tăng cường dữ liệu (chỉ khi huấn luyện) & Sinh thêm dữ liệu & \texttt{preprocess.py} \\
    \bottomrule
  \end{tabular}
\end{table}

Bước 1--5 thu hẹp tập train từ 3.875 tệp xuống 2.772 video sạch (Bảng~\ref{tab:pp_flow}). Bước 6, 7 và 10 không bỏ video nào mà quyết định video được đưa vào mô hình dưới dạng nào. Bước 8 và 9 chuẩn bị cho huấn luyện và đánh giá.

% =====================================================================
\subsubsection{Làm sạch dữ liệu}
\label{subsec:pp_clean}

\paragraph{Quét toàn bộ video} Notebook mở từng video bằng \texttt{cv2.VideoCapture}, đọc \emph{tất cả} khung hình và ghi lại: mã băm MD5 của tệp, số khung, và danh sách chỉ số các khung không đơn sắc (Mục~\ref{subsec:pp_frames}). Việc quét được chạy song song trên nhiều lõi CPU (\texttt{ProcessPoolExecutor}). Để kiểm tra bước quét, nhóm đối chiếu số khung đơn sắc của từng video với cột \texttt{solid\_frames} trong \texttt{video\_quality\_flags.csv} của bước EDA: cả 3.875/3.875 tệp đều khớp.

\paragraph{Bước 1 -- Kiểm tra toàn vẹn} Cả 3.875 tệp đều đọc được và có ít nhất 16 khung, nên không có tệp nào bị loại ở bước này. Tên lớp ở thư mục và khoá trong \texttt{label\_mapping.pkl} được chuẩn hoá về dạng Unicode NFC trước khi so khớp (Mục~\ref{subsec:eda_structure}); sau chuẩn hoá, 100/100 lớp có nhãn số 0--99.

\paragraph{Bước 2 -- Khử trùng lặp} Hai tệp có cùng mã MD5 là cùng một video. Với mỗi mã MD5, notebook giữ lại một tệp (ưu tiên tệp không có hậu tố \texttt{\_N}) và loại các tệp còn lại. Bước này loại 1.099 tệp, trong đó 1.069 tệp mang hậu tố \texttt{\_N}, và còn \textbf{2.776 video duy nhất}, khớp với kết quả EDA. Từ đây mỗi video trong tập sạch có một mã MD5 riêng; tính chất này giúp chia tập kiểm định không bị rò rỉ (Mục~\ref{subsec:pp_split}).

\paragraph{Bước 3 -- Xung đột nhãn} Hai nội dung video vừa nằm ở lớp \textit{Chào} vừa nằm ở lớp \textit{Chậm lại} (5 tệp: \texttt{Chào/392683}, \texttt{Chào/468203}, \texttt{Chào/792094}, \texttt{Chậm lại/254739}, \texttt{Chậm lại/284430}). Nhóm không xác minh được nhãn nào đúng, và giữ một video với nhãn có thể sai sẽ đưa nhiễu nhãn vào cả huấn luyện lẫn đánh giá, nên loại toàn bộ hai nội dung này. Notebook có sẵn biến \texttt{KEEP\_LABEL} để giữ lại các video này với nhãn đã xác minh nếu sau này có thông tin. Sau bước này còn 2.774 video. Hai lớp bị ảnh hưởng vẫn còn nhiều mẫu (\textit{Chào} 22 video, \textit{Chậm lại} 44 video).

\paragraph{Bước 4 -- Video rỗng} Hai video không có khung hợp lệ nào (\texttt{Biết/214576}: 16/16 khung đơn sắc; \texttt{Khai báo/885339}: 23/23 khung đơn sắc) bị loại vì không chứa thông tin ký hiệu. Tập train sạch còn \textbf{2.772 video} và vẫn đủ 100 lớp.

\begin{table}[H]
  \centering
  \caption{Số video còn lại sau từng bước làm sạch tập train.}
  \label{tab:pp_flow}
  \small
  \begin{tabular}{@{}p{0.3\textwidth}rp{0.45\textwidth}@{}}
    \toprule
    \textbf{Bước} & \textbf{Còn lại} & \textbf{Ghi chú} \\
    \midrule
    Ban đầu                          & 3.875 & Toàn bộ tệp trong \texttt{train/} \\
    1 -- Kiểm tra toàn vẹn           & 3.875 & Không có tệp lỗi hoặc thiếu khung \\
    2 -- Khử trùng lặp MD5           & 2.776 & Loại 1.099 bản sao \\
    3 -- Xung đột nhãn               & 2.774 & Loại 2 nội dung (5 tệp) \textit{Chào} $\leftrightarrow$ \textit{Chậm lại} \\
    4 -- Video rỗng                  & \textbf{2.772} & Loại \texttt{Biết/214576}, \texttt{Khai báo/885339} \\
    5 -- Khung đơn sắc               & 2.772 & 563 video có khung bị bỏ qua, không loại video \\
    \bottomrule
  \end{tabular}
\end{table}

Danh sách 1.103 tệp bị loại cùng lý do được lưu ở \texttt{processing/output/removed.csv}, để có thể kiểm tra lại từng trường hợp. Hình~\ref{fig:pp_class_counts} so sánh số video mỗi lớp trước và sau khi làm sạch: sau làm sạch mỗi lớp có từ 5 đến 51 video (trung bình 27,7; trung vị 32,5).

\begin{figure}[H]
  \centering
  \includegraphics[width=\textwidth]{processing/class_counts_before_after.png}
  \caption{Số video mỗi lớp trước (xám) và sau (xanh) khi làm sạch, sắp xếp theo số video sau làm sạch.}
  \label{fig:pp_class_counts}
\end{figure}

% =====================================================================
\subsubsection{Xử lý khung nhiễu và lấy mẫu theo thời gian}
\label{subsec:pp_frames}

\paragraph{Bước 5 -- Đánh dấu khung đơn sắc} Gọi $f_i$ là khung thứ $i$ của video và $\sigma_c(f_i)$ là độ lệch chuẩn giá trị điểm ảnh trên kênh màu $c\in\{R,G,B\}$. Khung $f_i$ được coi là \emph{đơn sắc} nếu
\begin{equation}
  \max_{c\in\{R,G,B\}} \sigma_c(f_i) < 3,
\end{equation}
tức gần như mọi điểm ảnh có cùng một màu; ngưỡng này giống với ngưỡng ở bước EDA. Các khung còn lại là \emph{khung hợp lệ}. Thay vì xoá khung khỏi tệp video, notebook lưu danh sách chỉ số khung hợp lệ $V=(v_1<v_2<\dots<v_m)$ của từng video vào cột \texttt{valid\_idx} của manifest; khi đọc video, \texttt{preprocess.py} chỉ lấy mẫu trong danh sách này.

Trong 2.772 video sạch, 563 video (20,3\%, thuộc 78 lớp) có khung đơn sắc, trung bình 12,2 khung mỗi video. Sau khi bỏ khung đơn sắc, mỗi video còn từ 9 đến 30 khung hợp lệ (trung vị 30); 222 video còn ít hơn 16 khung hợp lệ (Hình~\ref{fig:pp_valid_hist}).

\begin{figure}[H]
  \centering
  \includegraphics[width=0.75\textwidth]{processing/valid_frames_hist.png}
  \caption{Phân phối số khung hợp lệ mỗi video sau khi bỏ khung đơn sắc (trục tung thang log). Các video bên trái đường đứt nét có ít hơn $T=16$ khung hợp lệ.}
  \label{fig:pp_valid_hist}
\end{figure}

\paragraph{Bước 6 -- Lấy mẫu $T=16$ khung} Mô hình cần chuỗi đầu vào có độ dài cố định, nên mỗi video được thu gọn thành $T=16$ khung lấy đều theo thời gian trên danh sách khung hợp lệ:
\begin{equation}
  s_k = v_{\,1+\operatorname{round}\left(k\,\frac{m-1}{T-1}\right)}, \qquad k = 0, 1, \dots, T-1.
\end{equation}
Khung đầu và khung cuối của chuỗi luôn là khung hợp lệ đầu tiên và cuối cùng. Với video có $m<16$ khung hợp lệ, công thức trên tự lặp lại đều các khung (mỗi khung xuất hiện một hoặc hai lần) thay vì chỉ lặp khung cuối, nên chuỗi vẫn giữ đúng thứ tự thời gian. Nếu một video không có khung hợp lệ nào (chỉ có thể xảy ra ở dữ liệu kiểm tra hoặc demo) thì dùng toàn bộ khung.

Để thấy tác dụng của bước 5, nhóm so sánh với cách lấy mẫu đều trực tiếp trên toàn bộ khung, không lọc khung đơn sắc (Bảng~\ref{tab:pp_sampling}, Hình~\ref{fig:pp_sampling}). Nếu không lọc, cả 563 video có khung đơn sắc đều bị lẫn khung đơn sắc vào chuỗi đầu vào, trung bình 6,4/16 khung, và 222 video bị lẫn từ một nửa chuỗi trở lên; các số này khớp với mô phỏng ở bước EDA (Mục~\ref{subsec:eda_quality}). Khi lấy mẫu trên khung hợp lệ thì không còn khung đơn sắc nào trong đầu vào.

\begin{table}[H]
  \centering
  \caption{Số khung đơn sắc trong chuỗi 16 khung đầu vào, khi không lọc và khi lọc khung đơn sắc (2.772 video train sạch).}
  \label{tab:pp_sampling}
  \small
  \begin{tabular}{@{}lrr@{}}
    \toprule
    \textbf{Chỉ số} & \textbf{Không lọc} & \textbf{Có lọc} \\
    \midrule
    Video có $\geq$1 khung đơn sắc trong đầu vào   & 563 & 0 \\
    Số khung đơn sắc trung bình (trên các video đó) & 6,4/16 & 0 \\
    Video có $\geq$8/16 khung đơn sắc               & 222 & 0 \\
    \bottomrule
  \end{tabular}
\end{table}

\begin{figure}[H]
  \centering
  \includegraphics[width=\textwidth]{processing/sampling_naive_vs_filtered.png}
  \caption{Mười sáu khung đầu vào của video \textit{Nôn ói}/915809 (21/30 khung đơn sắc). Hàng trên: lấy mẫu trực tiếp, 11/16 khung là khung đơn sắc. Hàng dưới: lấy mẫu trên các khung hợp lệ $f_{21}$--$f_{29}$, các khung được lặp đều.}
  \label{fig:pp_sampling}
\end{figure}

% =====================================================================
\subsubsection{Chuẩn hoá}
\label{subsec:pp_norm}

\paragraph{Bước 7 -- Chuẩn hoá z-score} Giá trị điểm ảnh được đưa về $[0,1]$ rồi chuẩn hoá z-score trên từng kênh màu:
\begin{equation}
  x'_c = \frac{x_c/255 - \mu_c}{\sigma_c}, \qquad \mu = (0{,}485;\ 0{,}456;\ 0{,}406), \quad \sigma = (0{,}229;\ 0{,}224;\ 0{,}225),
\end{equation}
với $\mu,\sigma$ là trung bình và độ lệch chuẩn của tập ImageNet. Nhóm dùng thống kê của ImageNet thay vì tính lại trên tập train vì các mô hình ở Chương~\ref{sec:model} dùng mạng trích đặc trưng đã huấn luyện sẵn trên ImageNet; đầu vào cần có cùng phân phối với dữ liệu mà mạng đó đã học. Kết quả của mỗi video là một tensor kích thước $16\times3\times224\times224$. Độ phân giải đã đồng nhất $224\times224$ (Bảng~\ref{tab:eda_res}) nên không cần đổi kích thước.

% =====================================================================
\subsubsection{Chia tập huấn luyện và kiểm định}
\label{subsec:pp_split}

\paragraph{Bước 8 -- Chia phân tầng} Chỉ tập train có nhãn, nên tập kiểm định (\textit{validation}) phải tách ra từ chính tập này. Nhóm chia 2.772 video sạch thành 5 phần bằng \texttt{StratifiedKFold} của scikit-learn~\cite{sklearn} (\texttt{shuffle=True}, \texttt{random\_state=42}) và ghi số thứ tự phần vào cột \texttt{fold} của manifest. Phần \texttt{fold = 0} là tập kiểm định chung cho \emph{mọi} mô hình, các phần còn lại là tập huấn luyện; nhờ vậy kết quả của các mô hình so sánh được với nhau. Mô hình SVM có thể dùng cả 5 phần để kiểm định chéo (\textit{cross-validation}).

\begin{itemize}
  \item \textbf{Phân tầng theo nhãn:} mỗi phần giữ tỷ lệ các lớp gần giống toàn tập. Lớp nhỏ nhất có 5 video nên có đúng 1 video ở mỗi phần, nhờ đó cả 100 lớp đều có mặt ở tập kiểm định.
  \item \textbf{Không rò rỉ dữ liệu:} EDA khuyến nghị chia theo nhóm video gốc (\texttt{StratifiedGroupKFold}) để các bản sao không rơi vào cả hai phía. Vì bước 2 đã loại mọi bản sao, mỗi video sạch là một nhóm riêng, nên chia theo từng video cũng chính là chia theo nhóm. Notebook kiểm tra lại bằng \texttt{assert}: mã MD5 là duy nhất trong manifest và không có mã MD5 nào nằm ở cả tập huấn luyện lẫn tập kiểm định. Nhóm đã thử \texttt{StratifiedGroupKFold}, nhưng thuật toán này chỉ cân bằng gần đúng và tạo ra phần thiếu lớp.
\end{itemize}

\begin{table}[H]
  \centering
  \caption{Kết quả chia tập train sạch.}
  \label{tab:pp_split}
  \small
  \begin{tabular}{@{}lrrr@{}}
    \toprule
    \textbf{Tập} & \textbf{Số video} & \textbf{Số lớp} & \textbf{Số video mỗi lớp} \\
    \midrule
    Huấn luyện (\texttt{fold} 1--4) & 2.217 & 100 & 4--41 \\
    Kiểm định (\texttt{fold} 0)     & 555   & 100 & 1--10 \\
    \midrule
    Tổng                            & 2.772 & 100 & 5--51 \\
    \bottomrule
  \end{tabular}
\end{table}

% =====================================================================
\subsubsection{Xử lý mất cân bằng lớp}
\label{subsec:pp_weights}

\paragraph{Bước 9 -- Trọng số lớp} Trên tập huấn luyện, lớp nhiều nhất (\textit{Chạy}) có 41 video còn nhiều lớp chỉ có 4 video (\textit{Cách ly}, \textit{Ho}, \textit{Sốt}, \textit{Virus}, \textit{Xe máy},\ldots), chênh khoảng 10,2 lần. Nếu mọi mẫu có vai trò như nhau trong hàm mất mát, mô hình sẽ ưu tiên các lớp đông và Macro F1 giảm ở các lớp ít mẫu. Nhóm tính trọng số cho từng lớp $c$ theo kiểu \textit{balanced}:
\begin{equation}
  w_c = \frac{N}{K \cdot n_c},
\end{equation}
với $N=2.217$ là số video huấn luyện, $K=100$ là số lớp và $n_c$ là số video huấn luyện của lớp $c$. Lớp ít mẫu nhận trọng số lớn: trọng số dao động từ 0,54 (\textit{Chạy}) đến 5,54 (các lớp có 4 video), và một lớp trung bình ($n_c = N/K$) có trọng số bằng 1. Trọng số chỉ tính trên tập huấn luyện để không dùng thông tin của tập kiểm định. Các mô hình dùng trọng số này trong hàm mất mát (ví dụ tham số \texttt{weight} của \texttt{CrossEntropyLoss}) hoặc tham số \texttt{class\_weight} của SVM.

% =====================================================================
\subsubsection{Tăng cường dữ liệu}
\label{subsec:pp_aug}

\paragraph{Bước 10 -- Tăng cường} Tập huấn luyện nhỏ (trung bình khoảng 22 video mỗi lớp), nên mỗi lần đọc một video huấn luyện, \texttt{preprocess.py} tạo ra một biến thể ngẫu nhiên khác của video đó. Ba phép biến đổi được dùng, với cùng một tham số cho mọi khung của một video để chuyển động không bị méo:

\begin{itemize}
  \item \textbf{Cắt đoạn theo thời gian:} giữ một đoạn liên tục dài 80--100\% danh sách khung hợp lệ ở vị trí ngẫu nhiên, rồi mới lấy 16 khung. Phép này mô phỏng việc người ký bắt đầu, kết thúc sớm hay muộn và ký nhanh hay chậm khác nhau.
  \item \textbf{Cắt vùng ảnh ngẫu nhiên:} cắt một vùng có cạnh bằng 85--100\% cạnh khung hình ở vị trí ngẫu nhiên rồi phóng lại về $224\times224$. Phép này mô phỏng khác biệt về vị trí người ký và khoảng cách tới máy quay.
  \item \textbf{Thay đổi màu:} nhân độ sáng, độ tương phản và độ bão hoà với các hệ số ngẫu nhiên trong khoảng $[0{,}75;\ 1{,}25]$, giúp mô hình bền vững với khác biệt ánh sáng và màu nền giữa train và test (Mục~\ref{subsec:eda_shift}).
\end{itemize}

Nhóm \emph{không} dùng lật ngang vì tay trái và tay phải có thể làm đổi nghĩa ký hiệu. Tăng cường chỉ áp dụng cho tập huấn luyện; tập kiểm định và tập kiểm tra được đưa vào mô hình nguyên bản để kết quả đánh giá ổn định. Hình~\ref{fig:pp_aug} minh hoạ ba biến thể tạo ra từ cùng một video.

\begin{figure}[H]
  \centering
  \includegraphics[width=0.9\textwidth]{processing/augmentation_examples.png}
  \caption{Một video lớp \textit{An ủi} (hàng đầu) và ba biến thể tăng cường ngẫu nhiên; mỗi hàng hiển thị 8 trong 16 khung đầu vào.}
  \label{fig:pp_aug}
\end{figure}

% =====================================================================
\subsubsection{Dữ liệu kiểm tra}
\label{subsec:pp_test}

Hai tập kiểm tra (\texttt{public\_test}, \texttt{private\_test}) không có nhãn và cuộc thi yêu cầu dự đoán cho \emph{mọi} tệp, nên nhóm không loại tệp nào. Notebook chỉ lập bảng \texttt{test\_manifest.csv} gồm tên tệp, mã MD5 và cột \texttt{dup\_of\_train} đánh dấu video trùng với tập train: có 157/1.630 video public test trùng với train và 0/2.859 video private test trùng với train. Các video trùng vẫn được mô hình dự đoán bình thường; cột này chỉ dùng khi diễn giải điểm public test (Mục~\ref{subsec:eda_quality}).

Khi suy luận, mỗi video kiểm tra đi qua đúng các bước 5, 6 và 7 như video kiểm định: tự tìm khung đơn sắc, lấy 16 khung trên các khung hợp lệ và chuẩn hoá, không tăng cường. Notebook chạy thử hàm đọc video này trên 8 video kiểm tra ngẫu nhiên để xác nhận đầu ra có đúng kích thước $16\times224\times224\times3$.

% =====================================================================
\subsubsection{Kết quả và đầu ra}
\label{subsec:pp_output}

Bảng~\ref{tab:pp_outputs} liệt kê các tệp do notebook tạo ra trong thư mục \texttt{processing/output/}. Các mô hình chỉ cần \texttt{preprocess.py}, \texttt{manifest\_clean.csv}, \texttt{test\_manifest.csv} và \texttt{class\_weights.csv} (vài trăm KB) cùng tệp dữ liệu gốc, không cần chạy lại bước làm sạch.

\begin{table}[H]
  \centering
  \caption{Các tệp đầu ra của bước tiền xử lý.}
  \label{tab:pp_outputs}
  \small
  \begin{tabular}{@{}p{0.3\textwidth}p{0.62\textwidth}@{}}
    \toprule
    \textbf{Tệp} & \textbf{Nội dung} \\
    \midrule
    \texttt{manifest\_clean.csv} & 2.772 video train sạch: lớp, nhãn số, tên tệp, đường dẫn, MD5, số khung, số khung hợp lệ, chỉ số khung hợp lệ (\texttt{valid\_idx}), phần (\texttt{fold}) \\
    \texttt{removed.csv} & 1.103 tệp bị loại, kèm bước và lý do \\
    \texttt{test\_manifest.csv} & 4.489 video kiểm tra: tập, tên tệp, MD5, \texttt{dup\_of\_train} \\
    \texttt{class\_weights.csv} & Trọng số 100 lớp tính trên tập huấn luyện \\
    \texttt{processing\_summary.csv} & Số video còn lại sau từng bước (Bảng~\ref{tab:pp_flow}) \\
    \texttt{class\_counts\_before\_after.csv} & Số video mỗi lớp trước và sau làm sạch \\
    \texttt{processing\_config.json} & Tham số đã dùng (\texttt{SEED}, số phần, $T$, ngưỡng đơn sắc) \\
    \texttt{figures/} & Các hình trong mục này \\
    \bottomrule
  \end{tabular}
\end{table}

Tóm lại, sau tiền xử lý tập train không còn bản sao, video rỗng hay nhãn mâu thuẫn; không video nào còn khung đơn sắc trong chuỗi đầu vào; tập kiểm định đủ 100 lớp và không chứa video trùng với tập huấn luyện. Các bước này xử lý cả ba vấn đề chính mà EDA đã nêu: bản sao gây rò rỉ dữ liệu, khung đơn sắc chỉ có ở train và mất cân bằng lớp.

\paragraph{Hạn chế} Nhóm chưa xác định được nhãn đúng cho hai nội dung xung đột nên phải bỏ chúng, làm giảm 5 tệp dữ liệu có thể hữu ích. Ngưỡng đơn sắc chỉ phát hiện khung một màu hoàn toàn, chưa phát hiện khung mờ hoặc khung chỉ có nền. Khử trùng lặp bằng MD5 chỉ tìm được bản sao giống hệt từng byte, không tìm được video gần giống nhau (ví dụ cùng cảnh nhưng nén lại). Cuối cùng, tập kiểm định chỉ có 1--10 video mỗi lớp, nên F1 của từng lớp riêng lẻ trên tập này còn dao động lớn.
]
% Mã nguồn: processing/data_processing.ipynb, processing/preprocess.py
% Hình: image/processing/*.png (sinh từ processing/output/figures/)
%   3.2.1 Tổng quan quy trình
%   3.2.2 Làm sạch dữ liệu
%   3.2.3 Xử lý khung nhiễu và lấy mẫu theo thời gian
%   3.2.4 Chuẩn hoá
%   3.2.5 Chia tập huấn luyện và kiểm định
%   3.2.6 Xử lý mất cân bằng lớp
%   3.2.7 Tăng cường dữ liệu
%   3.2.8 Dữ liệu kiểm tra
%   3.2.9 Kết quả và đầu ra
% =====================================================================

\subsection{Tiền xử lý dữ liệu}
\label{sec:preprocess}

Mục này trình bày các bước biến dữ liệu thô thành đầu vào sạch và thống nhất cho mọi mô hình ở Chương~\ref{sec:model}. Mỗi bước xuất phát từ một vấn đề đã phát hiện ở bước EDA (Bảng~\ref{tab:eda_reco}). Nhóm dùng Python với \texttt{pandas}, \texttt{numpy}, \texttt{OpenCV} và \texttt{scikit-learn}; mã nguồn gồm hai phần:

\begin{itemize}
  \item \textbf{Notebook làm sạch} \texttt{processing/data\_processing.ipynb}: chạy \emph{một lần} trên toàn bộ tập train để kiểm tra, loại bỏ và đánh dấu dữ liệu lỗi, chia tập kiểm định và tính trọng số lớp. Kết quả được lưu thành một bảng chỉ mục (\textit{manifest}) dùng chung cho cả nhóm.
  \item \textbf{Module dùng chung} \texttt{processing/preprocess.py}: các phép biến đổi áp dụng \emph{mỗi lần} mô hình đọc một video (bỏ khung đơn sắc, lấy mẫu khung, tăng cường, chuẩn hoá). Mô hình SVM, CNN-LSTM, Vision Transformer và ứng dụng demo cùng gọi module này, nên dữ liệu ở huấn luyện, kiểm định, kiểm tra và demo được xử lý giống hệt nhau.
\end{itemize}

Nhóm không lưu sẵn khung hình đã xử lý ra đĩa vì hai lý do: lưu 16 khung $224\times224\times3$ cho 2.772 video tốn khoảng 6,7~GB, và phép tăng cường cần tạo ra phiên bản khác nhau ở mỗi lần đọc. Mọi bước có yếu tố ngẫu nhiên đều cố định hạt giống \texttt{SEED = 42} để kết quả tái lập được.

% =====================================================================
\subsubsection{Tổng quan quy trình}
\label{subsec:pp_overview}

Bảng~\ref{tab:pp_overview} liệt kê mười bước tiền xử lý theo thứ tự thực hiện và xếp mỗi bước vào nhóm tác vụ tiền xử lý dữ liệu tương ứng: làm sạch (\textit{data cleaning}), tích hợp (\textit{data integration}), thu gọn (\textit{data reduction}) và biến đổi dữ liệu (\textit{data transformation})~\cite{han2011}.

\begin{table}[H]
  \centering
  \caption{Các bước tiền xử lý, nhóm tác vụ và nơi thực hiện.}
  \label{tab:pp_overview}
  \small
  \begin{tabular}{@{}cp{0.33\textwidth}p{0.3\textwidth}p{0.2\textwidth}@{}}
    \toprule
    \textbf{Bước} & \textbf{Việc} & \textbf{Nhóm tác vụ} & \textbf{Thực hiện ở} \\
    \midrule
    1  & Kiểm tra toàn vẹn, chuẩn hoá tên lớp & Làm sạch: dữ liệu thiếu, lỗi & Notebook \\
    2  & Khử trùng lặp theo MD5 & Tích hợp: dư thừa dữ liệu & Notebook \\
    3  & Loại video xung đột nhãn & Làm sạch: dữ liệu không nhất quán & Notebook \\
    4  & Loại video rỗng & Làm sạch: dữ liệu thiếu & Notebook \\
    5  & Đánh dấu khung đơn sắc & Làm sạch: dữ liệu nhiễu & Notebook, \texttt{preprocess.py} \\
    6  & Lấy đều $T=16$ khung trên các khung hợp lệ & Thu gọn: lấy mẫu & \texttt{preprocess.py} \\
    7  & Chuẩn hoá z-score theo kênh màu & Biến đổi: chuẩn hoá & \texttt{preprocess.py} \\
    8  & Chia train/validation phân tầng & Đánh giá mô hình & Notebook \\
    9  & Trọng số lớp & Xử lý mất cân bằng lớp & Notebook \\
    10 & Tăng cường dữ liệu (chỉ khi huấn luyện) & Sinh thêm dữ liệu & \texttt{preprocess.py} \\
    \bottomrule
  \end{tabular}
\end{table}

Bước 1--5 thu hẹp tập train từ 3.875 tệp xuống 2.772 video sạch (Bảng~\ref{tab:pp_flow}). Bước 6, 7 và 10 không bỏ video nào mà quyết định video được đưa vào mô hình dưới dạng nào. Bước 8 và 9 chuẩn bị cho huấn luyện và đánh giá.

% =====================================================================
\subsubsection{Làm sạch dữ liệu}
\label{subsec:pp_clean}

\paragraph{Quét toàn bộ video} Notebook mở từng video bằng \texttt{cv2.VideoCapture}, đọc \emph{tất cả} khung hình và ghi lại: mã băm MD5 của tệp, số khung, và danh sách chỉ số các khung không đơn sắc (Mục~\ref{subsec:pp_frames}). Việc quét được chạy song song trên nhiều lõi CPU (\texttt{ProcessPoolExecutor}). Để kiểm tra bước quét, nhóm đối chiếu số khung đơn sắc của từng video với cột \texttt{solid\_frames} trong \texttt{video\_quality\_flags.csv} của bước EDA: cả 3.875/3.875 tệp đều khớp.

\paragraph{Bước 1 -- Kiểm tra toàn vẹn} Cả 3.875 tệp đều đọc được và có ít nhất 16 khung, nên không có tệp nào bị loại ở bước này. Tên lớp ở thư mục và khoá trong \texttt{label\_mapping.pkl} được chuẩn hoá về dạng Unicode NFC trước khi so khớp (Mục~\ref{subsec:eda_structure}); sau chuẩn hoá, 100/100 lớp có nhãn số 0--99.

\paragraph{Bước 2 -- Khử trùng lặp} Hai tệp có cùng mã MD5 là cùng một video. Với mỗi mã MD5, notebook giữ lại một tệp (ưu tiên tệp không có hậu tố \texttt{\_N}) và loại các tệp còn lại. Bước này loại 1.099 tệp, trong đó 1.069 tệp mang hậu tố \texttt{\_N}, và còn \textbf{2.776 video duy nhất}, khớp với kết quả EDA. Từ đây mỗi video trong tập sạch có một mã MD5 riêng; tính chất này giúp chia tập kiểm định không bị rò rỉ (Mục~\ref{subsec:pp_split}).

\paragraph{Bước 3 -- Xung đột nhãn} Hai nội dung video vừa nằm ở lớp \textit{Chào} vừa nằm ở lớp \textit{Chậm lại} (5 tệp: \texttt{Chào/392683}, \texttt{Chào/468203}, \texttt{Chào/792094}, \texttt{Chậm lại/254739}, \texttt{Chậm lại/284430}). Nhóm không xác minh được nhãn nào đúng, và giữ một video với nhãn có thể sai sẽ đưa nhiễu nhãn vào cả huấn luyện lẫn đánh giá, nên loại toàn bộ hai nội dung này. Notebook có sẵn biến \texttt{KEEP\_LABEL} để giữ lại các video này với nhãn đã xác minh nếu sau này có thông tin. Sau bước này còn 2.774 video. Hai lớp bị ảnh hưởng vẫn còn nhiều mẫu (\textit{Chào} 22 video, \textit{Chậm lại} 44 video).

\paragraph{Bước 4 -- Video rỗng} Hai video không có khung hợp lệ nào (\texttt{Biết/214576}: 16/16 khung đơn sắc; \texttt{Khai báo/885339}: 23/23 khung đơn sắc) bị loại vì không chứa thông tin ký hiệu. Tập train sạch còn \textbf{2.772 video} và vẫn đủ 100 lớp.

\begin{table}[H]
  \centering
  \caption{Số video còn lại sau từng bước làm sạch tập train.}
  \label{tab:pp_flow}
  \small
  \begin{tabular}{@{}p{0.3\textwidth}rp{0.45\textwidth}@{}}
    \toprule
    \textbf{Bước} & \textbf{Còn lại} & \textbf{Ghi chú} \\
    \midrule
    Ban đầu                          & 3.875 & Toàn bộ tệp trong \texttt{train/} \\
    1 -- Kiểm tra toàn vẹn           & 3.875 & Không có tệp lỗi hoặc thiếu khung \\
    2 -- Khử trùng lặp MD5           & 2.776 & Loại 1.099 bản sao \\
    3 -- Xung đột nhãn               & 2.774 & Loại 2 nội dung (5 tệp) \textit{Chào} $\leftrightarrow$ \textit{Chậm lại} \\
    4 -- Video rỗng                  & \textbf{2.772} & Loại \texttt{Biết/214576}, \texttt{Khai báo/885339} \\
    5 -- Khung đơn sắc               & 2.772 & 563 video có khung bị bỏ qua, không loại video \\
    \bottomrule
  \end{tabular}
\end{table}

Danh sách 1.103 tệp bị loại cùng lý do được lưu ở \texttt{processing/output/removed.csv}, để có thể kiểm tra lại từng trường hợp. Hình~\ref{fig:pp_class_counts} so sánh số video mỗi lớp trước và sau khi làm sạch: sau làm sạch mỗi lớp có từ 5 đến 51 video (trung bình 27,7; trung vị 32,5).

\begin{figure}[H]
  \centering
  \includegraphics[width=\textwidth]{processing/class_counts_before_after.png}
  \caption{Số video mỗi lớp trước (xám) và sau (xanh) khi làm sạch, sắp xếp theo số video sau làm sạch.}
  \label{fig:pp_class_counts}
\end{figure}

% =====================================================================
\subsubsection{Xử lý khung nhiễu và lấy mẫu theo thời gian}
\label{subsec:pp_frames}

\paragraph{Bước 5 -- Đánh dấu khung đơn sắc} Gọi $f_i$ là khung thứ $i$ của video và $\sigma_c(f_i)$ là độ lệch chuẩn giá trị điểm ảnh trên kênh màu $c\in\{R,G,B\}$. Khung $f_i$ được coi là \emph{đơn sắc} nếu
\begin{equation}
  \max_{c\in\{R,G,B\}} \sigma_c(f_i) < 3,
\end{equation}
tức gần như mọi điểm ảnh có cùng một màu; ngưỡng này giống với ngưỡng ở bước EDA. Các khung còn lại là \emph{khung hợp lệ}. Thay vì xoá khung khỏi tệp video, notebook lưu danh sách chỉ số khung hợp lệ $V=(v_1<v_2<\dots<v_m)$ của từng video vào cột \texttt{valid\_idx} của manifest; khi đọc video, \texttt{preprocess.py} chỉ lấy mẫu trong danh sách này.

Trong 2.772 video sạch, 563 video (20,3\%, thuộc 78 lớp) có khung đơn sắc, trung bình 12,2 khung mỗi video. Sau khi bỏ khung đơn sắc, mỗi video còn từ 9 đến 30 khung hợp lệ (trung vị 30); 222 video còn ít hơn 16 khung hợp lệ (Hình~\ref{fig:pp_valid_hist}).

\begin{figure}[H]
  \centering
  \includegraphics[width=0.75\textwidth]{processing/valid_frames_hist.png}
  \caption{Phân phối số khung hợp lệ mỗi video sau khi bỏ khung đơn sắc (trục tung thang log). Các video bên trái đường đứt nét có ít hơn $T=16$ khung hợp lệ.}
  \label{fig:pp_valid_hist}
\end{figure}

\paragraph{Bước 6 -- Lấy mẫu $T=16$ khung} Mô hình cần chuỗi đầu vào có độ dài cố định, nên mỗi video được thu gọn thành $T=16$ khung lấy đều theo thời gian trên danh sách khung hợp lệ:
\begin{equation}
  s_k = v_{\,1+\operatorname{round}\left(k\,\frac{m-1}{T-1}\right)}, \qquad k = 0, 1, \dots, T-1.
\end{equation}
Khung đầu và khung cuối của chuỗi luôn là khung hợp lệ đầu tiên và cuối cùng. Với video có $m<16$ khung hợp lệ, công thức trên tự lặp lại đều các khung (mỗi khung xuất hiện một hoặc hai lần) thay vì chỉ lặp khung cuối, nên chuỗi vẫn giữ đúng thứ tự thời gian. Nếu một video không có khung hợp lệ nào (chỉ có thể xảy ra ở dữ liệu kiểm tra hoặc demo) thì dùng toàn bộ khung.

Để thấy tác dụng của bước 5, nhóm so sánh với cách lấy mẫu đều trực tiếp trên toàn bộ khung, không lọc khung đơn sắc (Bảng~\ref{tab:pp_sampling}, Hình~\ref{fig:pp_sampling}). Nếu không lọc, cả 563 video có khung đơn sắc đều bị lẫn khung đơn sắc vào chuỗi đầu vào, trung bình 6,4/16 khung, và 222 video bị lẫn từ một nửa chuỗi trở lên; các số này khớp với mô phỏng ở bước EDA (Mục~\ref{subsec:eda_quality}). Khi lấy mẫu trên khung hợp lệ thì không còn khung đơn sắc nào trong đầu vào.

\begin{table}[H]
  \centering
  \caption{Số khung đơn sắc trong chuỗi 16 khung đầu vào, khi không lọc và khi lọc khung đơn sắc (2.772 video train sạch).}
  \label{tab:pp_sampling}
  \small
  \begin{tabular}{@{}lrr@{}}
    \toprule
    \textbf{Chỉ số} & \textbf{Không lọc} & \textbf{Có lọc} \\
    \midrule
    Video có $\geq$1 khung đơn sắc trong đầu vào   & 563 & 0 \\
    Số khung đơn sắc trung bình (trên các video đó) & 6,4/16 & 0 \\
    Video có $\geq$8/16 khung đơn sắc               & 222 & 0 \\
    \bottomrule
  \end{tabular}
\end{table}

\begin{figure}[H]
  \centering
  \includegraphics[width=\textwidth]{processing/sampling_naive_vs_filtered.png}
  \caption{Mười sáu khung đầu vào của video \textit{Nôn ói}/915809 (21/30 khung đơn sắc). Hàng trên: lấy mẫu trực tiếp, 11/16 khung là khung đơn sắc. Hàng dưới: lấy mẫu trên các khung hợp lệ $f_{21}$--$f_{29}$, các khung được lặp đều.}
  \label{fig:pp_sampling}
\end{figure}

% =====================================================================
\subsubsection{Chuẩn hoá}
\label{subsec:pp_norm}

\paragraph{Bước 7 -- Chuẩn hoá z-score} Giá trị điểm ảnh được đưa về $[0,1]$ rồi chuẩn hoá z-score trên từng kênh màu:
\begin{equation}
  x'_c = \frac{x_c/255 - \mu_c}{\sigma_c}, \qquad \mu = (0{,}485;\ 0{,}456;\ 0{,}406), \quad \sigma = (0{,}229;\ 0{,}224;\ 0{,}225),
\end{equation}
với $\mu,\sigma$ là trung bình và độ lệch chuẩn của tập ImageNet. Nhóm dùng thống kê của ImageNet thay vì tính lại trên tập train vì các mô hình ở Chương~\ref{sec:model} dùng mạng trích đặc trưng đã huấn luyện sẵn trên ImageNet; đầu vào cần có cùng phân phối với dữ liệu mà mạng đó đã học. Kết quả của mỗi video là một tensor kích thước $16\times3\times224\times224$. Độ phân giải đã đồng nhất $224\times224$ (Bảng~\ref{tab:eda_res}) nên không cần đổi kích thước.

% =====================================================================
\subsubsection{Chia tập huấn luyện và kiểm định}
\label{subsec:pp_split}

\paragraph{Bước 8 -- Chia phân tầng} Chỉ tập train có nhãn, nên tập kiểm định (\textit{validation}) phải tách ra từ chính tập này. Nhóm chia 2.772 video sạch thành 5 phần bằng \texttt{StratifiedKFold} của scikit-learn~\cite{sklearn} (\texttt{shuffle=True}, \texttt{random\_state=42}) và ghi số thứ tự phần vào cột \texttt{fold} của manifest. Phần \texttt{fold = 0} là tập kiểm định chung cho \emph{mọi} mô hình, các phần còn lại là tập huấn luyện; nhờ vậy kết quả của các mô hình so sánh được với nhau. Mô hình SVM có thể dùng cả 5 phần để kiểm định chéo (\textit{cross-validation}).

\begin{itemize}
  \item \textbf{Phân tầng theo nhãn:} mỗi phần giữ tỷ lệ các lớp gần giống toàn tập. Lớp nhỏ nhất có 5 video nên có đúng 1 video ở mỗi phần, nhờ đó cả 100 lớp đều có mặt ở tập kiểm định.
  \item \textbf{Không rò rỉ dữ liệu:} EDA khuyến nghị chia theo nhóm video gốc (\texttt{StratifiedGroupKFold}) để các bản sao không rơi vào cả hai phía. Vì bước 2 đã loại mọi bản sao, mỗi video sạch là một nhóm riêng, nên chia theo từng video cũng chính là chia theo nhóm. Notebook kiểm tra lại bằng \texttt{assert}: mã MD5 là duy nhất trong manifest và không có mã MD5 nào nằm ở cả tập huấn luyện lẫn tập kiểm định. Nhóm đã thử \texttt{StratifiedGroupKFold}, nhưng thuật toán này chỉ cân bằng gần đúng và tạo ra phần thiếu lớp.
\end{itemize}

\begin{table}[H]
  \centering
  \caption{Kết quả chia tập train sạch.}
  \label{tab:pp_split}
  \small
  \begin{tabular}{@{}lrrr@{}}
    \toprule
    \textbf{Tập} & \textbf{Số video} & \textbf{Số lớp} & \textbf{Số video mỗi lớp} \\
    \midrule
    Huấn luyện (\texttt{fold} 1--4) & 2.217 & 100 & 4--41 \\
    Kiểm định (\texttt{fold} 0)     & 555   & 100 & 1--10 \\
    \midrule
    Tổng                            & 2.772 & 100 & 5--51 \\
    \bottomrule
  \end{tabular}
\end{table}

% =====================================================================
\subsubsection{Xử lý mất cân bằng lớp}
\label{subsec:pp_weights}

\paragraph{Bước 9 -- Trọng số lớp} Trên tập huấn luyện, lớp nhiều nhất (\textit{Chạy}) có 41 video còn nhiều lớp chỉ có 4 video (\textit{Cách ly}, \textit{Ho}, \textit{Sốt}, \textit{Virus}, \textit{Xe máy},\ldots), chênh khoảng 10,2 lần. Nếu mọi mẫu có vai trò như nhau trong hàm mất mát, mô hình sẽ ưu tiên các lớp đông và Macro F1 giảm ở các lớp ít mẫu. Nhóm tính trọng số cho từng lớp $c$ theo kiểu \textit{balanced}:
\begin{equation}
  w_c = \frac{N}{K \cdot n_c},
\end{equation}
với $N=2.217$ là số video huấn luyện, $K=100$ là số lớp và $n_c$ là số video huấn luyện của lớp $c$. Lớp ít mẫu nhận trọng số lớn: trọng số dao động từ 0,54 (\textit{Chạy}) đến 5,54 (các lớp có 4 video), và một lớp trung bình ($n_c = N/K$) có trọng số bằng 1. Trọng số chỉ tính trên tập huấn luyện để không dùng thông tin của tập kiểm định. Các mô hình dùng trọng số này trong hàm mất mát (ví dụ tham số \texttt{weight} của \texttt{CrossEntropyLoss}) hoặc tham số \texttt{class\_weight} của SVM.

% =====================================================================
\subsubsection{Tăng cường dữ liệu}
\label{subsec:pp_aug}

\paragraph{Bước 10 -- Tăng cường} Tập huấn luyện nhỏ (trung bình khoảng 22 video mỗi lớp), nên mỗi lần đọc một video huấn luyện, \texttt{preprocess.py} tạo ra một biến thể ngẫu nhiên khác của video đó. Ba phép biến đổi được dùng, với cùng một tham số cho mọi khung của một video để chuyển động không bị méo:

\begin{itemize}
  \item \textbf{Cắt đoạn theo thời gian:} giữ một đoạn liên tục dài 80--100\% danh sách khung hợp lệ ở vị trí ngẫu nhiên, rồi mới lấy 16 khung. Phép này mô phỏng việc người ký bắt đầu, kết thúc sớm hay muộn và ký nhanh hay chậm khác nhau.
  \item \textbf{Cắt vùng ảnh ngẫu nhiên:} cắt một vùng có cạnh bằng 85--100\% cạnh khung hình ở vị trí ngẫu nhiên rồi phóng lại về $224\times224$. Phép này mô phỏng khác biệt về vị trí người ký và khoảng cách tới máy quay.
  \item \textbf{Thay đổi màu:} nhân độ sáng, độ tương phản và độ bão hoà với các hệ số ngẫu nhiên trong khoảng $[0{,}75;\ 1{,}25]$, giúp mô hình bền vững với khác biệt ánh sáng và màu nền giữa train và test (Mục~\ref{subsec:eda_shift}).
\end{itemize}

Nhóm \emph{không} dùng lật ngang vì tay trái và tay phải có thể làm đổi nghĩa ký hiệu. Tăng cường chỉ áp dụng cho tập huấn luyện; tập kiểm định và tập kiểm tra được đưa vào mô hình nguyên bản để kết quả đánh giá ổn định. Hình~\ref{fig:pp_aug} minh hoạ ba biến thể tạo ra từ cùng một video.

\begin{figure}[H]
  \centering
  \includegraphics[width=0.9\textwidth]{processing/augmentation_examples.png}
  \caption{Một video lớp \textit{An ủi} (hàng đầu) và ba biến thể tăng cường ngẫu nhiên; mỗi hàng hiển thị 8 trong 16 khung đầu vào.}
  \label{fig:pp_aug}
\end{figure}

% =====================================================================
\subsubsection{Dữ liệu kiểm tra}
\label{subsec:pp_test}

Hai tập kiểm tra (\texttt{public\_test}, \texttt{private\_test}) không có nhãn và cuộc thi yêu cầu dự đoán cho \emph{mọi} tệp, nên nhóm không loại tệp nào. Notebook chỉ lập bảng \texttt{test\_manifest.csv} gồm tên tệp, mã MD5 và cột \texttt{dup\_of\_train} đánh dấu video trùng với tập train: có 157/1.630 video public test trùng với train và 0/2.859 video private test trùng với train. Các video trùng vẫn được mô hình dự đoán bình thường; cột này chỉ dùng khi diễn giải điểm public test (Mục~\ref{subsec:eda_quality}).

Khi suy luận, mỗi video kiểm tra đi qua đúng các bước 5, 6 và 7 như video kiểm định: tự tìm khung đơn sắc, lấy 16 khung trên các khung hợp lệ và chuẩn hoá, không tăng cường. Notebook chạy thử hàm đọc video này trên 8 video kiểm tra ngẫu nhiên để xác nhận đầu ra có đúng kích thước $16\times224\times224\times3$.

% =====================================================================
\subsubsection{Kết quả và đầu ra}
\label{subsec:pp_output}

Bảng~\ref{tab:pp_outputs} liệt kê các tệp do notebook tạo ra trong thư mục \texttt{processing/output/}. Các mô hình chỉ cần \texttt{preprocess.py}, \texttt{manifest\_clean.csv}, \texttt{test\_manifest.csv} và \texttt{class\_weights.csv} (vài trăm KB) cùng tệp dữ liệu gốc, không cần chạy lại bước làm sạch.

\begin{table}[H]
  \centering
  \caption{Các tệp đầu ra của bước tiền xử lý.}
  \label{tab:pp_outputs}
  \small
  \begin{tabular}{@{}p{0.3\textwidth}p{0.62\textwidth}@{}}
    \toprule
    \textbf{Tệp} & \textbf{Nội dung} \\
    \midrule
    \texttt{manifest\_clean.csv} & 2.772 video train sạch: lớp, nhãn số, tên tệp, đường dẫn, MD5, số khung, số khung hợp lệ, chỉ số khung hợp lệ (\texttt{valid\_idx}), phần (\texttt{fold}) \\
    \texttt{removed.csv} & 1.103 tệp bị loại, kèm bước và lý do \\
    \texttt{test\_manifest.csv} & 4.489 video kiểm tra: tập, tên tệp, MD5, \texttt{dup\_of\_train} \\
    \texttt{class\_weights.csv} & Trọng số 100 lớp tính trên tập huấn luyện \\
    \texttt{processing\_summary.csv} & Số video còn lại sau từng bước (Bảng~\ref{tab:pp_flow}) \\
    \texttt{class\_counts\_before\_after.csv} & Số video mỗi lớp trước và sau làm sạch \\
    \texttt{processing\_config.json} & Tham số đã dùng (\texttt{SEED}, số phần, $T$, ngưỡng đơn sắc) \\
    \texttt{figures/} & Các hình trong mục này \\
    \bottomrule
  \end{tabular}
\end{table}

Tóm lại, sau tiền xử lý tập train không còn bản sao, video rỗng hay nhãn mâu thuẫn; không video nào còn khung đơn sắc trong chuỗi đầu vào; tập kiểm định đủ 100 lớp và không chứa video trùng với tập huấn luyện. Các bước này xử lý cả ba vấn đề chính mà EDA đã nêu: bản sao gây rò rỉ dữ liệu, khung đơn sắc chỉ có ở train và mất cân bằng lớp.

\paragraph{Hạn chế} Nhóm chưa xác định được nhãn đúng cho hai nội dung xung đột nên phải bỏ chúng, làm giảm 5 tệp dữ liệu có thể hữu ích. Ngưỡng đơn sắc chỉ phát hiện khung một màu hoàn toàn, chưa phát hiện khung mờ hoặc khung chỉ có nền. Khử trùng lặp bằng MD5 chỉ tìm được bản sao giống hệt từng byte, không tìm được video gần giống nhau (ví dụ cùng cảnh nhưng nén lại). Cuối cùng, tập kiểm định chỉ có 1--10 video mỗi lớp, nên F1 của từng lớp riêng lẻ trên tập này còn dao động lớn.
]
% Mã nguồn: processing/data_processing.ipynb, processing/preprocess.py
% Hình: image/processing/*.png (sinh từ processing/output/figures/)
%   3.2.1 Tổng quan quy trình
%   3.2.2 Làm sạch dữ liệu
%   3.2.3 Xử lý khung nhiễu và lấy mẫu theo thời gian
%   3.2.4 Chuẩn hoá
%   3.2.5 Chia tập huấn luyện và kiểm định
%   3.2.6 Xử lý mất cân bằng lớp
%   3.2.7 Tăng cường dữ liệu
%   3.2.8 Dữ liệu kiểm tra
%   3.2.9 Kết quả và đầu ra
% =====================================================================

\subsection{Tiền xử lý dữ liệu}
\label{sec:preprocess}

Mục này trình bày các bước biến dữ liệu thô thành đầu vào sạch và thống nhất cho mọi mô hình ở Chương~\ref{sec:model}. Mỗi bước xuất phát từ một vấn đề đã phát hiện ở bước EDA (Bảng~\ref{tab:eda_reco}). Nhóm dùng Python với \texttt{pandas}, \texttt{numpy}, \texttt{OpenCV} và \texttt{scikit-learn}; mã nguồn gồm hai phần:

\begin{itemize}
  \item \textbf{Notebook làm sạch} \texttt{processing/data\_processing.ipynb}: chạy \emph{một lần} trên toàn bộ tập train để kiểm tra, loại bỏ và đánh dấu dữ liệu lỗi, chia tập kiểm định và tính trọng số lớp. Kết quả được lưu thành một bảng chỉ mục (\textit{manifest}) dùng chung cho cả nhóm.
  \item \textbf{Module dùng chung} \texttt{processing/preprocess.py}: các phép biến đổi áp dụng \emph{mỗi lần} mô hình đọc một video (bỏ khung đơn sắc, lấy mẫu khung, tăng cường, chuẩn hoá). Mô hình SVM, CNN-LSTM, Vision Transformer và ứng dụng demo cùng gọi module này, nên dữ liệu ở huấn luyện, kiểm định, kiểm tra và demo được xử lý giống hệt nhau.
\end{itemize}

Nhóm không lưu sẵn khung hình đã xử lý ra đĩa vì hai lý do: lưu 16 khung $224\times224\times3$ cho 2.772 video tốn khoảng 6,7~GB, và phép tăng cường cần tạo ra phiên bản khác nhau ở mỗi lần đọc. Mọi bước có yếu tố ngẫu nhiên đều cố định hạt giống \texttt{SEED = 42} để kết quả tái lập được.

% =====================================================================
\subsubsection{Tổng quan quy trình}
\label{subsec:pp_overview}

Bảng~\ref{tab:pp_overview} liệt kê mười bước tiền xử lý theo thứ tự thực hiện và xếp mỗi bước vào nhóm tác vụ tiền xử lý dữ liệu tương ứng: làm sạch (\textit{data cleaning}), tích hợp (\textit{data integration}), thu gọn (\textit{data reduction}) và biến đổi dữ liệu (\textit{data transformation})~\cite{han2011}.

\begin{table}[H]
  \centering
  \caption{Các bước tiền xử lý, nhóm tác vụ và nơi thực hiện.}
  \label{tab:pp_overview}
  \small
  \begin{tabular}{@{}cp{0.33\textwidth}p{0.3\textwidth}p{0.2\textwidth}@{}}
    \toprule
    \textbf{Bước} & \textbf{Việc} & \textbf{Nhóm tác vụ} & \textbf{Thực hiện ở} \\
    \midrule
    1  & Kiểm tra toàn vẹn, chuẩn hoá tên lớp & Làm sạch: dữ liệu thiếu, lỗi & Notebook \\
    2  & Khử trùng lặp theo MD5 & Tích hợp: dư thừa dữ liệu & Notebook \\
    3  & Loại video xung đột nhãn & Làm sạch: dữ liệu không nhất quán & Notebook \\
    4  & Loại video rỗng & Làm sạch: dữ liệu thiếu & Notebook \\
    5  & Đánh dấu khung đơn sắc & Làm sạch: dữ liệu nhiễu & Notebook, \texttt{preprocess.py} \\
    6  & Lấy đều $T=16$ khung trên các khung hợp lệ & Thu gọn: lấy mẫu & \texttt{preprocess.py} \\
    7  & Chuẩn hoá z-score theo kênh màu & Biến đổi: chuẩn hoá & \texttt{preprocess.py} \\
    8  & Chia train/validation phân tầng & Đánh giá mô hình & Notebook \\
    9  & Trọng số lớp & Xử lý mất cân bằng lớp & Notebook \\
    10 & Tăng cường dữ liệu (chỉ khi huấn luyện) & Sinh thêm dữ liệu & \texttt{preprocess.py} \\
    \bottomrule
  \end{tabular}
\end{table}

Bước 1--5 thu hẹp tập train từ 3.875 tệp xuống 2.772 video sạch (Bảng~\ref{tab:pp_flow}). Bước 6, 7 và 10 không bỏ video nào mà quyết định video được đưa vào mô hình dưới dạng nào. Bước 8 và 9 chuẩn bị cho huấn luyện và đánh giá.

% =====================================================================
\subsubsection{Làm sạch dữ liệu}
\label{subsec:pp_clean}

\paragraph{Quét toàn bộ video} Notebook mở từng video bằng \texttt{cv2.VideoCapture}, đọc \emph{tất cả} khung hình và ghi lại: mã băm MD5 của tệp, số khung, và danh sách chỉ số các khung không đơn sắc (Mục~\ref{subsec:pp_frames}). Việc quét được chạy song song trên nhiều lõi CPU (\texttt{ProcessPoolExecutor}). Để kiểm tra bước quét, nhóm đối chiếu số khung đơn sắc của từng video với cột \texttt{solid\_frames} trong \texttt{video\_quality\_flags.csv} của bước EDA: cả 3.875/3.875 tệp đều khớp.

\paragraph{Bước 1 -- Kiểm tra toàn vẹn} Cả 3.875 tệp đều đọc được và có ít nhất 16 khung, nên không có tệp nào bị loại ở bước này. Tên lớp ở thư mục và khoá trong \texttt{label\_mapping.pkl} được chuẩn hoá về dạng Unicode NFC trước khi so khớp (Mục~\ref{subsec:eda_structure}); sau chuẩn hoá, 100/100 lớp có nhãn số 0--99.

\paragraph{Bước 2 -- Khử trùng lặp} Hai tệp có cùng mã MD5 là cùng một video. Với mỗi mã MD5, notebook giữ lại một tệp (ưu tiên tệp không có hậu tố \texttt{\_N}) và loại các tệp còn lại. Bước này loại 1.099 tệp, trong đó 1.069 tệp mang hậu tố \texttt{\_N}, và còn \textbf{2.776 video duy nhất}, khớp với kết quả EDA. Từ đây mỗi video trong tập sạch có một mã MD5 riêng; tính chất này giúp chia tập kiểm định không bị rò rỉ (Mục~\ref{subsec:pp_split}).

\paragraph{Bước 3 -- Xung đột nhãn} Hai nội dung video vừa nằm ở lớp \textit{Chào} vừa nằm ở lớp \textit{Chậm lại} (5 tệp: \texttt{Chào/392683}, \texttt{Chào/468203}, \texttt{Chào/792094}, \texttt{Chậm lại/254739}, \texttt{Chậm lại/284430}). Nhóm không xác minh được nhãn nào đúng, và giữ một video với nhãn có thể sai sẽ đưa nhiễu nhãn vào cả huấn luyện lẫn đánh giá, nên loại toàn bộ hai nội dung này. Notebook có sẵn biến \texttt{KEEP\_LABEL} để giữ lại các video này với nhãn đã xác minh nếu sau này có thông tin. Sau bước này còn 2.774 video. Hai lớp bị ảnh hưởng vẫn còn nhiều mẫu (\textit{Chào} 22 video, \textit{Chậm lại} 44 video).

\paragraph{Bước 4 -- Video rỗng} Hai video không có khung hợp lệ nào (\texttt{Biết/214576}: 16/16 khung đơn sắc; \texttt{Khai báo/885339}: 23/23 khung đơn sắc) bị loại vì không chứa thông tin ký hiệu. Tập train sạch còn \textbf{2.772 video} và vẫn đủ 100 lớp.

\begin{table}[H]
  \centering
  \caption{Số video còn lại sau từng bước làm sạch tập train.}
  \label{tab:pp_flow}
  \small
  \begin{tabular}{@{}p{0.3\textwidth}rp{0.45\textwidth}@{}}
    \toprule
    \textbf{Bước} & \textbf{Còn lại} & \textbf{Ghi chú} \\
    \midrule
    Ban đầu                          & 3.875 & Toàn bộ tệp trong \texttt{train/} \\
    1 -- Kiểm tra toàn vẹn           & 3.875 & Không có tệp lỗi hoặc thiếu khung \\
    2 -- Khử trùng lặp MD5           & 2.776 & Loại 1.099 bản sao \\
    3 -- Xung đột nhãn               & 2.774 & Loại 2 nội dung (5 tệp) \textit{Chào} $\leftrightarrow$ \textit{Chậm lại} \\
    4 -- Video rỗng                  & \textbf{2.772} & Loại \texttt{Biết/214576}, \texttt{Khai báo/885339} \\
    5 -- Khung đơn sắc               & 2.772 & 563 video có khung bị bỏ qua, không loại video \\
    \bottomrule
  \end{tabular}
\end{table}

Danh sách 1.103 tệp bị loại cùng lý do được lưu ở \texttt{processing/output/removed.csv}, để có thể kiểm tra lại từng trường hợp. Hình~\ref{fig:pp_class_counts} so sánh số video mỗi lớp trước và sau khi làm sạch: sau làm sạch mỗi lớp có từ 5 đến 51 video (trung bình 27,7; trung vị 32,5).

\begin{figure}[H]
  \centering
  \includegraphics[width=\textwidth]{processing/class_counts_before_after.png}
  \caption{Số video mỗi lớp trước (xám) và sau (xanh) khi làm sạch, sắp xếp theo số video sau làm sạch.}
  \label{fig:pp_class_counts}
\end{figure}

% =====================================================================
\subsubsection{Xử lý khung nhiễu và lấy mẫu theo thời gian}
\label{subsec:pp_frames}

\paragraph{Bước 5 -- Đánh dấu khung đơn sắc} Gọi $f_i$ là khung thứ $i$ của video và $\sigma_c(f_i)$ là độ lệch chuẩn giá trị điểm ảnh trên kênh màu $c\in\{R,G,B\}$. Khung $f_i$ được coi là \emph{đơn sắc} nếu
\begin{equation}
  \max_{c\in\{R,G,B\}} \sigma_c(f_i) < 3,
\end{equation}
tức gần như mọi điểm ảnh có cùng một màu; ngưỡng này giống với ngưỡng ở bước EDA. Các khung còn lại là \emph{khung hợp lệ}. Thay vì xoá khung khỏi tệp video, notebook lưu danh sách chỉ số khung hợp lệ $V=(v_1<v_2<\dots<v_m)$ của từng video vào cột \texttt{valid\_idx} của manifest; khi đọc video, \texttt{preprocess.py} chỉ lấy mẫu trong danh sách này.

Trong 2.772 video sạch, 563 video (20,3\%, thuộc 78 lớp) có khung đơn sắc, trung bình 12,2 khung mỗi video. Sau khi bỏ khung đơn sắc, mỗi video còn từ 9 đến 30 khung hợp lệ (trung vị 30); 222 video còn ít hơn 16 khung hợp lệ (Hình~\ref{fig:pp_valid_hist}).

\begin{figure}[H]
  \centering
  \includegraphics[width=0.75\textwidth]{processing/valid_frames_hist.png}
  \caption{Phân phối số khung hợp lệ mỗi video sau khi bỏ khung đơn sắc (trục tung thang log). Các video bên trái đường đứt nét có ít hơn $T=16$ khung hợp lệ.}
  \label{fig:pp_valid_hist}
\end{figure}

\paragraph{Bước 6 -- Lấy mẫu $T=16$ khung} Mô hình cần chuỗi đầu vào có độ dài cố định, nên mỗi video được thu gọn thành $T=16$ khung lấy đều theo thời gian trên danh sách khung hợp lệ:
\begin{equation}
  s_k = v_{\,1+\operatorname{round}\left(k\,\frac{m-1}{T-1}\right)}, \qquad k = 0, 1, \dots, T-1.
\end{equation}
Khung đầu và khung cuối của chuỗi luôn là khung hợp lệ đầu tiên và cuối cùng. Với video có $m<16$ khung hợp lệ, công thức trên tự lặp lại đều các khung (mỗi khung xuất hiện một hoặc hai lần) thay vì chỉ lặp khung cuối, nên chuỗi vẫn giữ đúng thứ tự thời gian. Nếu một video không có khung hợp lệ nào (chỉ có thể xảy ra ở dữ liệu kiểm tra hoặc demo) thì dùng toàn bộ khung.

Để thấy tác dụng của bước 5, nhóm so sánh với cách lấy mẫu đều trực tiếp trên toàn bộ khung, không lọc khung đơn sắc (Bảng~\ref{tab:pp_sampling}, Hình~\ref{fig:pp_sampling}). Nếu không lọc, cả 563 video có khung đơn sắc đều bị lẫn khung đơn sắc vào chuỗi đầu vào, trung bình 6,4/16 khung, và 222 video bị lẫn từ một nửa chuỗi trở lên; các số này khớp với mô phỏng ở bước EDA (Mục~\ref{subsec:eda_quality}). Khi lấy mẫu trên khung hợp lệ thì không còn khung đơn sắc nào trong đầu vào.

\begin{table}[H]
  \centering
  \caption{Số khung đơn sắc trong chuỗi 16 khung đầu vào, khi không lọc và khi lọc khung đơn sắc (2.772 video train sạch).}
  \label{tab:pp_sampling}
  \small
  \begin{tabular}{@{}lrr@{}}
    \toprule
    \textbf{Chỉ số} & \textbf{Không lọc} & \textbf{Có lọc} \\
    \midrule
    Video có $\geq$1 khung đơn sắc trong đầu vào   & 563 & 0 \\
    Số khung đơn sắc trung bình (trên các video đó) & 6,4/16 & 0 \\
    Video có $\geq$8/16 khung đơn sắc               & 222 & 0 \\
    \bottomrule
  \end{tabular}
\end{table}

\begin{figure}[H]
  \centering
  \includegraphics[width=\textwidth]{processing/sampling_naive_vs_filtered.png}
  \caption{Mười sáu khung đầu vào của video \textit{Nôn ói}/915809 (21/30 khung đơn sắc). Hàng trên: lấy mẫu trực tiếp, 11/16 khung là khung đơn sắc. Hàng dưới: lấy mẫu trên các khung hợp lệ $f_{21}$--$f_{29}$, các khung được lặp đều.}
  \label{fig:pp_sampling}
\end{figure}

% =====================================================================
\subsubsection{Chuẩn hoá}
\label{subsec:pp_norm}

\paragraph{Bước 7 -- Chuẩn hoá z-score} Giá trị điểm ảnh được đưa về $[0,1]$ rồi chuẩn hoá z-score trên từng kênh màu:
\begin{equation}
  x'_c = \frac{x_c/255 - \mu_c}{\sigma_c}, \qquad \mu = (0{,}485;\ 0{,}456;\ 0{,}406), \quad \sigma = (0{,}229;\ 0{,}224;\ 0{,}225),
\end{equation}
với $\mu,\sigma$ là trung bình và độ lệch chuẩn của tập ImageNet. Nhóm dùng thống kê của ImageNet thay vì tính lại trên tập train vì các mô hình ở Chương~\ref{sec:model} dùng mạng trích đặc trưng đã huấn luyện sẵn trên ImageNet; đầu vào cần có cùng phân phối với dữ liệu mà mạng đó đã học. Kết quả của mỗi video là một tensor kích thước $16\times3\times224\times224$. Độ phân giải đã đồng nhất $224\times224$ (Bảng~\ref{tab:eda_res}) nên không cần đổi kích thước.

% =====================================================================
\subsubsection{Chia tập huấn luyện và kiểm định}
\label{subsec:pp_split}

\paragraph{Bước 8 -- Chia phân tầng} Chỉ tập train có nhãn, nên tập kiểm định (\textit{validation}) phải tách ra từ chính tập này. Nhóm chia 2.772 video sạch thành 5 phần bằng \texttt{StratifiedKFold} của scikit-learn~\cite{sklearn} (\texttt{shuffle=True}, \texttt{random\_state=42}) và ghi số thứ tự phần vào cột \texttt{fold} của manifest. Phần \texttt{fold = 0} là tập kiểm định chung cho \emph{mọi} mô hình, các phần còn lại là tập huấn luyện; nhờ vậy kết quả của các mô hình so sánh được với nhau. Mô hình SVM có thể dùng cả 5 phần để kiểm định chéo (\textit{cross-validation}).

\begin{itemize}
  \item \textbf{Phân tầng theo nhãn:} mỗi phần giữ tỷ lệ các lớp gần giống toàn tập. Lớp nhỏ nhất có 5 video nên có đúng 1 video ở mỗi phần, nhờ đó cả 100 lớp đều có mặt ở tập kiểm định.
  \item \textbf{Không rò rỉ dữ liệu:} EDA khuyến nghị chia theo nhóm video gốc (\texttt{StratifiedGroupKFold}) để các bản sao không rơi vào cả hai phía. Vì bước 2 đã loại mọi bản sao, mỗi video sạch là một nhóm riêng, nên chia theo từng video cũng chính là chia theo nhóm. Notebook kiểm tra lại bằng \texttt{assert}: mã MD5 là duy nhất trong manifest và không có mã MD5 nào nằm ở cả tập huấn luyện lẫn tập kiểm định. Nhóm đã thử \texttt{StratifiedGroupKFold}, nhưng thuật toán này chỉ cân bằng gần đúng và tạo ra phần thiếu lớp.
\end{itemize}

\begin{table}[H]
  \centering
  \caption{Kết quả chia tập train sạch.}
  \label{tab:pp_split}
  \small
  \begin{tabular}{@{}lrrr@{}}
    \toprule
    \textbf{Tập} & \textbf{Số video} & \textbf{Số lớp} & \textbf{Số video mỗi lớp} \\
    \midrule
    Huấn luyện (\texttt{fold} 1--4) & 2.217 & 100 & 4--41 \\
    Kiểm định (\texttt{fold} 0)     & 555   & 100 & 1--10 \\
    \midrule
    Tổng                            & 2.772 & 100 & 5--51 \\
    \bottomrule
  \end{tabular}
\end{table}

% =====================================================================
\subsubsection{Xử lý mất cân bằng lớp}
\label{subsec:pp_weights}

\paragraph{Bước 9 -- Trọng số lớp} Trên tập huấn luyện, lớp nhiều nhất (\textit{Chạy}) có 41 video còn nhiều lớp chỉ có 4 video (\textit{Cách ly}, \textit{Ho}, \textit{Sốt}, \textit{Virus}, \textit{Xe máy},\ldots), chênh khoảng 10,2 lần. Nếu mọi mẫu có vai trò như nhau trong hàm mất mát, mô hình sẽ ưu tiên các lớp đông và Macro F1 giảm ở các lớp ít mẫu. Nhóm tính trọng số cho từng lớp $c$ theo kiểu \textit{balanced}:
\begin{equation}
  w_c = \frac{N}{K \cdot n_c},
\end{equation}
với $N=2.217$ là số video huấn luyện, $K=100$ là số lớp và $n_c$ là số video huấn luyện của lớp $c$. Lớp ít mẫu nhận trọng số lớn: trọng số dao động từ 0,54 (\textit{Chạy}) đến 5,54 (các lớp có 4 video), và một lớp trung bình ($n_c = N/K$) có trọng số bằng 1. Trọng số chỉ tính trên tập huấn luyện để không dùng thông tin của tập kiểm định. Các mô hình dùng trọng số này trong hàm mất mát (ví dụ tham số \texttt{weight} của \texttt{CrossEntropyLoss}) hoặc tham số \texttt{class\_weight} của SVM.

% =====================================================================
\subsubsection{Tăng cường dữ liệu}
\label{subsec:pp_aug}

\paragraph{Bước 10 -- Tăng cường} Tập huấn luyện nhỏ (trung bình khoảng 22 video mỗi lớp), nên mỗi lần đọc một video huấn luyện, \texttt{preprocess.py} tạo ra một biến thể ngẫu nhiên khác của video đó. Ba phép biến đổi được dùng, với cùng một tham số cho mọi khung của một video để chuyển động không bị méo:

\begin{itemize}
  \item \textbf{Cắt đoạn theo thời gian:} giữ một đoạn liên tục dài 80--100\% danh sách khung hợp lệ ở vị trí ngẫu nhiên, rồi mới lấy 16 khung. Phép này mô phỏng việc người ký bắt đầu, kết thúc sớm hay muộn và ký nhanh hay chậm khác nhau.
  \item \textbf{Cắt vùng ảnh ngẫu nhiên:} cắt một vùng có cạnh bằng 85--100\% cạnh khung hình ở vị trí ngẫu nhiên rồi phóng lại về $224\times224$. Phép này mô phỏng khác biệt về vị trí người ký và khoảng cách tới máy quay.
  \item \textbf{Thay đổi màu:} nhân độ sáng, độ tương phản và độ bão hoà với các hệ số ngẫu nhiên trong khoảng $[0{,}75;\ 1{,}25]$, giúp mô hình bền vững với khác biệt ánh sáng và màu nền giữa train và test (Mục~\ref{subsec:eda_shift}).
\end{itemize}

Nhóm \emph{không} dùng lật ngang vì tay trái và tay phải có thể làm đổi nghĩa ký hiệu. Tăng cường chỉ áp dụng cho tập huấn luyện; tập kiểm định và tập kiểm tra được đưa vào mô hình nguyên bản để kết quả đánh giá ổn định. Hình~\ref{fig:pp_aug} minh hoạ ba biến thể tạo ra từ cùng một video.

\begin{figure}[H]
  \centering
  \includegraphics[width=0.9\textwidth]{processing/augmentation_examples.png}
  \caption{Một video lớp \textit{An ủi} (hàng đầu) và ba biến thể tăng cường ngẫu nhiên; mỗi hàng hiển thị 8 trong 16 khung đầu vào.}
  \label{fig:pp_aug}
\end{figure}

% =====================================================================
\subsubsection{Dữ liệu kiểm tra}
\label{subsec:pp_test}

Hai tập kiểm tra (\texttt{public\_test}, \texttt{private\_test}) không có nhãn và cuộc thi yêu cầu dự đoán cho \emph{mọi} tệp, nên nhóm không loại tệp nào. Notebook chỉ lập bảng \texttt{test\_manifest.csv} gồm tên tệp, mã MD5 và cột \texttt{dup\_of\_train} đánh dấu video trùng với tập train: có 157/1.630 video public test trùng với train và 0/2.859 video private test trùng với train. Các video trùng vẫn được mô hình dự đoán bình thường; cột này chỉ dùng khi diễn giải điểm public test (Mục~\ref{subsec:eda_quality}).

Khi suy luận, mỗi video kiểm tra đi qua đúng các bước 5, 6 và 7 như video kiểm định: tự tìm khung đơn sắc, lấy 16 khung trên các khung hợp lệ và chuẩn hoá, không tăng cường. Notebook chạy thử hàm đọc video này trên 8 video kiểm tra ngẫu nhiên để xác nhận đầu ra có đúng kích thước $16\times224\times224\times3$.

% =====================================================================
\subsubsection{Kết quả và đầu ra}
\label{subsec:pp_output}

Bảng~\ref{tab:pp_outputs} liệt kê các tệp do notebook tạo ra trong thư mục \texttt{processing/output/}. Các mô hình chỉ cần \texttt{preprocess.py}, \texttt{manifest\_clean.csv}, \texttt{test\_manifest.csv} và \texttt{class\_weights.csv} (vài trăm KB) cùng tệp dữ liệu gốc, không cần chạy lại bước làm sạch.

\begin{table}[H]
  \centering
  \caption{Các tệp đầu ra của bước tiền xử lý.}
  \label{tab:pp_outputs}
  \small
  \begin{tabular}{@{}p{0.3\textwidth}p{0.62\textwidth}@{}}
    \toprule
    \textbf{Tệp} & \textbf{Nội dung} \\
    \midrule
    \texttt{manifest\_clean.csv} & 2.772 video train sạch: lớp, nhãn số, tên tệp, đường dẫn, MD5, số khung, số khung hợp lệ, chỉ số khung hợp lệ (\texttt{valid\_idx}), phần (\texttt{fold}) \\
    \texttt{removed.csv} & 1.103 tệp bị loại, kèm bước và lý do \\
    \texttt{test\_manifest.csv} & 4.489 video kiểm tra: tập, tên tệp, MD5, \texttt{dup\_of\_train} \\
    \texttt{class\_weights.csv} & Trọng số 100 lớp tính trên tập huấn luyện \\
    \texttt{processing\_summary.csv} & Số video còn lại sau từng bước (Bảng~\ref{tab:pp_flow}) \\
    \texttt{class\_counts\_before\_after.csv} & Số video mỗi lớp trước và sau làm sạch \\
    \texttt{processing\_config.json} & Tham số đã dùng (\texttt{SEED}, số phần, $T$, ngưỡng đơn sắc) \\
    \texttt{figures/} & Các hình trong mục này \\
    \bottomrule
  \end{tabular}
\end{table}

Tóm lại, sau tiền xử lý tập train không còn bản sao, video rỗng hay nhãn mâu thuẫn; không video nào còn khung đơn sắc trong chuỗi đầu vào; tập kiểm định đủ 100 lớp và không chứa video trùng với tập huấn luyện. Các bước này xử lý cả ba vấn đề chính mà EDA đã nêu: bản sao gây rò rỉ dữ liệu, khung đơn sắc chỉ có ở train và mất cân bằng lớp.

\paragraph{Hạn chế} Nhóm chưa xác định được nhãn đúng cho hai nội dung xung đột nên phải bỏ chúng, làm giảm 5 tệp dữ liệu có thể hữu ích. Ngưỡng đơn sắc chỉ phát hiện khung một màu hoàn toàn, chưa phát hiện khung mờ hoặc khung chỉ có nền. Khử trùng lặp bằng MD5 chỉ tìm được bản sao giống hệt từng byte, không tìm được video gần giống nhau (ví dụ cùng cảnh nhưng nén lại). Cuối cùng, tập kiểm định chỉ có 1--10 video mỗi lớp, nên F1 của từng lớp riêng lẻ trên tập này còn dao động lớn.


% =====================================================================
\subsection{Trích xuất đặc trưng}
\label{subsec:features}

Video là dữ liệu phi cấu trúc: mỗi video sau tiền xử lý là một tensor $16\times3\times224\times224$ (khoảng 2,4 triệu giá trị). Các thuật toán phân lớp cổ điển không làm việc trực tiếp được với dữ liệu này, nên nhóm biến mỗi khung thành một vector đặc trưng bằng EfficientNet-B0 (Mục~\ref{subsec:cnn}), rồi gộp 16 vector theo thời gian thành một vector cho cả video.

\paragraph*{Hai bộ đặc trưng}
\begin{itemize}
  \item \textbf{Đóng băng} (\texttt{models/extract\_features.ipynb}): EfficientNet-B0 với trọng số ImageNet, bỏ lớp phân loại, chạy ở chế độ suy luận trên GPU Tesla T4. Mỗi video cho ra mảng $16\times1.280$. Ngoài bản gốc, notebook tạo thêm 5 bản tăng cường cho mỗi video train (Mục~\ref{subsec:pp_aug}) với hạt giống cố định, dùng để huấn luyện CNN-LSTM giai đoạn A.
  \item \textbf{Tinh chỉnh} (\texttt{models/extract\_finetuned.ipynb}): EfficientNet-B0 đã được tinh chỉnh trong CNN-LSTM giai đoạn B (Mục~\ref{subsec:cnnlstm_design}). Cùng định dạng $16\times1.280$, kèm vector video 512 chiều sau attention và trọng số attention của từng khung.
\end{itemize}
Thứ tự dòng của mọi tệp đặc trưng khớp với \texttt{manifest\_clean.csv} và \texttt{test\_manifest.csv}. Đặc trưng lưu ở dạng số thực 16 bit (khoảng 300~MB mỗi bộ cho cả ba tập).

\paragraph*{Gộp theo thời gian} Với chuỗi vector khung $\mathbf{f}_1,\dots,\mathbf{f}_{16}\in\mathbb{R}^{1280}$, nhóm thử ba cách gộp cho các thuật toán cổ điển:
\begin{align}
  \text{F1:}\quad & \bar{\mathbf{f}} = \frac{1}{16}\sum_{t=1}^{16}\mathbf{f}_t \in\mathbb{R}^{1280} && \text{(tư thế trung bình, mất thứ tự)}\\
  \text{F2:}\quad & [\bar{\mathbf{f}};\ \mathbf{s}],\quad s_j=\sqrt{\tfrac{1}{16}\textstyle\sum_t(f_{t,j}-\bar{f}_j)^2} \in\mathbb{R}^{2560} && \text{(thêm mức độ thay đổi theo thời gian)}\\
  \text{F3:}\quad & [\bar{\mathbf{f}}^{(1)};\bar{\mathbf{f}}^{(2)};\bar{\mathbf{f}}^{(3)};\bar{\mathbf{f}}^{(4)}] \in\mathbb{R}^{5120} && \text{(trung bình 4 đoạn liên tiếp: đầu, giữa, cuối)}
\end{align}
Thí nghiệm với SVM (Mục~\ref{subsec:res_svm}) cho thấy F2 tốt nhất, nên phần so sánh thuật toán dùng F2. Nửa sau của F2 (độ lệch chuẩn) mô tả \emph{chuyển động}: một chiều đặc trưng thay đổi nhiều qua các khung thì có độ lệch chuẩn lớn.

\paragraph*{Lưu ý về rò rỉ thông tin} CNN tinh chỉnh đã học trên chính các video fold 1--4, nên đặc trưng tinh chỉnh của các video này đã ``thấy'' nhãn của chúng. Vì vậy cross-validation trên fold 1--4 với bộ đặc trưng tinh chỉnh cho điểm lạc quan và chỉ dùng để chọn siêu tham số; kết quả báo cáo dựa trên fold~0 (CNN chỉ dùng fold~0 để chọn epoch dừng) và tập private.

% =====================================================================
\subsection{Thiết kế mô hình}
\label{sec:model}

\subsubsection{Các thuật toán phân lớp cổ điển}
\label{subsec:classical_design}

Mọi thuật toán dùng chung pipeline \texttt{StandardScaler} $\rightarrow$ \texttt{PCA} $\rightarrow$ bộ phân lớp của scikit-learn~\cite{sklearn}. Pipeline được học lại từ đầu trong mỗi lần chia cross-validation, nên chuẩn hoá và PCA chỉ dùng phần huấn luyện của lần chia đó. Logistic Regression, Random Forest và SVM dùng trọng số lớp \texttt{balanced} (Mục~\ref{subsec:pp_weights}); kNN không hỗ trợ trọng số lớp.

\paragraph*{SVM (thí nghiệm riêng)} Notebook \texttt{models/svm/svm.ipynb} khảo sát rộng hơn cho SVM trên đặc trưng đóng băng: ba cách gộp F1, F2, F3; PCA giữ 95\% phương sai hoặc 256 thành phần; SVM tuyến tính (\texttt{LinearSVC}, $C\in\{0{,}001;\,0{,}01\}$) và SVM nhân RBF (\texttt{SVC}, $C\in\{1;\,10\}$, $\gamma\in\{\texttt{scale};\,0{,}001\}$), tổng cộng 36 cấu hình.

\paragraph*{So sánh bốn thuật toán} Notebook \texttt{models/classical/classical.ipynb} so sánh bốn thuật toán trên cả hai bộ đặc trưng (gộp F2, PCA 95\%) với lưới siêu tham số ở Bảng~\ref{tab:grids}.

\begin{table}[H]
  \centering
  \caption{Thuật toán và lưới siêu tham số trong phần so sánh.}
  \label{tab:grids}
  \small
  \begin{tabular}{@{}p{0.17\textwidth}p{0.12\textwidth}p{0.21\textwidth}p{0.42\textwidth}@{}}
    \toprule
    \textbf{Thuật toán} & \textbf{Vai trò} & \textbf{Tìm tham số} & \textbf{Giá trị thử} \\
    \midrule
    Logistic Regression & baseline & GridSearchCV & $C\in\{0{,}01;\,0{,}1;\,1\}$ \\
    kNN & dựa trên khoảng cách & GridSearchCV & $k\in\{1;3;5;10;20\}$; trọng số đều / theo khoảng cách; khoảng cách euclid / cosine (20 cấu hình) \\
    Random Forest & ensemble & RandomizedSearchCV (8 cấu hình) & số cây $\in\{300;\,600\}$; độ sâu tối đa $\in\{\text{không giới hạn};\,30\}$; số thuộc tính mỗi lần tách $\in\{\sqrt{d};\,5\%\}$; số mẫu tối thiểu ở lá $\in\{1;\,2\}$ \\
    SVM tuyến tính & lề cực đại & GridSearchCV & $C\in\{0{,}001;\,0{,}01\}$ \\
    \bottomrule
  \end{tabular}
\end{table}

\subsubsection{Mô hình CNN-LSTM}
\label{subsec:cnnlstm_design}

Mô hình đọc 16 khung theo thứ tự thời gian (Bảng~\ref{tab:cnnlstm_arch}). Phần sau CNN (gọi là \emph{phần đầu}) gồm khoảng 3,9 triệu tham số.

\begin{table}[H]
  \centering
  \caption{Kiến trúc CNN-LSTM.}
  \label{tab:cnnlstm_arch}
  \small
  \begin{tabular}{@{}p{0.16\textwidth}p{0.55\textwidth}p{0.2\textwidth}@{}}
    \toprule
    \textbf{Khối} & \textbf{Cấu hình} & \textbf{Đầu ra} \\
    \midrule
    CNN & EfficientNet-B0, chạy riêng từng khung & $B\times16\times1.280$ \\
    Chuẩn hoá & trừ trung bình, chia độ lệch chuẩn của từng chiều (tính trên fold 1--4) & $B\times16\times1.280$ \\
    Chiếu & Linear $\rightarrow$ 512, LayerNorm, ReLU, Dropout 0,3 & $B\times16\times512$ \\
    BiLSTM & 2 lớp, 256 đơn vị mỗi chiều, dropout 0,3 giữa hai lớp & $B\times16\times512$ \\
    Attention & $e_t=\mathbf{v}^\top\tanh(W\mathbf{h}_t+\mathbf{b})$, $W\in\mathbb{R}^{128\times512}$, softmax theo thời gian & $B\times512$ \\
    Phân loại & Dropout 0,4, Linear 512 $\rightarrow$ 100 & $B\times100$ \\
    \bottomrule
  \end{tabular}
\end{table}

\paragraph*{Giai đoạn A -- CNN đóng băng} Phần đầu học trên đặc trưng đóng băng đã trích sẵn, nên mỗi epoch chỉ mất vài giây. Mỗi epoch, mỗi video lấy ngẫu nhiên 1 trong 5 bản tăng cường, rồi được tăng cường thêm trên chuỗi đặc trưng: che ngẫu nhiên 0--3 khung, cộng nhiễu Gauss độ lệch chuẩn 0,1 và dropout đầu vào 0,1. Huấn luyện bằng AdamW (tốc độ học $10^{-3}$, suy giảm trọng số $10^{-2}$, batch 64), 3 epoch khởi động rồi giảm theo cosine, cross-entropy có trọng số lớp và label smoothing 0,1, cắt gradient 1,0, tối đa 60 epoch, dừng sớm sau 10 epoch Macro-F1 fold~0 không tăng. Hai biến thể:
\begin{itemize}
  \item \textbf{A0}: đầu vào là chuỗi đặc trưng gốc.
  \item \textbf{A2}: đầu vào \emph{chuyển động} $[\mathbf{f}_t;\ \mathbf{f}_t-\mathbf{f}_{t-1}]$ (2.560 chiều); phần không đổi theo thời gian như nền, trang phục bị triệt tiêu ở nửa sau.
\end{itemize}
Sau khi chọn epoch tốt nhất trên fold~0, mỗi biến thể được huấn luyện lại trên cả 2.772 video đúng số epoch đó để ghi tệp nộp.

\paragraph*{Giai đoạn B -- tinh chỉnh CNN} EfficientNet-B0 được nối vào trước phần đầu và huấn luyện cùng nhau trên video. Phần đầu khởi tạo từ mô hình A0 (đã học trên fold 1--4); stem và hai nhóm khối đầu của EfficientNet (học cạnh, màu cơ bản) giữ đóng băng, các lớp BatchNorm giữ thống kê ImageNet vì mỗi batch chỉ có 8 video. Mỗi epoch đọc lại video và tăng cường mới (giữ 80--100\% đoạn thời gian, cắt 75--100\% khung hình, đổi độ sáng, độ tương phản, độ bão hoà $\pm25\%$; không lật ngang vì tay trái và tay phải có thể đổi nghĩa ký hiệu). Tốc độ học $10^{-4}$ cho CNN và $5\cdot10^{-4}$ cho phần đầu, 1 epoch khởi động rồi cosine, tính toán số thực hỗn hợp 16/32 bit, tối đa 15 epoch, dừng sớm sau 4 epoch. Mô hình nộp bài là checkpoint có Macro-F1 fold~0 cao nhất (không huấn luyện lại trên toàn bộ train vì chi phí GPU).

% =====================================================================
\subsection{Thiết kế thí nghiệm và đánh giá}
\label{subsec:eval_design}

\paragraph*{Chia dữ liệu} Theo Mục~\ref{subsec:pp_split}, 2.772 video sạch được chia phân tầng thành 5 phần. \textbf{Fold~0} (555 video) được giữ riêng làm tập kiểm định, không dùng khi chọn siêu tham số. Siêu tham số của các thuật toán cổ điển được chọn bằng cross-validation 4 lần trên \textbf{fold 1--4} (2.217 video, dùng \texttt{PredefinedSplit} để giữ đúng cách chia), tiêu chí Macro-F1.

\paragraph*{Ba mức đánh giá}
\begin{enumerate}
  \item \textbf{Cross-validation trên fold 1--4}: chỉ để chọn siêu tham số.
  \item \textbf{Fold~0}: mô hình với tham số tốt nhất được huấn luyện trên fold 1--4 và đánh giá một lần bằng Macro-F1~\eqref{eq:macro_f1}, Macro-Precision, Macro-Recall, Accuracy~\eqref{eq:accuracy_multi}, Top-5 Accuracy và ROC-AUC (Mục~\ref{subsec:metrics}), kèm thời gian huấn luyện và dự đoán. Với SVM tuyến tính (không cho xác suất), điểm quyết định được đưa qua hàm softmax để tính Top-5 và ROC-AUC; thứ tự các lớp không đổi.
  \item \textbf{Tập private test (2.859 video)}: mô hình được huấn luyện lại trên toàn bộ 2.772 video, dự đoán tập private và nộp lên AIChallenge để nhận Macro-F1 do ban tổ chức chấm. Đây là đánh giá độc lập duy nhất trên video chưa từng được dùng.
\end{enumerate}

\paragraph*{Định dạng tệp nộp} Tệp \texttt{cv\_submission.csv} gồm hai cột \texttt{video\_name,label}, mỗi video private đúng một dòng, nén thành tệp \texttt{.zip}. Nhãn phải viết đúng như khoá trong \texttt{label\_mapping.pkl} (dạng Unicode NFD, Mục~\ref{subsec:eda_structure}); các notebook tự kiểm tra số dòng, tên tệp không trùng và dạng Unicode trước khi ghi.

\paragraph*{Kiểm tra tính đúng} 156 video public test giống hệt video train (cùng mã MD5). Mọi mô hình phải đoán đúng nhãn train cho gần hết các video này; nếu không, có lỗi lệch thứ tự dòng giữa đặc trưng và manifest. Phép kiểm tra chỉ để phát hiện lỗi, dự đoán không bị thay bằng nhãn train.

% =====================================================================
\subsection{Thiết kế ứng dụng demo}
\label{subsec:demo_design}

% TODO: cập nhật theo ứng dụng thực tế sau khi hoàn thành demo/
Ứng dụng web viết bằng Streamlit cho phép người dùng tải lên một video ký hiệu và nhận kết quả nhận diện. Luồng xử lý dùng lại đúng các hàm của \texttt{processing/preprocess.py} và mô hình CNN-LSTM giai đoạn B, nên dữ liệu ở bước demo được xử lý giống hệt lúc huấn luyện:
\begin{enumerate}
  \item Người dùng tải video (\texttt{.mp4}, \texttt{.avi}, \texttt{.mov}) hoặc chọn video mẫu.
  \item Ứng dụng bỏ khung đơn sắc, lấy đều 16 khung, chuẩn hoá ImageNet.
  \item Mô hình trả về 5 từ có xác suất cao nhất cùng xác suất tương ứng.
  \item Ứng dụng hiển thị các khung đã lấy mẫu kèm trọng số attention của từng khung, cho biết khung nào ảnh hưởng nhiều nhất đến quyết định.
\end{enumerate}

\subsection{Công cụ và môi trường}
\label{subsec:tools}

\begin{itemize}
  \item \textbf{Ngôn ngữ và thư viện}: Python; \texttt{pandas}, \texttt{numpy}, \texttt{OpenCV} (đọc video, xử lý dữ liệu); \texttt{scikit-learn}~\cite{sklearn} (pipeline, PCA, các thuật toán cổ điển, tìm siêu tham số, chỉ số); PyTorch và \texttt{torchvision} (EfficientNet-B0, CNN-LSTM); \texttt{matplotlib} (trực quan hoá); Streamlit (ứng dụng demo).
  \item \textbf{Phần cứng}: Kaggle Notebook với GPU NVIDIA Tesla T4 cho trích xuất đặc trưng và huấn luyện CNN-LSTM; máy tính cá nhân cho các thuật toán cổ điển và trích xuất đặc trưng tinh chỉnh.
  \item \textbf{Theo dõi thí nghiệm}: Weights \& Biases, các tệp \texttt{history.csv}, \texttt{summary.json} trong thư mục output của mỗi notebook.
  \item \textbf{Quản lý mã nguồn}: GitHub, kèm \texttt{README.md} và \texttt{requirements.txt}.
\end{itemize}
